\documentclass[11pt,a4paper]{article}


\usepackage[inline]{asymptote}
\usepackage{asypictureB}

\usepackage{qrcode}
%\usepackage[default]{frcursive}
\usepackage[legal,text={17.5cm,25cm},centering,top=2.5cm]{geometry} 
\usepackage{amsfonts}
%\usepackage[colorlinks=true,linkcolor=black]{hyperref}
%\usepackage{amsmath}
\usepackage{amssymb,amsmath,stackrel,mathrsfs,bm}
\usepackage{multicol}
\usepackage{pdflscape}% girar una pagina, hay que usar \begin{landscape}En este punto inicia la hoja en vertical \end{landscape}
%\usepackage[mathcal]{euscript}
%\usepackage{amssymb}
\usepackage{esvect}%para la flecha de vector 

\usepackage[T1]{fontenc}
\usepackage[spanish,es-noshorthands]{babel}
%\usepackage[latin1,ansinew]{inputenc}\usepackage{calligra}
%\usepackage[pdftex]{graphicx}
%\usepackage{etex}
\usepackage{calligra}
\usepackage{array}
 \usepackage{booktabs}
 \usepackage{ifsym}
\usepackage{pgf,tikz,pgfplots}
\usepackage{tkz-graph}%paquete (di)grafos
\usepackage{tkz-euclide}% Para realizar diagramas que coloreen conjuntos
\usetikzlibrary[decorations.text]% Decorar curvas
%\usetikzlibrary{matrix}
\usetikzlibrary{arrows,matrix,calc}
%\usetkzobj{all}

\pgfplotsset{compat=1.15}
\usepackage{mathrsfs}
\usetikzlibrary{arrows}
\usepackage{makecell} %Paquete para controlar altura de celdas
\usepackage{verbatim} 
\usetikzlibrary{babel}%para no general error al usar babel y tikz simultaneamente.
\usetikzlibrary{patterns,patterns.meta}%%%%%%%%%%%para  pintar de forma rrallada un rectangulo

\usepackage{pictex}
\usepackage{verbatim}%me sirve para usar comentarios con mucha lineas 
\usepackage{amstext}
\usepackage{enumerate}
%\usepackage{enumitem}
\usepackage{amsthm}
\usepackage{xcolor}
\usepackage{mathdots}% para tener mas opciones de puntos suspensivos

%\usepackage{pgf,tikz}%,pgfplots
%\pgfplotsset{compat=1.15}
\usepackage{mathrsfs}
%\usetikzlibrary{arrows}
%\usepackage{etex,tkz-euclide} 
%\usetkzobj{all}
\usepackage{color}

\usepackage{cancel}
\newcommand{\Cancelar}[2][black]{{\color{#1}\cancel{\color{black}#2}}}%%% para cancelar con colores recordar usar mayusculas
\usepackage{lscape}%paquete para girar hojas












\newcommand{\R}{\mathbb{R}}
\newcommand{\re}{\mathbb{R}}
\newcommand{\I}{\mathbb{I}}
\newcommand{\Q}{\mathbb{Q}}
\newcommand{\Z}{\mathbb{Z}}
\newcommand{\C}{\mathbb{C}}
\newcommand{\N}{\mathbb{N}}
\newcommand{\K}{\mathbb{K}}
%\newcommand{\B}{\mathcal{B}}
\newcommand{\V}{\mathbb{V}}
\newcommand{\can}{\mathcal{C}}
\newcommand{\M}{\mathbb{M}}
\newcommand{\Pol}{\mathcal{P}}
\newcommand{\iv}{\bm{\check{\imath}}}
\newcommand{\jv}{\bm{\check{\jmath}}}
\newcommand{\kv}{\check{\bm k}} 

\newcommand{\dis}{\displaystyle}



\newtheorem{ejer}{Ejercicio}[]

\newtheorem{teo}{Teorema}%[section]
\newtheorem{lema}[teo]{Lema}
\newtheorem{propo}[teo]{Proposición}
\newtheorem{defi}[teo]{Definición}
\newtheorem{ejem}[teo]{}
\newtheorem*{ejemsinn}{}
\newtheorem{obs}[teo]{Observación}%[section]
\newtheorem{acla}{Aclaración}%[section]
\newtheorem{act}{Actividad}%%%%%%%%%%%%%%%%
\newtheorem*{defsinn}{{\bf Definición}:}
\newtheorem*{ejemplosinn}{{\bf Ejemplo}:}
\newtheorem*{proposinn}{{\bf\sc Proposición}}
\newtheorem*{obssinn}{Observación:}
\newtheorem*{corosinn}{Corolario}
\newtheorem*{propsinn}{Propiedades:}
%\newenvironment{proof} { \noindent {\bf Demostración: }\rm}{$\quad \hfill \blacksquare $}
\usepackage{setspace} % (es una alternativa a \renewcommand{\baselinestretch}{1.3}) usando \doublespacing \onehalfspace \singlespace \spacing{1.5}  se va cambiando segun uno quiera

%\newenvironment{proof}[Demostración:]
\newenvironment{sol}{ \noindent {\bf Solución:} \rm}{%\hfill \protect\tikz[overlay,yshift=-4pt]\protect\node[xshift=8mm, yshift=13mm, blue]() {\huge $\checkmark$ };
}

%%%%%%%%%%%%%%%%%%%%%%%%%%%%%%%%%%%%%%%%


%\renewcommand{\baselinestretch}{1.5}
%%%%%%%%%%%%%%%%%%%%%%%%%%%%%%%%%%%%%%%%%%%%%%%%%%%%%%%%%%%%%%%%%%%%%%%%%%%%%%%%%%%%%%%%%%%%%%%%%%%%%%%%%%%%%%%%%%%%%%%%%%%%%%%%%%%%%%%%%%%%%%%%%%%%%%%%%%%%%%%%%%%%%%%%%%%%%%%%%%%%%%%%%%%%%%%%%%%%%%%%%%%%%%%%%%%%%%%%%%%%%%%%%%%%%%%%%%%%%%%%%%%%%%%%%%%%%%%%%%%%%%%%%%%%%%%%%%%%%%%%
\usepackage[framemethod=TikZ]{mdframed}%Paquete para estilos de los bordes de los ambientes Importante este paquete (mdframed) define un entorno  que trata automáticamente los saltos de página en el texto enmarcado.
            \mdfsetup{skipabove=1\baselineskip, skipbelow=1.5\baselineskip ,             frametitlefont= \sffamily\bfseries , needspace=4\baselineskip , splittopskip=1.5\baselineskip} %
            \mdfsetup{apptotikzsetting={\tikzset{mdfbackground/.append style={draw=none}}}} 
%%%%%%%%%%%%%%%%%%%%%%%%%%%%%%%%%%%%%%%%%%%%%%%%%%%%%%%%%%%%%%
%
%Estilo para Cuentas auxiliares
%
%%%%%%%%%%%%%%%%%%%%%%%%%%%%%%%%%%%%%%%%%%%%%%%%%%%%%%%%%%%%%%
\mdfdefinestyle{Cuentaauxi}{%skipabove=10pt, skipbelow=0,
userdefinedwidth=1.02\textwidth,align=center,
  everyline=true,
  ignorelastdescenders=true,
  middlelinewidth=1pt,
  linecolor=black!20,
  outerlinewidth=1pt,
  %leftline=false,topline=false
  %rightline=false,bottomline=false,
  %backgroundcolor=black!20,
  innerleftmargin=2mm, innerrightmargin=2mm,
  innerbottommargin=3mm,
  innertopmargin=1mm,
  frametitleaboveskip=2mm,
  frametitlebelowskip=2mm,
  %frametitlerule=false,
  %repeatframetitle=false,
 % outerlinewidth=6pt,outerlinecolor=blue!20,
 pstricksappsetting={\addtopsstyle{mdfouterlinestyle}{linestyle=dashed}},
 %firstextra={\node[xshift=5cm,right,draw=White, line width=1.5pt,%star,star points=10,star point ratio=1.2, 
 %minimum size=15mm, fill=white] at (P-|O) {\tikzpicture\draw[draw=black!20,decoration=snake, fill=white] (0,0) --(2.5,0) decorate{--(2.5,1)}--(0,1)decorate{--cycle};%\includegraphics[width=10mm]{.pdf}
 %  \endtikzpicture
 %}; },
 %singleextra={\node[xshift=5cm,right,draw=White, line width=1.5pt,%star,star points=10,star point ratio=1.2, minimum size=15mm,
  %fill=white] at (P-|O) {\tikzpicture\draw[draw=black!20,decoration=snake, fill=white] (0,0) --(2.5,0) decorate{--(2.5,1)}--(0,1)decorate{--cycle};%\includegraphics[width=10mm]{.pdf}    
 %  \endtikzpicture
 %}; }
}
%%%%%%%%%%%%%%%%%%%%%%%%%%%%%%%%%%%%%%%%%%%%%%%%%%%%%%%%%%%%%%
%
%Estilo para Citar
%
%%%%%%%%%%%%%%%%%%%%%%%%%%%%%%%%%%%%%%%%%%%%%%%%%%%%%%%%%%%%%%
\mdfdefinestyle{citar}{
userdefinedwidth=.83\textwidth ,align=center,
skipabove=5pt, skipbelow=5pt,
  everyline=true,
  ignorelastdescenders=true,
  middlelinewidth=1pt,
  linecolor=blue!60,
  outerlinewidth=2pt,
  %leftline=false,
  topline=false,
  rightline=false,
  bottomline=false,
  %backgroundcolor=black!20,
  innerleftmargin=5mm, innerrightmargin=5mm,
  innerbottommargin=3mm,
  innertopmargin=1mm,
  leftmargin= 14mm,
  %frametitleaboveskip=2mm,
  %frametitlebelowskip=2mm,
  %frametitlerule=false,
  %repeatframetitle=false,
 % outerlinewidth=6pt,outerlinecolor=blue!20,
 pstricksappsetting={\addtopsstyle{mdfouterlinestyle}{linestyle=dashed}},
 %firstextra={\node[xshift=5cm,right,draw=White, line width=1.5pt,%star,star points=10,star point ratio=1.2, 
 %minimum size=15mm, fill=white] at (P-|O) {\tikzpicture\draw[draw=black!20,decoration=snake, fill=white] (0,0) --(2.5,0) decorate{--(2.5,1)}--(0,1)decorate{--cycle};%\includegraphics[width=10mm]{.pdf}
 %  \endtikzpicture
 %}; },
 %singleextra={\node[xshift=5cm,right,draw=White, line width=1.5pt,%star,star points=10,star point ratio=1.2, minimum size=15mm,
  %fill=white] at (P-|O) {\tikzpicture\draw[draw=black!20,decoration=snake, fill=white] (0,0) --(2.5,0) decorate{--(2.5,1)}--(0,1)decorate{--cycle};%\includegraphics[width=10mm]{.pdf}    
 %  \endtikzpicture
 %}; }
}
%%%%%%%%%%%%%%%%%%%%%%%%%%%%%%%%%%%%%%%%%%%%%%%%%%%%%%%%%%%%%%%%%%%%%%%%%%%%%%%%%%%%%%%%%%%%%%%%%%%%%%%%%%%%%%%%%%%%%%%%%%%%%%%%%%%%%%%%%%%%%%%%%%%%%%%%%%%%%%%%%%%%%%%%%%%%%%%%%%%%%%%%
%%%%%%%%%%%%%%%%%%%%%%%%%%%%%%%%%%%%%%%%%%%%%%%%%%%%%%%%%%%%%%%%%%%%%%%%%%%%%%%%%%%%%%%%%%%%%%%%%%%%%%%%%%%%%%%%%%%%%%%%%%%%%%%%%%%%%%%%%%%%%%%%%%%%%%%%%%%%%%%%%%%%%%%%%%%%%%%%%%%%%%%%%%%%%%%%%%%%%%%%%%%%%%%%%%%%%%%%%%%%%%%%%%%%%%
%%%%%%%%%%%%%%%%%%%%%%%%%%%%%%%%%%%%%%%%%%%%%%%%%%%%%%%%%%%%%%%%%%%%%%%%%%%%%%%%%%%%%%%%%%%%%%%%%%%%%%%%%%%%%%%%%%%%%%%%%%%%%%%%%%%%%%%%%%%%%%%%%%%%%%%%%%%%%%%%%%%%%%%%%%%%%%%%%%%%%%%%%%%%%%%%%%%%%%%%%%
\newcommand{\Dem}{\protect\tikz[overlay,yshift=-4pt]\protect\draw[line width=1.5pt,line cap=round,color=gray] (0,0)[out=14,in=180] to node [shift=(90:6.pt),color=black] {\large\bf Dem:} (7:1); \qquad \quad }
\DeclareRobustCommand{\michi}{{\mathpalette\irchi\relax}}
\newcommand{\irchi}[1]{\raisebox{\depth}{\large$#1\chi$}} % inner command, used by \rchi



%%%%%%%%%%%%%%%%%%%%%%%%%%%%%%%%%%%%%%%%%%%%%%%%%%%%%%%%%%%%%%%%%%%%%%%%%%%%%%%%%%%%%%%%%%%%%%%%%%%%%%%%%%%%%%%%%%%%%%%%%%%%%%%%%%%%%%%%%%%%%%%%%%%%%%%%%%%%%%%%%%%%%%%%%%%%%%%%%%%%%%%%%%%%%%%%%%%%%%%%%%%%%%%%%%%%%%%%%%%%%%%%%%%%%%%%%%%%%%%%%%%%%%%%%%%%%%%%%%%%%%%%%%%%%%%%%%%%%%%%%%%%%%%%%%%%%%%%%%%%%%%%%%%%%%%%%%%%%%%%%%%%%%%%%%%%%%%%%%%%%%%%%%%%%%%%%%%%%%%%%%%%%%%%%%%%%%%%%%%%%%%%%%%%%%%%%%%%%%%%%%%%%%%%%%%%%%%%%%%%%%%%%%%%%%%%%%%%%%%%%%%%%%%%%%%%%%%%%%%%%%%%%%%%%%%%%%%%%%%%%%%%%%%%%%%%%%%%%%%%%%%%%%%%%%%%
% comando aclaraciones  

\newcommand{\aclaracion}[2]{

\begin{center}
\definecolor{gris}{rgb}{0.5,0.5,0.5}
\begin{tikzpicture}[line cap=round,line join=round,scale=1]
%\draw [color=gris,dash pattern=on 1pt off 1pt, xstep=1cm,ystep=1cm] (0.1,0.1) grid (6.3,6.3);
%\draw[color=black] (0.,0.) -- (\linewidth,0.);
\node at (.23,{.9*#2}) [color=black,above right] {\bf #1};
\foreach \y in {1,2,...,#2}
\draw[shift={(0,{.9*\y})},color=gris,opacity=.5] (0.,0.) -- ({.99*\linewidth},0.);
\end{tikzpicture} 
\end{center}
}





%%%%%%%%%%%%%%%%%%%%%%%%%%%%%%%%%%%%%%%%%%%%%%%%%%%%%%%%%%%%%%%%%%%%%%%%%%%%%%%%%%%%%%%%%%%%%%%%%%%%%%%%%%%%%%%%%%%%%%%%%%%%%%%%%%%%%%%%%%%%%%%%%%%%%%%%%%%%%%%%%%%%%%%%%%%%%%%%%%%%%%%%%%%%%%%%%%%%%%%%%%%%%%%%%%%%%%%%%%%%%%%%%%%%%%%%%%%%%%%%%%%%%%%%%%%%%%%%%%%%%%%%%%%%%%%%%%%%%%%%%%%%%%%%%%%%%%%%%%%%%%%%%%%%%%%%%%%%%%%%%%%%%%%%%%%%%%%%%%%%%%%%%%%%%%%%%%%%%%%%%%%%%%%%%%%%%%%%%%%%%%%%%%%%%%%%%%%%%%%%%%%%%%%%%%%%%%%%%%%%%%%%%%%%%%%%%%%%%%%%%%%%%%%%%%%%%%%%%%%%%%%%%%%%%%%%%%%%%%%%%%%%%%%%%%%%%%%%%%%%%%%%%%%%%%%%
% comando respuesta 

\newcommand{\res}[1]{

\begin{center}
\definecolor{gris}{rgb}{0.5,0.5,0.5}
\begin{tikzpicture}[line cap=round,line join=round,scale=1]
%\draw [color=gris,dash pattern=on 1pt off 1pt, xstep=1cm,ystep=1cm] (0.1,0.1) grid (6.3,6.3);
%\draw[color=black] (0.,0.) -- (\linewidth,0.);
\node at (.23,{.9*#1}) [color=black,above right] {\bf respuesta:};
\foreach \y in {1,2,...,#1}
\draw[shift={(0,{.9*\y})},color=gris,opacity=.5] (0.,0.) -- ({.99*\linewidth},0.);
\end{tikzpicture} 
\end{center}
}

%%%%%%%%%%%%%%%%%%%%%%%%%%%%%%%%%%%%%%%%%%%%%%%%%%%%%%%%%%%%%%%%%%%%%%%%%%%%%%%%%%%%%%%%%%%%%%%%%%%%%%%%%%%%%%%%%%%%%%%%%%%%%%%%%%%%%%%%%%%%%%%%%%%%%%%%%%%%%%%%%%%%%%%%%%%%%%%%%%%%%%%%%%%%%%%%%%%%%%%%%%%%%%%%%%%%%%%%%%%%%%%%%%%%%%%%%%%%%%%%%%%%%%%%%%%%%%%%%%%%%%%%%%%%%%%%%%%%%%%%%%%%%%%%%%%%%%%%%%%%%%%%%%%%%%%%%%%%%%%%%%%%%%%%%%%%%%%%%%%%%%%%%%%%%%%%%%%%%%%%%%%%%%%%%%%%%%%%%%%%%%%%%%%%%%%%%%%%%%%%%%%%%%%%%%%%%%%%%%%%%%%%%%%%%%%%%%%%%%%%%%%%%%%%%%%%%%%%%%%%%%%%%%%%%%%%%%%%%%%%%%%%%%%%%%%%%%%%%%%%%%%%%%%%%%%%
% comando ejemplo para completar 

\newcommand{\ejparacompletar}[1]{

\begin{center}
\definecolor{gris}{rgb}{0.5,0.5,0.5}
\begin{tikzpicture}[line cap=round,line join=round,scale=1]
%\draw [color=gris,dash pattern=on 1pt off 1pt, xstep=1cm,ystep=1cm] (0.1,0.1) grid (6.3,6.3);
%\draw[color=black] (0.,0.) -- (\linewidth,0.);
\node at (.23,{.9*#1}) [color=black,above right] {\bf Ejemplo:};
\foreach \y in {1,2,...,#1}
\draw[shift={(0,{.9*\y})},color=gris,opacity=.5] (0.,0.) -- ({.99*\linewidth},0.);
\end{tikzpicture} 
\end{center}
}


%%%%%%%%%%%%%%%%%%%%%%%%%%%%%%%%%%%%%%%%%%%%%%%%%%%%%%%%%%%%%%%%%%%%%%%%%%%%%%%%%%%%%%%%%%%%%%%%%%%%%%%%%%%%%%%%%%%%%%%%%%%%%%%%%%%%%%%%%%%%%%%%%%%%%%%%%%%%%%%%%%%%%%%%%%%%%%%%%%%%%%%%%%%%%%%%%%%%%%%%%%%%%%%%%%%%%%%%%%%%%%%%%%%%%%%%%%%%%%%%%%%%%%%%%%%%%%%%%%%%%%%%%%%%%%%%%%%%%%%%%%%%%%%%%%%%%%%%%%%%%%%%%%%%%%%%%%%%%%%%%%%%%%%%%%%%%%%%%%%%%%%%%%%%%%%%%%%%%%%%%%%%%%%%%%%%%%%%%%%%%%%%%%%%%%%%%%%%%%%%%%%%%%%%%%%%%%%%%%%%%%%%%%%%%%%%%%%%%%%%%%%%%%%%%%%%%%%%%%%%%%%%%%%%%%%%%%%%%%%%%%%%%%%%%%%%%%%%%%%%%%%%%%%%%%%%
% comando conclusión 

\newcommand{\conclus}[1]{

\begin{center}
\definecolor{gris}{rgb}{0.5,0.5,0.5}
\begin{tikzpicture}[line cap=round,line join=round,scale=1]
%\draw [color=gris,dash pattern=on 1pt off 1pt, xstep=1cm,ystep=1cm] (0.1,0.1) grid (6.3,6.3);
%\draw[color=black] (0.,0.) -- (\linewidth,0.);
\node at (.23,{.9*#1}) [color=black,above right] {\bf Conclusión:};
\foreach \y in {1,2,...,#1}
\draw[shift={(0,{.9*\y})},color=gris,opacity=.5] (0.,0.) -- ({.99*\linewidth},0.);
\end{tikzpicture} 
\end{center}
}

%%%%%%%%%%%%%%%%%%%%%%%%%%%%%%%%%%%%%%%%%%%%%%%%%%%%%%%%%%%%%%%%%%%%%%%%%%%%%%%%%%%%%%%%%%%%%%%%%%%%%%%%%%%%%%%%%%%%%%%%%%%%%%%%%%%%%%%%%%%%%%%%%%%%%%%%%%%%%%%%%%%%%%%%%%%%%%%%%%%%%%%%%%%%%%%%%%%%%%%%%%%%%%%%%%%%%%%%%%%%%%%%%%%%%%%%%%%%%%%%%%%%%%%%%%%%%%%%%%%%%%%%%%%%%%%%%%%%%%%%%%%%%%%%%%%%%%%%%%%%%%%%%%%%%%%%%%%%%%%%%%%%%%%%%%%%%%%%%%%%%%%%%%%%%%%%%%%%%%%%%%%%%%%%%%%%%%%%%%%%%%%%%%%%%%%%%%%%%%%%%%%%%%%%%%%%%%%%%%%%%%%%%%%%%%%%%%%%%%%%%%%%%%%%%%%%%%%%%%%%%%%%%%%%%%%%%%%%%%%%%%%%%%%%%%%%%%%%%%%%%%%%%%%%%%%%
% boton calculadora 

\newcommand{\boton}[2]{
\tikzpicture{  \node[draw,color=#1,fill=#1,text=white,,rounded corners=3mm,text width=7mm,text height =3mm,align=center,midway]  at (0,0) { \large \bf #2 };
}
\endtikzpicture
}

\usepackage[framemethod=TikZ]{mdframed}%Paquete para estilos de los bordes de los ambientes Importante este paquete (mdframed) define un entorno  que trata automáticamente los saltos de página en el texto enmarcado.
            \mdfsetup{skipabove=2\baselineskip,skipbelow=2\baselineskip,frametitlefont=\sffamily\bfseries, needspace=4\baselineskip, splittopskip=1.5\baselineskip} 
            \mdfsetup{apptotikzsetting={\tikzset{mdfbackground/.append style={draw=none}}}} 
            
            
%%%%%%%%%%%%%%%%%%%%%%%%%%%%%%%%%%%%%%%%%%%%%%%%%%%%%%%%%%%%%%
%
%Estilo  para cajas
%
%%%%%%%%%%%%%%%%%%%%%%%%%%%%%%%%%%%%%%%%%%%%%%%%%%%%%%%%%%%%%%

\mdfdefinestyle{teoSty}{backgroundcolor=blue!10, linecolor=blue, linewidth=3pt,
skipabove=4pt,
skipbelow=3pt, 
innerleftmargin=3ex, innerrightmargin=3ex, innertopmargin=3mm, innerbottommargin=3mm, innermargin=15pt, outermargin=-15pt, needspace={0.2\textheight},roundcorner=10pt}

\mdfdefinestyle{ObsSty}{backgroundcolor=red!10, linecolor=red, linewidth=2pt,
skipabove=4pt,
skipbelow=3pt, 
innerleftmargin=3ex, innerrightmargin=3ex, innertopmargin=3mm, innerbottommargin=3mm, innermargin=5pt, outermargin=-5pt, needspace={0.2\textheight},roundcorner=10pt}

\mdfdefinestyle{EjerSty}{backgroundcolor=blue!5!green!5!, linecolor=blue!50!green!50!, linewidth=2pt,
skipabove=4pt,
skipbelow=3pt, 
innerleftmargin=3ex, innerrightmargin=3ex, innertopmargin=3mm, innerbottommargin=3mm, innermargin=5pt, outermargin=-5pt, needspace={0.2\textheight},roundcorner=10pt}



%%%%%%%%%%%%%%%%%%%%%%%%%%%%%%%%%%%%%%%%%%%%%%%%%%%%%%%%%%%%%%
%
%Estilo  para lineas de relleno usando el paquete patterns y la libreria  patterns.meta
%
%%%%%%%%%%%%%%%%%%%%%%%%%%%%%%%%%%%%%%%%%%%%%%%%%%%%%%%%%%%%%%

\tikzdeclarepattern{
  name=mylines,
  parameters={
      \pgfkeysvalueof{/pgf/pattern keys/size},
      \pgfkeysvalueof{/pgf/pattern keys/angle},
      \pgfkeysvalueof{/pgf/pattern keys/line width},
  },
  bounding box={
    (0,-0.5*\pgfkeysvalueof{/pgf/pattern keys/line width}) and
    (\pgfkeysvalueof{/pgf/pattern keys/size},
0.5*\pgfkeysvalueof{/pgf/pattern keys/line width})},
  tile size={(\pgfkeysvalueof{/pgf/pattern keys/size},
\pgfkeysvalueof{/pgf/pattern keys/size})},
  tile transformation={rotate=\pgfkeysvalueof{/pgf/pattern keys/angle}},
  defaults={
    size/.initial=5pt,
    angle/.initial=45,
    line width/.initial=.4pt,
  },
  code={
      \draw [line width=\pgfkeysvalueof{/pgf/pattern keys/line width}]
        (0,0) -- (\pgfkeysvalueof{/pgf/pattern keys/size},0);
  },
}





\newcommand{\semejante}[1]{\stackrel[#1]{}{\longleftrightarrow}}% \thicksim para decir la Semejanza de matrices la variable #1  indica la operacion elemental 
\newcommand{\OptipoUno}[1]{\stackrel[#1]{}{=}}



\newcommand{\indicador}[4]{
{\tikz[remember picture,overlay]\coordinate (#1) at #2;}{\tikz[remember picture,overlay]\coordinate (#3) at #4;} 
}

\newcommand{\marka}[2]{
{\tikz[remember picture,overlay]\coordinate (#1) at #2;}
}




\newcommand{\cof}{ \mbox{cof}}

\newcommand{\paralela}{\rotatebox[origin=c]{-10}{\Large$\parallel$}}

\usepackage{fancyhdr}
\fancypagestyle{miestilo}{
\fancyhead[L]{}
\fancyhead[C]{}
\fancyhead[R]{}
\fancyfoot[L]{}
\fancyfoot[C]{}
\fancyfoot[R]{pág \thepage}
\renewcommand{\headrulewidth}{0pt}
\renewcommand{\footrulewidth}{0pt}
}
%%%%%%%%%%%%
%%%%%%%%%%%%
\renewcommand{\footrulewidth}{1pt}
\pagestyle{fancy}\markboth{}{}
\renewcommand{\headrulewidth}{.05pt}
\renewcommand{\headrule}{{\color{gray}%
\hrule width\headwidth height\headrulewidth \vskip-\headrulewidth}}
\renewcommand{\footrule}{{\color{gray}%
\vskip-\footruleskip\vskip-\footrulewidth
\hrule width\headwidth height\footrulewidth\vskip\footruleskip}}
\renewcommand{\footrulewidth}{.50pt}
\fancyhead[L]{\color{gray}{\footnotesize \sc ALGEBRA Y GEOMETRIA I} }
\fancyhead[C]{\color{gray}{\sc \footnotesize Primer cuatrimestre 2025}}
\fancyhead[R]{\color{gray}{\sc \footnotesize  Universidad Nacional Del Comahue}}
\fancyfoot[R]{\color{gray}{pág. \thepage}}
\fancyfoot[C]{\color{gray}{\sc \footnotesize Ingenieria en petroleo}}
\setlength{\headheight}{15pt}


%\usepackage{wrapfig}

\newcommand{\co}{\mathbb{C}}
%\newcommand{\re}{\mathbb{R}}
\newcommand{\q}{\mathbb{Q}}
\newcommand{\z}{\mathbb{Z}}
\newcommand{\n}{\mathbb{N}}
\newcommand{\ve}{\mathbb{V}}
\newcommand{\w}{\mathbb{W}}
\newcommand{\be}{{\cal B}}
\newcommand{\ca}{{\cal C}}
\newcommand{\pqx}{\mathbb{Q}[x]}
\newcommand{\prx}{\mathbb{R}[x]}
\newcommand{\pcx}{\mathbb{C}[x]}
%\newcommand{\izg}{\big\langle}
%\newcommand{\rangle}{\big\rangle}
%%%%%%%%%%%%%%%%%%%%%%%%%%%%%5

%%%%%%%%%%%%%%%%%%%%%%%%%%%%%%%%%%%%5
\newcommand{\materias}{\sc ALGEBRA Y GEOMETRIA I}%separadas con '&'
%%%%%%%%%%%%%%%%%%%%%%%%%%%%%%%%%%%%%%%
% se define el encabezado, 
% el parametro es el texto que va despues de: trabajo practico $N^\circ$... (agregar numero y titulo)
%
%%%%%%%%%%%%%%%%%%%%%%%%%%%%%%%%%\newcommand{\encabezado}[1]{%
\fancyfoot[L]{\vspace*{-3.5mm}\color{gray}{\small T. P. 4: Vectores}}
\thispagestyle{miestilo} 
\newcommand{\encabezado}[1]{%
\fancyfoot[L]{\vspace*{-3.5mm}\color{gray}{\small T. P. #1}}
%\thispagestyle{miestilo}%
%--------------------logodpto //texto // logouni------------------
%--------------------logodpto------------------

\vspace*{-18mm}\hspace*{-11mm} \begin{minipage}{2.7cm}
\begin{center} 
     \includegraphics[width=2.7cm, height=2.7cm]{logodpto}
\end{center}  
\end{minipage}\hspace*{-15mm} \begin{minipage}{15.9cm}
\begin{center}
\begin{tabular}{c @{ - } c}
 \materias \\
\carreras
\end{tabular}

\cuatri

\medskip

TRABAJO PRÁCTICO N$^\circ$#1 

\end{center}

\end{minipage}\hspace{-14mm} \begin{minipage}{2.5cm}
  \begin{center} 
    \includegraphics[width=2.5cm, height=2.5cm]{logouni}
  \end{center}  
\end{minipage}
%-----------------------------------------------------

%----------------------------------------------------------------
\hrule
%---------------------------------------------------------------
}%%



\pagestyle{fancy}
\def\identacion{-2.5mm}%mueve la identacion de la enumeracion de los ejercicios.


\usepackage{asypictureB}
\begin{asyheader}
settings.outformat="png";
settings.render=8;
//import graph3;
//import smoothcontour3;

pen bpen = rgb(0.75, 0.07, 0.01);
//material m = material(diffusepen=0.7bpen, ambientpen=bpen+opacity(0.8), emissivepen=0.3*bpen,specularpen=0.999white, shininess=1.0);      
\end{asyheader}


\usepackage{multicol}
\setlength{\columnsep}{7mm}





\newcommand{\Or}{%
    \tikz[scale=.17%baseline=-0.75ex, 
    ]{ % Cambia "scale" para ajustar el tamaño
        \draw[thick] plot[smooth cycle, tension=1.2] coordinates {(-0.4, 0) (0, -0.6) (0.4, 0) (0, 0.6)};
        \draw[thick,domain=pi:2*pi,samples=200] 
        plot ({0.1*\x * cos(\x r)}, {0.07*\x * sin(\x r)+0.45});
        %\draw[thick,domain=1.5*pi:2*pi,samples=200] plot ({0.08*\x * cos(\x r)-0.9}, {0.09*\x * sin(\x r)-0.1});
    }%
    \,\, 
}

\newcommand{\Noo}{%Es una cruz de que esa solución se descarta
 \tikz[overlay,remember picture] { 
  %%%% cruz en (-, 0)
\draw[red,thick,yshift=2mm] (-.15,\h-.25) -- (+.15,\h+.25); % Línea diagonal ascendente
\draw[red,thick,yshift=2mm] (-.15,\h+.25) -- (+.15,\h-.25); % Línea diagonal descendente
}
}






\begin{document}


\setlength{\headheight}{15pt}

\vspace*{-2.5cm}
%--------------------logodpto------------------
\begin{figure}[htbp]
  \begin{minipage}[r][][t]{0.15\textwidth}
\begin{center} 
 % Universidad Nacional del Comahue.  
 \includegraphics[width=2.2cm, height=2.2cm]{imagen/inglogo.jpg}
\vspace{1.8cm}
\end{center}  
\end{minipage}\quad
% texto central
\begin{minipage}[r][][t]{0.31\textwidth}
\begin{center}
\vspace*{-19mm}
{\bf Universidad Nacional
del Comahue} \\
Facultad de Ingeniería
\end{center}
\end{minipage}\quad\begin{minipage}[r][][t]{0.31\textwidth}
\begin{center}
\vspace*{-20mm}
{\bf  Álgebra y Geometría I}
Ingeniería en Petróleo
Primer Cuatrimestre 2024
\end{center}
\end{minipage}\quad
%----------------------logouni-------------------------
\begin{minipage}[r][][t]{0.15\textwidth}
  \begin{center} 
     \includegraphics[width=2.2cm, height=2.2cm]{imagen/Logo_UNco.jpg}
     \vspace{1.8cm}
  \end{center}  
\end{minipage}
%-----------------------------------------------------
\end{figure}
\vspace*{-22mm}

\begin{center}
{\bf \Large{Resolución Trabajo Práctico 4: Vectores}}
\end{center}


\begin{enumerate}[\hspace{-5mm} 1) ]
\item Dado el siguiente gráfico y considerando los elementos marcados: vector $\overrightarrow{PQ}$ y los puntos $A$ y $B$.  







\begin{center}
    \vspace{-4mm} 
    
\definecolor{rojo}{rgb}{0.7,0.4,0.4}
\definecolor{azul}{rgb}{.4,.4,0.9}
 
\definecolor{gris}{rgb}{0.5,0.5,0.5}
\begin{tikzpicture}[line cap=round,line join=round,>=triangle 45,scale=.8]

% extremos para los ejes
\xdef\ext{6}
\xdef\f{-2}
\pgfmathtruncatemacro{\c}{\f+1}
\draw[->] (\f-.3,0) -- (12+.3,0) ;

\draw[->] (0,\f-.3) -- (0,\ext+.3) ;
\draw [color=gris!30!,dash pattern=on 1pt off 1pt, xstep=1cm,ystep=1cm] (\f-.3,\f-.3) grid (12.3,\ext+.3);

\foreach \x in {\f,...,-1,1,2,...,12}
\draw[shift={(\x,0)},color=black] (0pt,2pt) -- (0pt,-2pt) node[below] {\tiny $\x$};

\foreach \y in {\f,...,-1,1,2,...,\ext }
\draw[shift={(0,\y)},color=black] (2pt,0pt) -- (-2pt,0pt) node[left] {\tiny $\y$};


%\clip(-3.2,-1.3) rectangle (3.45,4.3);

\coordinate (A) at (6,-2);
\coordinate (B) at (2,4);
\coordinate (P) at (4,2);
\coordinate (Q) at (7,4); 
\xdef\ang{80}


%%%%%Grafico de los puntos
\draw[fill=black] (P) circle (1.5pt) node[shift=(-\ang: 5pt),right] {\scriptsize $P=(4,2)$};

\draw[fill=black] (Q) circle (1.5pt) node[shift=(\ang: 5pt),right] {\scriptsize $Q=(7,4)$};

\draw[fill=black] (A) circle (1.5pt) node[shift=(\ang: 5pt),below right] {\scriptsize $A=(6,-2)$};

\draw[fill=black] (B) circle (1.5pt) node[shift=(\ang: 5pt),right] {\scriptsize $B=(2,4)$};


%vectores
\draw [->,line width=1.5pt, color=blue] (P) -- (Q);

\end{tikzpicture}  

\vspace{-4mm}
\end{center}  

\begin{enumerate}
    \item Encontrar gráficamente:
   \begin{enumerate}[\hspace{-4mm} i)]
        \item Un vector $\vec{v}$ equivalente al vector $\overrightarrow{PQ}$ con origen en el punto $(0, 0)$.
        \item Un vector, con origen en el punto $A$, que tenga igual dirección y el doble de longitud que el vector $\vec{v}$. ¿Es único? ¿Por qué?
        \item Un vector $\vec{z}$, con origen en el punto $B$, de igual dirección, sentido opuesto y la mitad de la longitud que $\vec{v}$.
        \item Un vector simétrico al vector $\vec{v}$ con respecto al eje $x$ y otro vector simétrico con respecto al eje $y$.
    \end{enumerate}
\item  Encontrar analíticamente:
\begin{enumerate}[\hspace{-4mm} i)]
    \item  El vector $\vec{v}$ hallado en el inciso a)i). ¿qué representa en la teoría de vectores? Indicar sus elementos.
    \item  Las componentes del vector con extremo inicial en el punto $B$ y extremo final en el punto $P$.
\end{enumerate}
\end{enumerate}


{\bf Solución:}


\begin{enumerate}
    \item \begin{enumerate}[i)]
    \begin{multicols}{7}
        \item \, 
        \item \,
        \item \,
        \item \, 
    \end{multicols}
    \end{enumerate}

\begin{center}


\definecolor{rojo}{rgb}{0.7,0.4,0.4}
\definecolor{azul}{rgb}{.4,.4,0.9}
 
\definecolor{gris}{rgb}{0.5,0.5,0.5}
\begin{tikzpicture}[line cap=round,line join=round,>=triangle 45,scale=.9]
% Configuración de los ejes
\xdef\ext{6}
\xdef\f{-2}
\pgfmathtruncatemacro{\c}{\f+1}
\draw[->] (\f-.3,0) -- (12+.3,0);
\draw[->] (0,\f-.3) -- (0,\ext+.3);
\draw [color=gris!30!,dash pattern=on 1pt off 1pt, xstep=1cm,ystep=1cm] (\f-.3,\f-.3) grid (12.3,\ext+.3);

% Etiquetas de los ejes
\foreach \x in {\f,...,-1,1,2,...,12}
\draw[shift={(\x,0)},color=black] (0pt,2pt) -- (0pt,-2pt) node[below] {\tiny $\x$};

\foreach \y in {\f,...,-1,1,2,...,\ext }
\draw[shift={(0,\y)},color=black] (2pt,0pt) -- (-2pt,0pt) node[left] {\tiny $\y$};

% Coordenadas
\coordinate (A) at (6,-2);
\coordinate (B) at (2,4);
\coordinate (P) at (4,2);
\coordinate (Q) at (7,4);



% Puntos y vectores
\draw[fill=black] (P) circle (1.5pt) node[above right] {\scriptsize $P=(4,2)$};
\draw[fill=black] (Q) circle (1.5pt) node[above right] {\scriptsize $Q=(7,4)$};
\draw[fill=black] (A) circle (1.5pt) node[below right] {\scriptsize $A=(6,-2)$};
\draw[fill=black] (B) circle (1.5pt) node[above right] {\scriptsize $B=(2,4)$};

% Vector PQ
\draw[->, line width=1.5pt, color=azul] (P) -- (Q) node[midway,above] {\scriptsize $\overrightarrow{PQ}$};



%Respuestas --------
\coordinate (V) at (3,2); 
 \coordinate (UU) at ($(A)+2*(V)$);
 \coordinate (ZZ) at ($(B)-.5*(V)$);
 \coordinate (Vx) at (3,-2); 
 \coordinate (Vy) at (-3,2); 
% Vector v
\draw[->, line width=1.5pt, color=green] (0,0) -- (V) node[midway,above] {\scriptsize $\overrightarrow{\vec{v}}$};
% Vector ejercicio b)
\draw[->, line width=1.5pt, color=red] (A) -- (UU) node[midway,above] {\scriptsize $\overrightarrow{\vec{u}}$};
% Vector ejercicio b)
\draw[->, line width=1.5pt, color=teal] (B) -- (ZZ) node[midway,above] {\scriptsize $\vec{z}$};
% Vector v_x
\draw[->, line width=1.5pt, color=orange] (0,0) -- (Vx) node[midway,above] {\scriptsize $\overrightarrow{\vec{v}}_x$};
% Vector v_y
\draw[->, line width=1.5pt, color=cyan] (0,0) -- (Vy) node[midway,above] {\scriptsize $\overrightarrow{\vec{v}}_y$};

\end{tikzpicture}
\end{center}





\item  Encontrar analíticamente:
\begin{enumerate}[\hspace{-4mm} i)]
    \item  El vector $\vec{v}$ hallado en el inciso a)i). ¿qué representa en la teoría de vectores? Indicar sus elementos.


 El vector $\overrightarrow{PQ}$ tiene componentes:
    \[
    \overrightarrow{PQ}  = \overrightarrow{OQ} -  \overrightarrow{OP} = (7,4)-(4,2) = (7-4, 4-2) = (3,2).
    \]
    Entonces, el vector $\vec{v}$ con origen en $(0,0)$ es:
    \[
    \vec{v} = (3,2)
    \]
es el vector libre, es decir es el representante a todos los vectores equivalentes al vector $\overrightarrow{PQ}$.









    
    \item  Las componentes del vector con extremo inicial en el punto $B$ y extremo final en el punto $P$.


    \[
\overrightarrow{BP} = \overrightarrow{OP} -  \overrightarrow{OB}  = (4,2) - (2,4) = (4 - 2,\ 2 - 4) = (2,\ -2).
\]

Por lo tanto, las componentes del vector $\vec{BP}$ son:

\[
\overrightarrow{BP} = (2,\ -2).
\]




\end{enumerate}
\end{enumerate}





    


    
  

\item Dados los vectores pertenecientes a $\mathbb{R}^2$, $\vec{u} = (6, 4)$, $\vec{v} = (-4, 2)$ y $\vec{w} = (-3, -2)$:

    \begin{enumerate}[ a)] 
        \item Representarlos en un mismo sistema de ejes cartesianos.
        \item Calcular analítica y gráficamente las siguientes operaciones:
            \begin{enumerate}[ i)]
            \begin{minipage}{.1\textwidth} 
                \item $\vec{u} + \vec{v}$.
            \end{minipage}\quad\begin{minipage}{.15\textwidth}
                \item $\tfrac{1}{2} \left(\vec{v} + \vec{u}\right)$.
            \end{minipage}\quad\begin{minipage}{.15\textwidth}
                \item $-\vec{v} + \vec{w}$.
            \end{minipage}\quad\begin{minipage}{.13\textwidth}
                \item $ \vec{u} - \tfrac{1}{2}\vec{v} $. 
            \end{minipage}\quad\begin{minipage}{.2\textwidth}
                \item $\vec{u} + 3\vec{v} - \vec{w}$.
            \end{minipage}\quad\begin{minipage}{.2\textwidth}
                \item $\frac{1}{2} \vec{v}  + \tfrac{1}{2} \vec{u}$
            \end{minipage} 
            \end{enumerate}
        \item ¿Qué propiedad relaciona los incisos {\scriptsize$II)$ }y {\scriptsize$VI)$}?
    \end{enumerate}





\textbf{Solución:}

\textbf{Enunciado:} Dados los vectores pertenecientes a $\mathbb{R}^2$, $\vec{u} = (6, 4)$, $\vec{v} = (-4, 2)$ y $\vec{w} = (-3, -2)$:

\begin{enumerate}[a) ]
    \item Representarlos en un mismo sistema de ejes cartesianos. 


\definecolor{rojo}{rgb}{0.7,0.4,0.4}
\definecolor{azul}{rgb}{.4,.4,0.9}
 
\definecolor{gris}{rgb}{0.5,0.5,0.5}
\begin{center}
\begin{tikzpicture}[line cap=round,line join=round,>=stealth,scale=.9]
% extremos para los ejes
\xdef\ext{7}
\xdef\f{-5}
\xdef\exty{4}
\xdef\fy{-2}
\pgfmathtruncatemacro{\c}{\f+1} 
\pgfmathtruncatemacro{\d}{\fy+1} 
% Configuración de los ejes

\draw [color=gris!30!,very thin] (-5.5,-2.5) grid (7.5,4.5);
\draw[->] (-5.5,0) -- (7.5,0) node[below] {$x$};
\draw[->] (0,-2.5) -- (0,4.5) node[left] {$y$};

% Etiquetas de los ejes
\foreach \x in {\f,...,-1,1,2,...,\ext}{
\draw[shift={(\x,0)},color=black] (0pt,2pt) -- (0pt,-2pt) node[below] {\tiny $\x$};}

\foreach \y in {\fy,...,-1,1,2,...,\exty }{
\draw[shift={(0,\y)},color=black] (2pt,0pt) -- (-2pt,0pt) node[left] {\tiny $\y$};}

%vectores
\coordinate (U) at (6,4);
\coordinate (V) at (-4,2);
\coordinate (W) at (-3,-2); 

% Vector u
\draw[->, thick, color=azul] (0,0) -- (U) node[midway, above right] {\scriptsize $\vec{u}$};

% Vector v
\draw[->, thick, color=rojo] (0,0) -- (V) node[midway, above left] {\scriptsize $\vec{v}$};

% Vector w
\draw[->, thick, color=green] (0,0) -- (W) node[midway, below left] {\scriptsize $\vec{w}$};

\end{tikzpicture}
\end{center}
 
    \item Calcular analítica y gráficamente las siguientes operaciones:
 

\begin{enumerate}[i)]
    \item $\vec{u} + \vec{v}$.  %$\vec{u} = (6, 4)$, $\vec{v} = (-4, 2)$ y $\vec{w} = (-3, -2)$:
    \[
    \vec{u} + \vec{v} = (6, 4) + (-4, 2) = (6 - 4 , 4 + 2) = (2, 6).
    \]
    Gráficamente:
    \begin{center}
    \definecolor{rojo}{rgb}{0.7,0.4,0.4}
\definecolor{azul}{rgb}{.4,.4,0.9}
 
\definecolor{gris}{rgb}{0.5,0.5,0.5}
\begin{tikzpicture}[line cap=round,line join=round,>=stealth,scale=.5]
% extremos para los ejes
\xdef\ext{7}
\xdef\f{-5}
\xdef\exty{6}
\xdef\fy{-1}
\pgfmathtruncatemacro{\c}{\f+1} 
\pgfmathtruncatemacro{\d}{\fy+1} 
% Configuración de los ejes

\draw [color=gris!30!,very thin] ({\f-.3},{\fy-.3}) grid ({\ext+.5},{\exty+.5});
\draw[->] ({\f-.3},0) -- ({\ext+.5},0) node[below] {$x$};
\draw[->] (0,{\fy-.3}) -- (0,{\exty+.5}) node[left] {$y$};

% Etiquetas de los ejes
\foreach \x in {\f,...,-1,1,2,...,\ext}{
\draw[shift={(\x,0)},color=black] (0pt,2pt) -- (0pt,-2pt) node[below] {\tiny $\x$};}

\foreach \y in {\fy,...,-1,1,2,...,\exty }{
\draw[shift={(0,\y)},color=black] (2pt,0pt) -- (-2pt,0pt) node[left] {\tiny $\y$};}

%vectores
\coordinate (U) at (6,4);
\coordinate (V) at (-4,2);
\coordinate (W) at (-3,-2); 
 \coordinate (Y) at ($(U)+(V)$); 

% Vector u
\draw[->, thick, color=azul] (0,0) -- (U) node[midway, above right] {\scriptsize $\vec{u}$};

% Vector v
\draw[->, thick, color=rojo] (0,0) -- (V) node[midway, above left] {\scriptsize $\vec{v}$};

 
    

     %Vector u
    \draw[->, thick, color=azul] (V) -- (Y) node[midway, below left] {\scriptsize $\vec{u}$};
    

    % Vector u + v
    \draw[->, thick, color=red] (0,0) -- (Y) node[midway, above] {\scriptsize $\vec{u} + \vec{v}$};
    \end{tikzpicture}
    \end{center}





    

    




    \item $\tfrac{1}{2}(\vec{v} + \vec{u})$: %$\vec{u} = (6, 4)$, $\vec{v} = (-4, 2)$ y $\vec{w} = (-3, -2)$:
    \[
    \tfrac{1}{2}(\vec{v} + \vec{u}) = \tfrac{1}{2}((-4, 2) + (6, 4)) = \tfrac{1}{2}(2, 6) = (1, 3).
    \]









\begin{center}
    \begin{tikzpicture}[line cap=round,line join=round,>=stealth,scale=.5]
    % extremos para los ejes
\xdef\ext{7}
\xdef\f{-5}
\xdef\exty{6}
\xdef\fy{-1}
\pgfmathtruncatemacro{\c}{\f+1} 
\pgfmathtruncatemacro{\d}{\fy+1} 
    % Configuración de los ejes% Configuración de los ejes
    \draw [color=gris!30!,very thin] ({\f-.3},{\fy-.3}) grid ({\ext+.5},{\exty+.5});
    \draw[->] ({\f-.3},0) -- ({\ext+.5},0) node[below] {$x$};
    \draw[->] (0,{\fy-.3}) -- (0,{\exty+.5}) node[left] {$y$};
    % Etiquetas de los ejes
\foreach \x in {\f,...,-1,1,2,...,\ext}{
\draw[shift={(\x,0)},color=black] (0pt,2pt) -- (0pt,-2pt) node[below] {\tiny $\x$};}

\foreach \y in {\fy,...,-1,1,2,...,\exty }{
\draw[shift={(0,\y)},color=black] (2pt,0pt) -- (-2pt,0pt) node[left] {\tiny $\y$};}


    %vectores
\coordinate (U) at (6,4);
\coordinate (V) at (-4,2);
\coordinate (W) at (-3,-2); 
\coordinate (X) at ($.5*(U) + .5*(V)$);    
% Dibujar el paralelogramo
  
  \coordinate (F) at ($(U)+(V)$);  
  \draw[dashed,,rounded corners=0mm,opacity=.5] (U) -- (F) -- (V) ;
  \draw[dashed,,rounded corners=0mm,opacity=.5] (0,0) -- (F)  ;
  \draw[dashed,,rounded corners=0mm,opacity=.5] (U) -- (V) ;


    % Vector u
    \draw[->, thick, color=blue] (0,0) -- (U) node[midway, below left] {\scriptsize $\vec{u}$};

    % Vector v
    \draw[->, thick, color=green] (0,0) -- (V) node[midway, below left] {\scriptsize $\vec{v}$};
    
     %Vector tfrac{1}{2}(u+V
    \draw[->, thick, color=red] (0,0) -- (X) node[midway, below right] {\scriptsize $\tfrac{1}{2}(\vec{u}+ \vec{v})$};
 
    \end{tikzpicture}
    \end{center}


















    \item $-\vec{v} + \vec{w}$:%$\vec{u} = (6, 4)$, $\vec{v} = (-4, 2)$ y $\vec{w} = (-3, -2)$:
    \[
    -\vec{v} + \vec{w} = -(-4, 2) + (-3, -2) =(4, -2) + (-3, -2) = (1, -4).
    \]






\begin{center}
    \begin{tikzpicture}[line cap=round,line join=round,>=stealth,scale=.8]
    % extremos para los ejes
\xdef\ext{6}
\xdef\f{-8}
\xdef\exty{2}
\xdef\fy{-4}
\pgfmathtruncatemacro{\c}{\f+1} 
\pgfmathtruncatemacro{\d}{\fy+1} 
    % Configuración de los ejes
    \draw [color=gris!30!,very thin] (-8.5,-4.5) grid (6.5,2.5);
    \draw[->] (-8.5,0) -- (6.5,0) node[below] {$x$};
    \draw[->] (0,-4.5) -- (0,2.5) node[left] {$y$};
    % Etiquetas de los ejes
\foreach \x in {\f,...,-1,1,2,...,\ext}{
\draw[shift={(\x,0)},color=black] (0pt,2pt) -- (0pt,-2pt) node[below] {\tiny $\x$};}

\foreach \y in {\fy,...,-1,1,2,...,\exty }{
\draw[shift={(0,\y)},color=black] (2pt,0pt) -- (-2pt,0pt) node[left] {\tiny $\y$};}


    %vectores
\coordinate (U) at (6,4);
\coordinate (V) at (-4,2);
\coordinate (W) at (-3,-2); 
\coordinate (X) at ($-1*(V)$); 
 \coordinate (Y) at ($-1*(V) + (W)$);    
    % Vector W
    \draw[->, thick, color=teal] (0,0) -- (W) node[midway, above left] {\scriptsize $\vec{w}$};
    \draw[->, thick, color=teal] (X) -- (Y) node[midway, above left] {\scriptsize $\vec{w}$};

    % Vector v
    \draw[->, thick, color=red] (0,0) -- (V) node[midway, below left] {\scriptsize $\vec{v}$};
    
     %Vector tfrac{1}{2}u
    \draw[->, thick, color=rojo] (0,0) -- (X) node[midway, above right] {\scriptsize $- \vec{v}$};

    % Vector -v + u
    \draw[->, thick, color=blue] (0,0) -- (Y) node[midway, below left] {\scriptsize $-\vec{v} + \vec{w}$};
 
    \end{tikzpicture}
    \end{center}

























 



\item $\vec{u} -\tfrac{1}{2} \vec{v}$: %$\vec{u} = (6, 4)$, $\vec{v} = (-4, 2)$ y $\vec{w} = (-3, -2)$:
    \[
    \vec{u} -\tfrac{1}{2}\vec{v} = (6, 4) -\tfrac{1}{2} (-4, 2) = (6,4) - (-2, 1) = (8,3).
    \]




\begin{center}
    \begin{tikzpicture}[line cap=round,line join=round,>=stealth,scale=.9]
% extremos para los ejes
\xdef\ext{9}
\xdef\f{-5}
\xdef\exty{4}
\xdef\fy{-1}
\pgfmathtruncatemacro{\c}{\f+1} 
\pgfmathtruncatemacro{\d}{\fy+1} 
% Configuración de los ejes

\draw [color=gris!30!,very thin] ({\f-.3},{\fy-.3}) grid ({\ext+.5},{\exty+.5});
\draw[->] ({\f-.3},0) -- ({\ext+.5},0) node[below] {$x$};
\draw[->] (0,{\fy-.3}) -- (0,{\exty+.5}) node[left] {$y$};

% Etiquetas de los ejes
\foreach \x in {\f,...,-1,1,2,...,\ext}{
\draw[shift={(\x,0)},color=black] (0pt,2pt) -- (0pt,-2pt) node[below] {\tiny $\x$};}

\foreach \y in {\fy,...,-1,1,2,...,\exty }{
\draw[shift={(0,\y)},color=black] (2pt,0pt) -- (-2pt,0pt) node[left] {\tiny $\y$};}



    %vectores
\coordinate (U) at (6,4);
\coordinate (V) at (-4,2);
\coordinate (W) at (-3,-2); 
\coordinate (X) at ($-.5*(V)$); 
 \coordinate (Y) at ($(U)-.5*(V) $);    
    % Vector u
    \draw[->, thick, color=blue] (0,0) -- (U) node[midway, below left] {\scriptsize $\vec{u}$};

    % Vector v
    \draw[->, thick, color=red] (0,0) -- (V) node[midway, below left] {\scriptsize $\vec{v}$};
    
     %Vector tfrac{1}{2}u
    \draw[->, thick, color=red] (U) -- (Y) node[midway,above] {\scriptsize $-\tfrac{1}{2}\vec{u}$};

    % Vector \tfrac{1}{2} u
    \draw[->, thick, color=red] (0,0) -- (X) node[midway, below ] {\scriptsize $-\tfrac{1}{2}\vec{u}$};

    %% Vector \tfrac{1}{2} u
    \draw[->, thick, color=teal,line width=2pt] (0,0) -- (Y) node[midway, below right] {\scriptsize $\vec{u} -\tfrac{1}{2} \vec{v}$};
    \end{tikzpicture}
    \end{center}




   \item $\vec{u} + 3\vec{v} - \vec{w}$:%$\vec{u} = (6, 4)$, $\vec{v} = (-4, 2)$ y $\vec{w} = (-3, -2)$:
    \[
    \vec{u} + 3\vec{v} - \vec{w} = (6, 4) + 3(-4, 2) - (-3, -2) = (6, 4) + (-12, 6) + (3, 2) = (-3, 12).
    \] 



\definecolor{rojo}{rgb}{0.7,0.4,0.4}
\definecolor{azul}{rgb}{.4,.4,0.9}
 \definecolor{verde}{rgb}{0.4,0.7,0.4}

\definecolor{gris}{rgb}{0.5,0.5,0.5}
\begin{center}
\begin{tikzpicture}[line cap=round,line join=round,>=stealth,scale=.5]
% extremos para los ejes
\xdef\ext{7}
\xdef\f{-6}
\xdef\exty{10}
\xdef\fy{-2}
\pgfmathtruncatemacro{\c}{\f+1} 
\pgfmathtruncatemacro{\d}{\fy+1} 
% Configuración de los ejes

\draw [color=gris!30!,very thin] (-6.5,-2.5) grid (7.5,12.5);
\draw[->] (-6.5,0) -- (7.5,0) node[below] {$x$};
\draw[->] (0,-2.5) -- (0,12.5) node[left] {$y$};

% Etiquetas de los ejes
\foreach \x in {\f,...,-1,1,2,...,\ext}{
\draw[shift={(\x,0)},color=black] (0pt,2pt) -- (0pt,-2pt) node[below] {\tiny $\x$};}

\foreach \y in {\fy,...,-1,1,2,...,\exty }{
\draw[shift={(0,\y)},color=black] (2pt,0pt) -- (-2pt,0pt) node[left] {\tiny $\y$};}

%vectores
\coordinate (U) at (6,4);
\coordinate (V) at (-4,2);
\coordinate (W) at (-3,-2); 
\coordinate (X) at ($(U)+3*(V)$);
%\coordinate (Y) at ($$);
\coordinate (S) at ($(U)+3*(V)-(W)$);

% Vector u
\draw[->, thick, color=azul] (0,0) -- (U) node[midway, above,pos=.8 ] {\scriptsize $\vec{u}$};

% Vector v
\draw[->, thick, color=rojo] (0,0) -- (V) node[midway, above left] {\scriptsize $\vec{v}$};

% Vector 3v
\draw[->, thick, color=rojo,opacity=.5] (U) -- (X) node[midway, above ,pos=.8 ] {\scriptsize $3\vec{v}$};

% Vector w
\draw[->, thick, color= blue!30!green!90!] (0,0) -- (W) node[midway, above,pos=.8 ] {\scriptsize $\vec{w}$};

% Vector -w
\draw[->, thick, color=verde] (X) -- (S) node[midway, below,pos=.8 ] {\scriptsize $-\vec{w}$};

% Vector u + 3 v - w
\draw[->, thick, color=teal] (0,0) -- (S) node[midway, below left] {\scriptsize $\vec{u} + 3\vec{v} - \vec{w}$};

\end{tikzpicture}
\end{center}







    \item $\tfrac{1}{2}\vec{v} + \tfrac{1}{2}\vec{u}$: %$\vec{u} = (6, 4)$, $\vec{v} = (-4, 2)$ y $\vec{w} = (-3, -2)$:
    \[
    \tfrac{1}{2}\vec{v} + \tfrac{1}{2}\vec{u} = \tfrac{1}{2}(-4, 2) + \tfrac{1}{2}(6, 4) = (-2, 1) + (3, 2) = (1, 3).
    \]







\begin{center}
    \begin{tikzpicture}[line cap=round,line join=round,>=stealth,scale=.5]
    % extremos para los ejes
\xdef\ext{7}
\xdef\f{-5}
\xdef\exty{6}
\xdef\fy{-1}
\pgfmathtruncatemacro{\c}{\f+1} 
\pgfmathtruncatemacro{\d}{\fy+1} 
    % Configuración de los ejes% Configuración de los ejes
    \draw [color=gris!30!,very thin] ({\f-.3},{\fy-.3}) grid ({\ext+.5},{\exty+.5});
    \draw[->] ({\f-.3},0) -- ({\ext+.5},0) node[below] {$x$};
    \draw[->] (0,{\fy-.3}) -- (0,{\exty+.5}) node[left] {$y$};
    % Etiquetas de los ejes
\foreach \x in {\f,...,-1,1,2,...,\ext}{
\draw[shift={(\x,0)},color=black] (0pt,2pt) -- (0pt,-2pt) node[below] {\tiny $\x$};}

\foreach \y in {\fy,...,-1,1,2,...,\exty }{
\draw[shift={(0,\y)},color=black] (2pt,0pt) -- (-2pt,0pt) node[left] {\tiny $\y$};}



    %vectores
\coordinate (U) at (6,4);
\coordinate (V) at (-4,2);
\coordinate (W) at (-3,-2); 
\coordinate (A) at ($.5*(U) $); 
\coordinate (B) at ($.5*(V)$); 
\coordinate (X) at ($.5*(U) + .5*(V)$); 
% Dibujar el paralelogramo
   
  \draw[dashed,,rounded corners=0mm,opacity=.5] (A) -- (X) -- (B) ;
  

    % Vector u
    \draw[->, thick, color=blue] (0,0) -- (U) node[midway, below left] {\scriptsize $\vec{u}$};

    % Vector v
    \draw[->, thick, color=green] (0,0) -- (V) node[midway, below left] {\scriptsize $\vec{v}$};
    
     %Vector tfrac{1}{2}(u+V
    \draw[->, thick, color=red] (0,0) -- (X) node[midway, above ,pos=1.1] {\scriptsize $\tfrac{1}{2}\vec{v}+ \tfrac{1}{2}\vec{u}$};
 
    \end{tikzpicture}
    \end{center}










    
\end{enumerate}

    
    \item ¿Qué propiedad relaciona los incisos {\scriptsize$III)$ }y {\scriptsize$V)$}? \\
    Los resultados de los incisos {\scriptsize$II)$ }y {\scriptsize$VI)$} se relacionan por la propiedad distributiva y se obtiene el mismo vector resultado.
\end{enumerate}















    \item  Dados los vectores $\vec{v} = (4, -8, 1)$, $\vec{w} = \left(\tfrac{1}{10}, 0, -1\right)$ y $\vec{z} = (-4, 2, 0)$:

    \begin{enumerate}[a)]
        \item Calcular analíticamente cada una de las operaciones que se indican:
            \begin{enumerate}[i)]
            \begin{multicols}{2}
                \item $3(\vec{v} - 2\vec{w})$.
                \item $-2\left(\tfrac{1}{2}\vec{z} + \vec{v}\right) - 3\vec{w}$.                 
            \end{multicols}
            \end{enumerate}
        \item Determinar las componentes del vector $\vec{u}$ que satisfaga: $2\vec{v} - \vec{u} + 3\vec{z} = \vec{w} - 5\vec{u}$.
    \end{enumerate} 

\bigskip

\textbf{Solución:}



\begin{enumerate}[a)]

\item \textbf{Resolución del ejercicios combinados:} primero resolvemos los paréntesis, después la multiplicaciones y finalmente las sumas y restas.



\begin{enumerate}[i)]
\begin{minipage}{.4\textwidth}   

        \item   \,  

        
        \hspace{-3mm}\begin{minipage}{.36\textwidth}\begin{align*}
        &\  \ \ \ \  3(\vec{v} - 2\vec{w}) =
        \\
        &= 
        3 \left( \rule{0mm}{4.5mm} (4,-8,1) - 2  \left(\tfrac{1}{10}, 0, -1\right) \right)
        \\[3mm]
        &= 
        3 \left( \rule{0mm}{4.5mm} (4,-8,1) -   \left(\tfrac{1}{5}, 0, -2\right) \right) 
        \\[3mm] 
        &= 
        3
        \left(4 - \tfrac{1}{5},\, -8 - 0, \, 1 - (-2)\right)
        \\[3mm]   
        &= 
         \quad  3    \left(\tfrac{19}{5}, -8, 3\right) 
         \\[3mm]   
        &= 
         \quad    \left(\tfrac{57}{5}, -24, 9\right) 
        \end{align*}
        \end{minipage}


\end{minipage}\qquad\begin{minipage}{.5\textwidth}


        \item  \, 

        \begin{align*}
        \qquad & \ \, \ 
        -2\left(\tfrac{1}{2}\vec{z} + \vec{v}\right) - 3\vec{w} 
        \\  
        &= 
        -2\left( \rule{0mm}{4.5mm} \tfrac{1}{2} (-4, 2, 0)  + (4, -8, 1) \right) - 3 \left(\tfrac{1}{10}, 0, -1\right)
        \\[4mm]
        &= 
        -2\left( \rule{0mm}{4.5mm} (-2, 1, 0)  + (4, -8, 1) \right) -  \left(\tfrac{3}{10}, 0, -3\right)
        \\[4mm] 
        &= 
        -2\left( \rule{0mm}{4mm}-2+4 , 1 -8, 0 +1    \right) -  \left(\tfrac{3}{10}, 0, -3\right)
        \\[4mm]   
        &= 
         -2\left( \rule{0mm}{4mm}2 , -7,1   \right) -  \left(\tfrac{3}{10}, 0, -3\right)
         \\[4mm]   
        &= 
          \left( \rule{0mm}{4mm}-4 , \, 14, \, -2   \right) -  \left(\tfrac{3}{10}, 0, -3\right)
        \\[4mm]   
        &= 
        \left(-4 - \tfrac{3}{10}, 14 - 0, -2 - (-3)\right)
        \\[4mm]   
        &= 
        \left(-\tfrac{43}{10}, 14, 1\right).
        \end{align*}
    
\end{minipage}       
    \end{enumerate}

    \item Determinar las componentes del vector $\vec{u}$ que satisfaga: 
    
    \begin{align*}
     2\vec{v} - \vec{u} + 3\vec{z} \quad &= \quad \vec{w} - 5\vec{u}  \hspace{70mm}
    \intertext{Procedemos a despejar  $\vec{u}$ }%%%%%%%%%%%_____---------____-----___
    2\vec{v} + 3\vec{z} - \vec{w}\quad &= \quad -5\vec{u} + \vec{u}  
    \\
    2\vec{v} + 3\vec{z} - \vec{w}  \quad &= \quad -4\vec{u}   
    \\
    -\tfrac{1}{4}\left(2\vec{v} + 3\vec{z} - \vec{w} \right)    \quad &= \qquad   \vec{u}  
    \intertext{remplazando:}
    -\tfrac{1}{4}\left(2 (4,-8,1) + 3 (-4,2,0) - \left(\tfrac{1}{10}, 0,-1 \right) \rule{0mm}{5mm} \right)      \quad &= \qquad    \vec{u}  
    \\
    -\tfrac{1}{4}\left( (8,-16,2) +  (-12,6,0) - \left(\tfrac{1}{10}, 0,-1 \right) \rule{0mm}{5mm} \right)     \quad &= \qquad \vec{u}  
    \\
    -\tfrac{1}{4} \left( \rule{0mm}{4.5mm}8 - 12 - \tfrac{1}{10}, -16 + 6 + 0, 2 + 0 + 1\right)   \quad &= \qquad   \vec{u}  
    \\
    \tfrac{1}{-4} \quad\left(\rule{0mm}{4.5mm} -\tfrac{41}{10}, -10, 3\right) \hspace{20mm} \quad &= \qquad  \vec{u} 
     \\
     \left(\rule{0mm}{4.5mm} \tfrac{41}{40}, \tfrac{5}{2}, -\tfrac{3}{4}\right) \hspace{15mm}     \quad &= \qquad   \vec{u} 
    \end{align*}
    



\end{enumerate}










\item  Usar interpretación geométrica de las operaciones indicadas y obtener el vector resultante en cada caso:
\begin{enumerate}
    \item \begin{minipage}{.6\textwidth}
    \begin{tikzpicture}[scale=1,>=stealth]

    \coordinate (A) at (0,2);
    \coordinate (B) at (1,1);
    \coordinate (C) at (1.5,0);
    \coordinate (D) at (.7,-.7); 
    % Vector a
    \draw[->, thick] (0,0) -- node[left] {$\vec{a}$}(A) ;
    
    % Suma
    \node[red] at (0.6,1) {\Large$+$};
    
    % Vector b 
    \coordinate (X1) at (1,.5);
    \coordinate (BB) at ($(B)+(X1)$);
    
    \draw[->, thick ] (X1) --node[above left] {$\vec{b}$} (BB) ;
    
    % Suma
    \node[red] at (2.7,1) {\Large$+$};
    
    % Vector c 
    \coordinate (X2) at (3.1,.9);
    \coordinate (CC) at ($(C)+(X2)$);
    \draw[->, thick ] (X2) --node[above] {$\vec{c}$} (CC);
    
    % Suma
    \node[red] at (5.8,1) {\Large$+$};
    
    % Vector d
    \coordinate (X3) at (6.3,1.5);
    \coordinate (DD) at ($(D)+(X3)$);
    \draw[->, thick ] (X3) --node[above right] {$\vec{d}$} (DD) ;
\end{tikzpicture}
\end{minipage} \begin{minipage}{.3\textwidth}


    \item \begin{minipage}{.7\textwidth}\begin{tikzpicture}[scale=1,>=stealth]
    \coordinate (A) at (-2,2);
    \coordinate (B) at (1,1);
    % Vector a
    \draw[->, thick] (0,0) --node[left] {$\vec{a}$} (A) ;
    
    % Resta
    \node[red] at (0.5,1) {\Large$-$};
    
    % Vector b     
    \coordinate (X1) at (.9,.5);
    \coordinate (BB) at ($(B)+(X1)$);
    \draw[->, thick,  ] (X1) --node[above left] {$\vec{b}$} (BB) ;
\end{tikzpicture}
\end{minipage} 
\end{minipage} 





\medskip
    
    \item  \begin{minipage}{.7\textwidth}\begin{tikzpicture}[scale=1,>=stealth]


    \coordinate (A) at (0,2);
    \coordinate (B) at (1,1);
    \coordinate (C) at (1.5,0);
    \coordinate (D) at (-.7,.7); 
    
    % Vector a
    \draw[->, thick, shift={(0,-.25)}] (0,0) --node[left] {$\vec{a}$} (0,1.5);
    
    % Suma
    \node[red] at (0.5,.5) {\Large$+$};
    
    % Vector b    
    \coordinate (X1) at (1,0);
    \coordinate (BB) at ($(B)+(X1)$);
    \draw[->, thick ] (X1) --node[ left] {$\vec{b}$}  (BB);
    
    % Resta
    \node[red] at (2.5,0.5) {\Large$-$};
    
    % Vector c
    \coordinate (X2) at (3,0.35);
    \coordinate (CC) at ($(C)+(X2)$);
    \draw[->, thick ] (X2) -- node[above] {$\vec{c}$}(CC);
    
    % Suma
    \node[red] at (5.5,.5) {\Large $+$};
    
    % Vector d 
    \coordinate (X3) at (6.7,0);
    \coordinate (DD) at ($(D)+(X3)$);
    \draw[->, thick ] (X3) --node[right] {$\vec{d}$} (DD) ;
\end{tikzpicture}
\end{minipage} 
\end{enumerate}










\textbf{Solución:}


\begin{enumerate}
    \item \begin{minipage}{.36\textwidth}
    \begin{tikzpicture}[scale=1,>=stealth]

    \coordinate (A) at (0,2);
    \coordinate (B) at (1,1);
    \coordinate (C) at (1.5,0);
    \coordinate (D) at (.7,-.7); 
    % Vector a
    \draw[->, thick] (0,0) -- node[left] {$\vec{a}$}(A) ;
    
 
    % Vector b 
    \coordinate (X1) at (A);
    \coordinate (BB) at ($(B)+(X1)$);
    
    \draw[->, thick ] (X1) --node[above left] {$\vec{b}$} (BB) ;
     
    
    % Vector c 
    \coordinate (X2) at ($(A)+(B)$);
    \coordinate (CC) at ($(C)+(X2)$);
    \draw[->, thick ] (X2) --node[above] {$\vec{c}$} (CC);
     
    
    % Vector d
    \coordinate (X3) at ($(A)+(B)+(C)$);
    \coordinate (DD) at ($(D)+(X3)$);

    
    \draw[->, thick ] (X3) --node[above right] {$\vec{d}$} (DD) ;

    % Vector suma
    \coordinate (X3) at (0,0);
    \coordinate (OP) at ($(A)+(B)+(C)+(D)$);
    \draw[->, thick,teal ] (X3) --node[right] {$\vec{a}+\vec{b}+\vec{c}+\vec{d}$} (OP) ;
\end{tikzpicture}
\end{minipage} \begin{minipage}{.28\textwidth}


    \item \begin{minipage}{.7\textwidth}\begin{tikzpicture}[scale=1,>=stealth]
    \coordinate (A) at (-2,2);
    \coordinate (B) at (1,1);
    % Vector a
    \draw[->, thick] (0,0) --node[left] {$\vec{a}$} (A) ;

    
    % Vector b     
    \coordinate (X1) at (A);
    \coordinate (BB) at ($(X1)-(B)$);
    \draw[->, thick, red ] (X1) --node[left] {$-\vec{b}$} (BB) ;


    % Vector suma
    \coordinate (X3) at (0,0);
    \coordinate (OP) at ($(A)-(B)$);
    \draw[->, thick,teal ] (X3) --node[below left] {$\vec{a}-\vec{b}$} (OP) ;
\end{tikzpicture}
\end{minipage} 
\end{minipage}\begin{minipage}{.3\textwidth}
    \item  
    \begin{minipage}{.7\textwidth}\begin{tikzpicture}[scale=1,>=stealth]


    \coordinate (A) at (0,1.5);
    \coordinate (B) at (1,1);
    \coordinate (C) at (1.5,0);
    \coordinate (D) at (-.7,.7); 
    
    % Vector a
    \draw[->, thick, shift={(0,-.25)}] (0,0) --node[pos=.7,right] {$\vec{a}$} (A);
     
    
    % Vector b    
    \coordinate (X1) at (A);
    \coordinate (BB) at ($(B)+(X1)$);
    \draw[->, thick ] (X1) --node[ left] {$\vec{b}$}  (BB);
    
    
    
    % Vector c
    \coordinate (X2) at ($(A)+(B)$);
    \coordinate (CC) at ($(X2)-(C)$);
    \draw[->, thick,red ] (X2) -- node[above] {$-\vec{c}$}(CC);
     
    
    % Vector d 
    \coordinate (X3) at ($(A)+(B)-(C)$);
    \coordinate (DD) at ($(D)+(X3)$);
    \draw[->, thick ] (X3) --node[right] {$\vec{d}$} (DD) ;

    % Vector suma
    \coordinate (X3) at (0,0);
    \coordinate (OP) at ($(A)+(B)-(C)+(D)$);
    \draw[->, thick,teal ] (X3) --node[left] {$\vec{a}+\vec{b}-\vec{c}+\vec{d}$} (OP) ;
\end{tikzpicture}
\end{minipage} 
\end{minipage}
\end{enumerate} 


\bigskip 



\item   Dados los puntos el espacio $A = (1, 3, 1)$, $B = (1, 2, 1)$ y $C = (0, 0, 2)$: 

    \begin{enumerate}[ a)] 
        \item Determinar las componentes de:
            \begin{enumerate}[i)]
                \item El vector $\overrightarrow{AB}$.
                \item El vector $\vec{v}$ que tiene a los puntos $B$ como origen y $C$ como extremo final.
            \end{enumerate}
        \item El punto $D$ tal que $\overrightarrow{AD}$ resulte equivalente al vector $\overrightarrow{BC}$.
        \item  Un punto $M$ tal que el vector $\overrightarrow{CM}$ resulte paralelo al vector $\overrightarrow{AB}$. 
        \item Considerando resultados anteriores y la definición de módulo (o norma) de un vector, calcular:
        \begin{enumerate}[i)]
        \begin{minipage}{.18\textwidth}         \item $\|\overrightarrow{AB}\|$
        \end{minipage}\quad\begin{minipage}{.18\textwidth}
            \item $\|\vec{v}\|$
        \end{minipage}\quad\begin{minipage}{.18\textwidth}
            \item $\|-\vec{v}\|$
        \end{minipage}\quad\begin{minipage}{.3\textwidth}
            \item $\|\sqrt{2} \cdot \overrightarrow{m}\|$ \quad siendo $\overrightarrow{m}= \overrightarrow{AC}$.
        \end{minipage}\ 
        \end{enumerate}
    \end{enumerate}








\textbf{Solución:}
\begin{enumerate}[a)]
\item  Las componentes del vector $\vec{v}= \overrightarrow{PQ}$ se obtiene como la diferencia de las coordenadas del punto final $Q$ y el punto inicial $P$, es decir:
\[
\vec{v} = \overrightarrow{OQ} -\overrightarrow{OP}.
\] 
    \begin{enumerate}[ i)] 

    \item Para $A= (1,3,1)$ y $B =(1,-2,1)$:
    \[
    \overrightarrow{AB}= \overrightarrow{OB} - \overrightarrow{OA} = (1,-2,1) -(1,3,1) = (1 -1, -2 - 3, 1-1 ) = (0, -5, 0).
    \]
    Por lo tanto, las componentes del vector son:\qquad $\vec{v} = (0, -5,0).$
    

    \item Para $B= (1,-2,1)$ y $C= (1,0,-2)$ El vector $\vec{v}$ que tiene a $B$ como origen y $C$ como extremo final:
    \[
    \vec{v} = \overrightarrow{BC} = \overrightarrow{OC} -\overrightarrow{OB}
    = 
    \left( \rule{0mm}{4mm}  1,0,-2 \right) - \left( \rule{0mm}{4mm} 1,-2,1 \right) 
    = 
    \left( \rule{0mm}{4mm} 1 - 1 , 0 - (-2), -2-1 \right)
     = (0, 2, -3).
    \]
    Por lo tanto, las componentes del vector son: \qquad $\vec{v} = (0, 2, -3).$

    \end{enumerate}

\medskip 


    \item  El punto $D$ tal que $\overrightarrow{AD}$ resulte equivalente al vector $ \overrightarrow{BC}$. 



    \bigskip    
    
    
    Recordemos que $\overrightarrow{AD}$ resulte equivalente al vector $ \overrightarrow{BC}$  implica que las componentes de ambos vectores son iguales, es decir $\overrightarrow{AD} = \overrightarrow{BC}$. Esto se expresa como:
    \begin{align*}
        \overrightarrow{AD} &= \overrightarrow{BC}\\
        \overrightarrow{OD} -\overrightarrow{OA} &= \overrightarrow{OC}-\overrightarrow{OB}        
        \intertext{Remplazando:}
        D - (1, 3, 1)  &=   (1,0,-2) - (1, -2, 1) 
        \intertext{Despejando $D$, tenemos:} 
        D &=   (1,0,-2) - (1, -2, 1) + (1, 3, 1)  
        \\
        &= \left( \rule{0mm}{4mm} 1-1+1,\, 0-(-2) + 3, \, -2-1+1 \right) \hspace{-15mm}
        \\
        &= (1, \, 5 ,\, -2 )
    \end{align*}

   %-----------------------------------------------__-----------------------------------------------------------------------------------------------------------------------------------------
    \item Un punto $M$ tal que $\overrightarrow{CM} \paralela \overrightarrow{AB}$, entonces  debe cumplirse que:
     
   

    
    \begin{align*}
        \overrightarrow{CM} &= k \cdot \overrightarrow{AB},\marka{b}{(-.3,-.2)} & k \in \mathbb{R}-\{0\}. 
        \intertext{Entonces:}
        \overrightarrow{CM} &= k \cdot (0, -1, 0) \marka{c}{(-.7,.47)}  
        \\[2mm]
        &= (0, -5k, 0)
        \intertext{Lo que implica que:}
        M-C &= (0,-5k,0) 
        \\[2mm]
        M &= (0,-5k,0) +C
        \\[2mm]
        &= (0,-5k,0) + (1,0,-2) \hspace{-9mm}
        \\[2mm]
        &= (1,-5k,-2)  & \text{ con   } k\in \R-\{0\}
    \end{align*}

\tikz[overlay,remember picture] { 
\coordinate (M) at ($1/2*(b)+1/2*(c)$);
\coordinate (a) at ($(M)+(1.8,0)$);
\draw[thick,magenta!30! , line width=1pt] (a)|-(M)--(b);
\draw[thick,magenta!30!,-> , line width=1pt] (M)--(c);
\node[fill=blue!30!white!30!,
rectangle,rounded corners=1mm,
%xshift=-7mm,
yshift=4mm
] at (a) [
text width=34mm, 
below right
]{\scriptsize 
 Calculamos   $\overrightarrow{AB}$:
    \begin{align*}
    \overrightarrow{AB} &= B - A 
    \\
    &= (1, -2, 1) - (1, 3, 1) 
    \\
    &= (0, -5, 0).
    \end{align*}
    };
}%%%%%%%%%%%%%%%%%%%%------------------------------$$$$------------------------------$$$$------------------------------$$$$------------------------------$$$$------------------------------$$$$------------------------------$$$$------------------------------$$$$------------------------------$$$$------------------------------$$$$------------------------------$$$$
Esto implica que el punto $M$ no es único, ya que el valor de $k$ puede variar. % El conjunto de todos los puntos $M$ que cumplen la condición forma una recta paralela a $\overrightarrow{AB}$ que pasa por $C$.  

\item Considerando resultados anteriores y la definición de módulo (o norma) de un vector, calcular:
        \begin{enumerate}[i)]
            \item Como $\overrightarrow{AB}$ es igual a $(0, -5, 0)$ tenemos:
            $$\|\overrightarrow{AB}\|=  \| (0, -5, 0)\|= \sqrt{0^2 +(-5)^2+ 0^2}= \sqrt{25\,} = 5.$$ 
            
            \item Como $\vec{v}= \overrightarrow{BC}$ es igual a $(0, 2,-3)$ tenemos:
            $$\|\vec{v}\|= \|(0, 2,-3)\| = \sqrt{0^2 + 2^2+ (-3)^2}= \sqrt{ 13\,} \approx 3,605$$
            
            \item Por propiedad tenemos que $\|-v\|= \|v\|$ con lo cual:
            $$\|-\vec{v}\| = \sqrt{ 13\,} \approx 3,605 $$
            
            \item Tenemos que: 
            \begin{itemize}
                \item $\overrightarrow{m} = \overrightarrow{AC} = C-A = (1,0,-2) -(1,3,1) = (1-1,\, 0-3,\, -2-1) = (0,\, -3,\, -3)$,
            \end{itemize}
            luego tenemos:
            $$\|\sqrt{2} \cdot \overrightarrow{m}\| =   \|(0,\, -\sqrt{2}\cdot 3,\, -\sqrt{2}\cdot 3 )\| = \cdot \sqrt{0^2 +(-\sqrt{2}\cdot3)^2+(-\sqrt{2}\cdot3)^2\, } = \sqrt{ 18+18} = \sqrt{36} = 6$$  
        \end{enumerate}
\end{enumerate}


















\item  Para los siguientes puntos $P = (1, -3)$, $Q = (2, 2)$ y $R = (6, -2) \in \mathbb{R}^2$:

      
\begin{enumerate}[a)]
    \item Calcular la distancia entre los puntos:
    \begin{enumerate}[i)]
    \begin{multicols}{3}
        \item $P$ y $Q$
        \item $Q$ y $R$
        \item $R$ y $P$
    \end{multicols}
    \end{enumerate}
    
    \item Si los tres puntos forman un triángulo, calcular su perímetro y clasificarlo según sus lados.
\end{enumerate}

\begin{flushright}
\small (Recordar: un triángulo se clasifica según sus lados en escaleno, isósceles o equilátero.)
\end{flushright}
 

{\bf Solución:}


\begin{enumerate}[a)]
    \item \textbf{Calcular la distancia entre los puntos:}
    
    Para $A=(x_1,y_1)$ y $B=(x_2,y_2)$ y utilizamos la fórmula de distancia entre dos puntos en el plano tenemos:
    \[
    d(A, B) = \| \overrightarrow{AB}\| = \sqrt{(x_2 - x_1)^2 + (y_2 - y_1)^2}
    \]

    \begin{enumerate}[i)]
        \item Entre $P$ y $Q$:
        \[
        d(P, Q)  = \| \overrightarrow{PQ}\| = \sqrt{(2 - 1)^2 + (2 - (-3))^2} = \sqrt{1^2 + 5^2} = \sqrt{1 + 25} = \sqrt{26}
        \]

        \item Entre $Q$ y $R$:
        \[
        d(Q, R)  = \| \overrightarrow{QR}\| = \sqrt{(6 - 2)^2 + (-2 - 2)^2} = \sqrt{4^2 + (-4)^2} = \sqrt{16 + 16} = \sqrt{32} = 4\sqrt{2}
        \]

        \item Entre $R$ y $P$:
        \[
        d(R, P)  = \| \overrightarrow{RP}\| = \sqrt{(1-6)^2 + (-3- (-2))^2} = \sqrt{5^2 + 1^2} = \sqrt{25 + 1} = \sqrt{26}
        \]
    \end{enumerate}

    \item  \begin{itemize}
        \item El perímetro del triángulo es la suma de las longitudes de sus lados:
        \[
        \text{Perímetro} = d(P, Q) + d(Q, R) + d(R, P) = \sqrt{26} + 4\sqrt{2} + \sqrt{26} = 2\sqrt{26} + 4\sqrt{2}
        \]

        \item Clasificación según sus lados:

        Observamos que:
        \[
        d(P, Q) = \sqrt{26}, \quad d(Q, R) = 4\sqrt{2}, \quad d(R, P) = \sqrt{26}
        \]

        Como dos lados son iguales a $\sqrt{26}$, el triángulo es \textbf{isósceles}.
    \end{itemize}
\end{enumerate}





\item  Determinar todos los valores de $k \in \mathbb{R}$ para los cuales resulta:

    \begin{itemize}
    \begin{multicols}{2}
        \item[a)] $\|\vec{v}\| = 2$, siendo $\vec{v} = (1, k, 0)$.
        \item[b)] $\|\vec{v}\| = 1$, siendo $\vec{v} = k(2, 2, 1)$.
    \end{multicols}
    \end{itemize}





\textbf{Solución:}


Recordar que la norma de $\vec{v}$ se calcula como: \qquad $\|\vec{v}\| = \sqrt{x^2 + y^2 + z^2},$

    

\begin{itemize}
    \item[a)] Determinemos los valores de $k$ para los cuales $\|\vec{v}\| = 2$, siendo $\vec{v} = (1, k, 0)$  En este caso: 
    $$
    \|\vec{v}\| = \sqrt{1^2 + k^2 + 0^2} = \sqrt{1 + k^2}. $$

    \medskip
     
    Igualamos la norma a 2 y resolvemos la ecuación:

    \begin{minipage}{.45\textwidth} 
    \begin{align*}
        \sqrt{1 + k^2}& = 2
        \\[3mm]
       \left( \sqrt{1 + k^2}  \rule{0mm}{4.5mm}\right)^2 &= 2^2
        \\[3mm]
        1 + k^2 &= 4 
        \\[3mm] 
        k^2 &= 3 
        \\[3mm]
        \sqrt{k^2} &= \sqrt{ 3 \, }
        \\[3mm]%%%%%% %%%%%%%% %%%%%%% %%%%%%% %%%%%% %%%%%%%% %%%%%%% %%%%%%% %%%%%% %%%%%%%% %%%%%%% %%%%%%% %%%%%% %%%%%%%% %%%%%%% %%%%%%% %%%%%% %%%%%%%% %%%%%%% %%%%%%% %%%%%% %%%%%%%% %%%%%%% %%%%%%% %%%%%% %%%%%%%% %%%%%%% %%%%%%% 
     k =  - \sqrt{3} \quad&o    \quad k=\sqrt{3} 
    \end{align*}
     \end{minipage}
      


    Por lo tanto, los valores de $k$ son: \qquad $    k = \sqrt{3} \quad \text{y} \quad k = -\sqrt{3}.$ 


\bigskip



    \item[b)] Determinemos los valores de $k$ para los cuales $\|\vec{v}\| = 1$, siendo $\vec{v} = k(2, 2, 1)$. \\
    Primero, escribimos $\vec{v}$ en función de $k$: \quad $\vec{v} = (2k, 2k, k).$
   
    La norma de $\vec{v}$ es:
    \[
    \|\vec{v}\| = \sqrt{(2k)^2 + (2k)^2 + k^2} = \sqrt{4k^2 + 4k^2 + k^2} = \sqrt{9k^2}.
    \]
    Igualamos la norma a 1 y despejamos:

\begin{minipage}{.45\textwidth}
      \begin{align*}
    \sqrt{9k^2} &= 1
    \\[3mm]
     \sqrt{9 } \cdot   \sqrt{k^2} &= 1  
    \\[3mm]   
    3 |k| &= 1  
    \\[3mm]
    |k| &= \tfrac{1}{3}
    \\[3mm]%%%%%% %%%%%%%% %%%%%%% %%%%%%% %%%%%% %%%%%%%% %%%%%%% %%%%%%% %%%%%% %%%%%%%% %%%%%%% %%%%%%% %%%%%% %%%%%%%% %%%%%%% %%%%%%% %%%%%% %%%%%%%% %%%%%%% %%%%%%% %%%%%% %%%%%%%% %%%%%%% %%%%%%% %%%%%% %%%%%%%% %%%%%%% %%%%%%% 
     k = - \tfrac{1}{3}\ \quad &o \qquad  \ k=\tfrac{1}{3} 
    \end{align*}
     

\end{minipage}



    Por lo tanto, los valores de $k$ son: \quad $k = \tfrac{1}{3} \quad \text{y} \quad k = -\tfrac{1}{3}$
    
\end{itemize}

\medskip 












    \item  Sean los vectores $\vec{u} = (3, 0, 1)$ y $\vec{p} = (1, -1, 0)$, hallar:

\begin{enumerate}[a)]
    \item Un vector paralelo a $\vec{p}$ que tenga norma cinco. ¿Cuántos vectores existen con estas condiciones?
    \item El vector opuesto a $\tfrac{1}{2}\vec{u}$.
    \item El versor asociado a $\vec{p}$.
    \item  Un vector unitario y de sentido contrario a $(\vec{u} + \vec{p})$. 
\end{enumerate}

\textbf{Resolución:}

\begin{enumerate}
    \item Buscamos un vector $\vec{w}$ que satisfaga las dos condiciones mencionadas. Realizamos el siguiente análisis:

\bigskip 


\begin{itemize}
\hspace{-12mm}
\begin{minipage}{.45\textwidth}
    \item Como $\vec{w} \paralela \vec{p} \ \Longrightarrow \ 
    \vec{w} =\indicador{A}{(-.1,-.15)}{B}{(-.1,-.45)} \marka{C}{(1,-.9)}  k (1,-1,0). 
    $

 

    \bigskip

    \item $\|\vec{w}\| = 5$ \marka{D}{(.1,.05)} 
\end{minipage}
\end{itemize}

\tikz[overlay,remember picture] {
\draw [overlay,->,very thick,red,opacity=.5]
(A) -| (B)to[out=0,in=180]% node [pos=0.62,text width=35mm,above,sloped] {\scriptsize reemplazando \\  $\vec{w} = k (1,-1,0)$ en $\|\vec{w}\|=5$ } 
(C);
\draw [overlay,very thick,red,opacity=.5]
(D)-|(B);
\node  at (C) [text width=40mm, below right ,yshift= 9mm]{
%%%%%%%%%%%%%%%%
\[
\begin{array}{rl}   
\|k (1,-1,0)\|
&=%%%%%%%%%%%%%%%%%%%%%%%%%5
5 \marka{a}{(.1,.1)}
\\[4mm] %%%%%%%%%%%%%%%%%%%%%%%%%%%%%%%%%%%%%%%%%%%%%%%%%%%%%%%%%%%%%%%%%%%%%%%%%%%%%%%%%%%%%%%%%%%%%%%%%%%%%%%%%%%%%%%%%%%%%%%%%%%%%%%%%%%%%%%%%%%%%%%%%%%%%%%%%%%%%%%%%%%%%%%%%%%%%%%%%%%%%%%%%%%%%%%%%%%%%%%%%%%%%%%%%%%%%%%%%%%%%%%%%%%%%%%%%%%%%%%%%%%%%%%%%%%%%%%%%%%%%%%% 
\hspace{-8mm}
\marka{b2}{(-.2,.05)} 
|k| \cdot \| (1,-1,0)\|
&=%%%%%%%%%%%%%%%%%%%%%5
5  \marka{b1}{(.1,.2)}   
\\[4mm] %%%%%%%%%%%%%%%%%%%%%%%%%%%%%%%%%%%%%%%%%%%%%%%%%%%%%%%%%%%%%%%%%%%%%%%%%%%%%%%%%%%%%%%%%%%%%%%%%%%%%%%%%%%%%%%%%%%%%%%%%%%%%%%%%%%%%%%%%%%%%%%%%%%%%%%%%%%%%%%%%%%%%%%%%%%%%%%%%%%%%%%%%%%%%%%%%%%%%%%%%%%%%%%%%%%%%%%%%%%%%%%%%%%%%%%%%%%%%%%%%%%%%%%%%%%%%%%%%%%%%%%%
\marka{c}{(-.2,.15)}  
| k |\cdot \sqrt{ 2 } 
&=%%%%%%%%%%%%%%%%%%%%%5
5  
\\[4mm] %%%%%%%%%%%%%%%%%%%%%%%%%%%%%%%%%%%%%%%%%%%%%%%%%%%%%%%%%%%%%%%%%%%%%%%%%%%%%%%%%%%%%%%%%%%%%%%%%%%%%%%%%%%%%%%%%%%%%%%%%%%%%%%%%%%%%%%%%%%%%%%%%%%%%%%%%%%%%%%%%%%%%%%%%%%%%%%%%%%%%%%%%%%%%%%%%%%%%%%%%%%%%%%%%%%%%%%%%%%%%%%%%%%%%%%%%%%%%%%%%%%%%%%%%%%%%%%%%%%%%%%%
 |k| 
&=%%%%%%%%%%%%%%%%%%%%%5
\dfrac{5}{\sqrt{2} }
\\[4mm]%%%%%% %%%%%%%% %%%%%%% %%%%%%% %%%%%% %%%%%%%% %%%%%%% %%%%%%% %%%%%% %%%%%%%% %%%%%%% %%%%%%% %%%%%% %%%%%%%% %%%%%%% %%%%%%% %%%%%% %%%%%%%% %%%%%%% %%%%%%% %%%%%% %%%%%%%% %%%%%%% %%%%%%% %%%%%% %%%%%%%% %%%%%%% %%%%%%% 
     k  = \tfrac{5}{\sqrt{ 2} } \quad & o \quad k= - \tfrac{5}{\sqrt{ 2 } } 
\end{array}
\]
%%%%%%%%%%%55
};
}

 
\tikz[overlay,remember picture] {
  \draw[thick,red,->, line width=1pt] (a) to[out=0,in=0] node [pos=0.45,text width=25mm,fill=blue!5!green!5!, rectangle,rounded corners=3mm ,right] {\footnotesize Propiedad: \\ $ \|\alpha \cdot \Vec{u}\| = |\alpha| \, \|\vec{u}\|$  } (b1) ;
  \draw[thick,red,->, line width=1pt] (b2) to[out=180,in=180] node [pos=0.7, left] {\marka{CA2}{(0,-3mm)} \quad \marka{CA1}{(0,0)}   } (c) ;
} 

\begin{tikzpicture}[overlay,remember picture]
\node[fill=red!30!blue!20!, rectangle,rounded corners=3mm] at (CA2) [text width=58mm, below left ,yshift= 9mm]{ \scriptsize $\begin{array}{rl}
 \text{C.A: } \  \| (1,-1,01)\|
&=%%%%%%%%%%%%%%%%%%%%%5
\sqrt{ 1^2 + (-1)^2 + 0^2 } 
\\%%%%%%%%%%%%%%%%%%%%%5%%%%%%%%%%%%%%%%%%%%%5%%%%%%%%%%%%%%%%%%%%%&=%%%%%%%%%%%%%%%%%%%%%5
&=
\sqrt{ 1 + 1 +0} 
\\%
&=%%%%%%%%%%%%%%%%%%%%%5
\sqrt{ 2 }  
\end{array}$ 
\[     
\therefore \| (1,-1,0) \|=\sqrt{2\, }   
\]
}   ;
%\draw [overlay,->,very thick,red,opacity=.5](CA2)  to[out=60, in= 180] (CA1);
\end{tikzpicture}



\vspace{27mm}






    Por lo tanto, $k$ puede tomar dos valores y los vectores resultantes son:
    \[
    \vec{w}_1 = \tfrac{5}{\sqrt{2}} \cdot (1, -1, 0) = \left(\tfrac{5}{\sqrt{2}} , -\tfrac{5}{\sqrt{2}} , 0\right),
    \qquad \qquad 
    \vec{w}_2 = -\tfrac{5}{\sqrt{2}}  \cdot (1, -1, 0) = \left(-\tfrac{5}{\sqrt{2}} , \tfrac{5}{\sqrt{2}} , 0\right).
    \]
    Por lo tanto, existen exactamente dos vectores con estas condiciones.

    \item El vector opuesto a $\tfrac{1}{2}\vec{u}$: \\
    El vector opuesto a $\vec{v}$ se obtiene cambiando el signo de cada componente, es decir: $-\vec{v}.$
     
    Calculamos :
    \[
    -\tfrac{1}{2}\vec{u} = -\tfrac{1}{2} \cdot (3, 0, 1) = \left( -\tfrac{3}{2}, 0, -\tfrac{1}{2}\right).
    \]


    \item \textbf{El versor asociado a $\vec{p}$.}

    El versor asociado a un vector $\vec{p}$ es el vector unitario con la misma dirección y sentido que $\vec{p}$, y se obtiene dividiendo el vector por su módulo: \quad $\check{p} = \frac{\vec{p}}{\|\vec{p}\|} $

    Primero, calculamos el módulo de $\vec{p}$:

    \[
    \|\vec{p}\| = \sqrt{1^2 + (-1)^2 + 0^2} = \sqrt{1 + 1 + 0} = \sqrt{2}
    \]

    Entonces, el versor asociado a $\vec{p}$ es:

    \[
    \check{p} = \frac{1}{\sqrt{2}}(1, -1, 0) = \left( \tfrac{1}{\sqrt{2}}, \, \tfrac{-1}{\sqrt{2}}, \, 0 \right)
    \]


    \item Un versor paralelo a $(\vec{u} + \vec{p})$ y con sentido contrario: \\
    Primero, calculamos $\vec{u} + \vec{p}$:
    \[
    \vec{u} + \vec{p} = (3, 0, 1) + (1, -1, 0) = (4, -1, 1).
    \]
    Calculamos la norma de $\vec{u} + \vec{p}$:
    \[
    \|\vec{u} + \vec{p}\| = \sqrt{4^2 + (-1)^2 + 1^2} = \sqrt{16 + 1 + 1} = \sqrt{18} = 3\sqrt{2}.
    \]
    El versor paralelo a $\vec{u} + \vec{p}$ se obtiene dividiendo el vector por su norma y cambiando el signo para tener sentido contrario:
    \[
    \text{Versor} = -\tfrac{\vec{u} + \vec{p}}{\|\vec{u} + \vec{p}\|} = -\dfrac{ \,  (4, -1, 1)\, }{3\sqrt{2}} = \left(-\tfrac{4}{3\sqrt{2}}, \tfrac{1}{3\sqrt{2}}, -\tfrac{1}{3\sqrt{2}}\right).
    \]
\end{enumerate}

\textbf{Resultados:}
\begin{enumerate}
    \item Los vectores paralelos a $\vec{p}$ con norma cinco son:\quad $\vec{v}_1 = \left(\tfrac{5\sqrt{2}}{2}, -\tfrac{5\sqrt{2}}{2}, 0\right), \quad \vec{v}_2 = \left(-\tfrac{5\sqrt{2}}{2}, \tfrac{5\sqrt{2}}{2}, 0\right).$
    


    
    \item El vector opuesto a $\tfrac{1}{2}\vec{u}$ es: \qquad $ -\tfrac{1}{2}\vec{u} = \left(-\tfrac{3}{2}, 0, -\tfrac{1}{2}\right).
    $ 

    \item \textbf{El versor asociado a $\vec{p}$} es: \quad $
    \check{p} = \frac{1}{\sqrt{2}}(1, -1, 0) = \left( \tfrac{1}{\sqrt{2}}, \, \tfrac{-1}{\sqrt{2}}, \, 0 \right)
    $



    
    \item El versor paralelo a $(\vec{u} + \vec{p})$ y con sentido contrario es:\qquad  Versor$ = \left(-\tfrac{4}{3\sqrt{2}}, \tfrac{1}{3\sqrt{2}}, -\tfrac{1}{3\sqrt{2}}\right).$
     
\end{enumerate}
    
    
    








\item  Dados los vectores $\vec{u} = \left(-2, \tfrac{9}{2} \right)$ y $\vec{a} = (0, 10, -5)$ del plano y del espacio respectivamente.

\begin{enumerate}[a)]
    \item Expresarlos según la dirección de los versores canónicos correspondientes y representar en un gráfico.

    \item En cada caso, justificar correctamente y determinar si $\vec{u}$ puede descomponerse en las direcciones de los vectores: \quad\qquad \begin{minipage}{.7\textwidth}\begin{multicols}{2}
    \begin{enumerate}[i)]
        \item $\vec{v} = \left(1, \tfrac{3}{2} \right)$ y $\vec{w} = (8, -3)$.
        \item $\vec{v} = (2, -1)$ y $\vec{w} = \left(0, -\tfrac{11}{6} \right)$.
    \end{enumerate}    
    \end{multicols}    
    \end{minipage}   

    \item ¿El vector $\vec{a}$ es combinación lineal de los vectores $\vec{b} = (2, 4, -1)$ y $\vec{c} = (1, -3, 2)$?
\end{enumerate}








\textbf{Solución:}  

\begin{enumerate}[a)]
    \item Recordemos que en $\mathbb{R}^2$, los versores canónicos son $\iv = (1, 0)$ y $\jv = (0, 1)$, mientras que en $\mathbb{R}^3$ tenemos: $\iv = (1, 0,0)$,  $\jv = (0, 1,0)$ y $\kv = (0, 0, 1)$. Un vector $\vec{w} = (x, y, z)$ puede expresarse como:
\[
\vec{w} = x\iv + y\jv + z\kv.
\]

\begin{enumerate}[a)]
\begin{multicols}{2}
    \item Descomposición de $\vec{u} = \left(-2, \frac{9}{2}\right)$: \\
    En $\mathbb{R}^2$, tenemos:
    \[
    \vec{u} = -2\iv + \tfrac{9}{2}\jv.
    \]
   o alternativamente:
    \[
    \vec{u} = -2(1, 0) + \tfrac{9}{2}(0, 1).
    \]

\columnbreak % Forzar cambio de columna si es necesario
    \item Descomposición de $\vec{v} = (0, 10 , -5 )$: \\
    En $\mathbb{R}^3$, tenemos:
    \[
    \vec{v} = 0 \iv + 10 \jv - 5\kv.
    \]
    o alternativamente:
    \[
    \vec{v} = 0\,(1, 0, 0) + 10 \, (0, 1, 0) - 4 (0, 0, 1).
    \]

\end{multicols}
\end{enumerate}

\textbf{Resultados:}
\begin{enumerate} 
\begin{multicols}{2}
    \item  $\vec{u} = -2\iv + \frac{9}{2}\jv$.
    \item $\vec{v} = 0\iv + 10\jv - 5\kv$.
\end{multicols}
\end{enumerate}



    \item Recordemos que un vector $\vec{u}$ puede descomponerse en las direcciones de $\vec{v}$ y $\vec{w}$ si existen escalares $\alpha$ y $\beta$ tales que:

$$
\vec{u} = \alpha \vec{v} + \beta \vec{w}
$$

    \begin{enumerate}[i)]
        \item  $\vec{u} = \left(-2, \frac{9}{2}\right), \quad \vec{v} = \left(1, \frac{3}{2}\right), \quad \vec{w} = (8, -3)$

        Planteamos la definición de descomposición de $\vec{u}$ en la dirección de  $\vec{v}$ y $\vec{w}$ tenemos:
\begin{align*}
    \vec{u} &= \alpha \vec{v} + \beta \vec{w}
    \\
    \left(-2, \frac{9}{2}\right) &= \alpha \left(1, \frac{3}{2}\right) + \beta  (8, -3)
    \\
    \left(-2, \frac{9}{2}\right) &= \left( \alpha + 8\beta ,\, \tfrac{3}{2} \alpha - 3\beta \right)
\end{align*}
Al aplicar la igualdad entre vectores obtenemos el siguiente sistema de ecuación:

$$
\begin{cases}
\alpha + 8\beta = -2 \\
\frac{3}{2} \alpha - 3\beta = \frac{9}{2}
\end{cases}
$$

Multiplicamos la segunda ecuación por 2 para eliminar denominadores y resolvemos el sistema de ecuación:

\medskip
$
\left\{ \begin{array}{rl}
    \alpha + 8\beta &= -2   \indicador{x}{(.1,.15)}{y}{(1.3,.15)}
    \\[4mm]
    3\alpha - 6\beta &= 9   \marka{a}{(.1,.15)}
\end{array}\right.
$




\tikz[overlay,remember picture] {
\draw [overlay,->,very thick,red,opacity=.5](x)--(y);
\node  at (y) [ right ]{
$\alpha
= \marka{XX}{(-.25,-.15)}%%%%%
-2 - 8 \beta
$ \indicador{x}{(.1,.15)}{y}{(1.5,-.25)}\marka{z}{(2.3,-.25)}
};
%%%%%%%%%%%%%%%%%%%%%%%%
\draw [overlay,->,very thick,red,opacity=.5](x)-|(y)--(z);
\draw [overlay,very thick,red,opacity=.5](a)-|(y);
\node  at (z) [text width=40mm, below right , yshift= 8mm]{ 
\[
\begin{array}{rl}
3 (-2-8\beta)-6 \beta  
&=%%%%%%%%%%%%%%%%%%%%%%
9
 \\[2mm]%%%%%%%%%%%%%%%%%%%%%%%%%%%%%%%%%%%%%%%%%%%%%%%%%%%%%%%%%%%%%%%%%%
-6 -24 \beta  -6 \beta 
  &=%%%%%%%%%%%%%%%%%%%%%%
 9
  \\[2mm]%%%%%%%%%%%%%%%%%%%%%%%%%%%%%%%%%%%%%%%%%%%%%%%%%%%%%%%%%%%%%%%%%%
-30\beta
  &=%%%%%%%%%%%%%%%%%%%%%%
  15
  \\[2mm]%%%%%%%%%%%%%%%%%%%%%%%%%%%%%%%%%%%%%%%%%%%%%%%%%%%%%%%%%%%%%%%%% 
 \beta
  &=%%%%%%%%%%%%%%%%%%%%%%
  15: (-30) 
  \\[3mm]%%%%%%%%%%%%%%%%%%%%%%%%%%%%%%%%%%%%%%%%%%%%%%%%%%%%%%%%%%%%%%%%%% 
  \marka{x1}{(-.1,-.2)}
  \beta
  &= \indicador{YY}{(-.25,-.2)}{ZZ}{(-5,-.4)}\marka{TT}{(-3.5,-1.3)}%%%%%%%%%%%%%%%%%%%%%%
  -\tfrac{1}{2} \marka{x2}{(.1,.5)}  
\end{array}
\]
};
%%%%%%%%%%%%%%%%%%%%%%%%%%%%%%%%%%%%%%%%%%%%%%%%%%%%%%%%%%%%%%%%%%%%%%%%%%%%%%%%%%%%%%%%%%%%%%%%%%%%%%%%%%%%%%%%%%%%%%%%%%%%%%%%%%%%%%%%%%%%%%%%%%%%%%%%%%%%%%%%%%%%%%%%%%%%%%%%%%%%
\draw[color=black!50!,rounded corners=1mm] (x1) rectangle (x2);
%%%%%%%%%%%%%%%%%%%%%%%%%%%%%%%%%%%%%%%%%%%%%%%%%%%%%%%%%%%%%%%%%%%%%%%%%%%%%%%%%%%%%%%%%%%%%%%%%%%%%%%%%%%%%%%%%%%%%%%%%%%%%%%%%%%%%%%%%%%%%%%%%%%%%%%%%%%%%%%%%%%%%%%%%%%%%%%%%%%%
\draw [overlay,->,very thick,blue,opacity=.5](YY)|-(ZZ)|-(TT);
\draw [overlay,very thick,blue,opacity=.5](XX)|-(ZZ);
\node  at (TT) [text width=33mm, below right , yshift= 9mm ]{
\[
\begin{array}{rl}
 \alpha
 &=%%%%%%%%%%%%%%%%%%%%%%
  -2 -8 \cdot (- \tfrac{1}{2})
 \\[4mm]%%%%%%%%%%%%%%%%%%%%%%%%%%%%%%%%%%%%%%%%%%%%%%%%%%%%%%%%%%%%%%%%%%
 \alpha
 &=%%%%%%%%%%%%%%%%%%%%%%
  -2 +4
 \\[4mm]%%%%%%%%%%%%%%%%%%%%%%%%%%%%%%%%%%%%%%%%%%%%%%%%%%%%%%%%%%%%%%%%%%
  \marka{a1}{(-.1,-.2)} 
   \alpha
  &=\indicador{YY}{(-.15,-.15)}{ZZ}{(-3,-.4)}\marka{TT}{(-2.5,-.9)}%%%%%%%%%%%%%%%%%%%%%%
  2 \marka{a2}{(.1,.5)}
\end{array}
\]
};
\draw[color=black!50!,rounded corners=1mm] (a1) rectangle (a2);
%%%%%%%%%%%%%%%%%%%%%%%%%%%%%%%%%%%%%%%%%%%%%%%%%%%%%%%%%%%%%%%%%%%%%%%%%%%%%%%%%%%%%%%%%%%%%%%%%%%%%%%%%%%%%%%%%%%%%%%%%%%%%%%%%%%%%%%%%%%%%%%%%%%%%%%%%%%%%%%%%%%%%%%%%%%%%%%%%%%%
}
 \vspace{54mm}

























$\therefore  \quad \left(-2, \frac{9}{2}\right) = 2\,\left(1, \frac{3}{2}\right) -\tfrac{1}{2}  (8, -3)$\quad con lo cual  $\vec{u}$ se puede descomponer en las direcciones de  $\vec{v}$ y $\vec{w}$.

\bigskip 

        \item $\vec{v} = (2, -1)$ y $\vec{w} = \left(0, -\tfrac{11}{6} \right)$.

Planteamos la definición de descomposición de descomposición de $\vec{u}$ en la dirección de  $\vec{v}$ y $\vec{w}$ tenemos:
\begin{align*}
    \vec{u} &= \alpha \vec{v} + \beta \vec{w}
    \\
    \left(-2, \tfrac{9}{2}\right) &= \alpha  (2,-1 ) + \beta  (\left(0,\, -\frac{11}{6}\right) 
    \\
    \left(-2, \tfrac{9}{2}\right) &= \left( 2\alpha + 0\cdot\beta  ,\, - \alpha - \tfrac{11}{6}\beta \right)
\end{align*}
Al aplicar la igualdad entre vectores obtenemos el siguiente sistema de ecuación:

\medskip 

$\left\{\begin{array}{rl}
    2\alpha  &= -2 \indicador{x}{(.1,.15)}{y}{(1.3,.15)}
    \\[4mm]
    - \alpha - \frac{11}{6}\beta &= \frac{9}{2} \marka{a}{(.1,.15)}
\end{array}\right.$







\tikz[overlay,remember picture] {
\draw [overlay,->,very thick,red,opacity=.5](x)--(y);
\node  at (y) [ right ]{
$\alpha
= \marka{XX}{(-.25,-.15)}%%%%%
-1
$ \indicador{x}{(.1,.15)}{y}{(1.5,-.25)}\marka{z}{(2.3,-.25)}
};
%%%%%%%%%%%%%%%%%%%%%%%%
\draw [overlay,->,very thick,red,opacity=.5](x)-|(y)--(z);
\draw [overlay,very thick,red,opacity=.5](a)-|(y);
\node  at (z) [text width=40mm, below right , yshift= 8mm ]{ 
\[
\begin{array}{rl}
1 - \frac{11}{6}\beta  
&=%%%%%%%%%%%%%%%%%%%%%%
- \frac{11}{6}\beta  
 \\[3mm]%%%%%%%%%%%%%%%%%%%%%%%%%%%%%%%%%%%%%%%%%%%%%%%%%%%%%%%%%%%%%%%%%%
- \frac{11}{6}\beta 
  &=%%%%%%%%%%%%%%%%%%%%%%
  \frac{9}{2} - 1
  \\[3mm]%%%%%%%%%%%%%%%%%%%%%%%%%%%%%%%%%%%%%%%%%%%%%%%%%%%%%%%%%%%%%%%%%%
- \frac{11}{6}\beta 
  &=%%%%%%%%%%%%%%%%%%%%%%
   \frac{7}{2}
  \\[3mm]%%%%%%%%%%%%%%%%%%%%%%%%%%%%%%%%%%%%%%%%%%%%%%%%%%%%%%%%%%%%%%%%%% 
 \beta  
  &=%%%%%%%%%%%%%%%%%%%%%%
  \frac{7}{2}: \left(- \tfrac{11}{6}\right)
  \\[4mm]%%%%%%%%%%%%%%%%%%%%%%%%%%%%%%%%%%%%%%%%%%%%%%%%%%%%%%%%%%%%%%%%%% 
  \marka{x1}{(-.1,-.2)}
  \beta
  &=  
   -\frac{21}{11} \marka{x2}{(.1,.5)}  
\end{array}
\]
};
%%%%%%%%%%%%%%%%%%%%%%%%%%%%%%%%%%%%%%%%%%%%%%%%%%%%%%%%%%%%%%%%%%%%%%%%%%%%%%%%%%%%%%%%%%%%%%%%%%%%%%%%%%%%%%%%%%%%%%%%%%%%%%%%%%%%%%%%%%%%%%%%%%%%%%%%%%%%%%%%%%%%%%%%%%%%%%%%%%%%
\draw[color=black!50!,rounded corners=1mm] (x1) rectangle (x2);
%%%%%%%%%%%%%%%%%%%%%%%%%%%%%%%%%%%%%%%%%%%%%%%%%%%%%%%%%%%%%%%%%%%%%%%%%%%%%%%%%%%%%%%%%%%%%%%%%%%%%%%%%%%%%%%%%%%%%%%%%%%%%%%%%%%%%%%%%%%%%%%%%%%%%%%%%%%%%%%%%%%%%%%%%%%%%%%%%%%% 
}
 \vspace{27mm} 

$\therefore  \quad \left(-2, \tfrac{9}{2}\right) = -1\,\left(2, 1\right) -\tfrac{21}{11}  \left( 0 , -\tfrac{11}{6} \right) $\quad con lo cual  $\vec{u}$ también se puede descomponer en estas direcciones.











        
    \end{enumerate}  








    \item ¿$\vec{a} = (0, 10, -5)$ es combinación lineal de $\vec{b} = (2, 4, -1)$ y $\vec{c} = (1, -3, 2)$?
Planteamos la definición de descomposición de descomposición de $\vec{u}$ en la dirección de  $\vec{v}$ y $\vec{w}$ y buscamos $\alpha, \beta$ tales que:
$$
\begin{array}{cl}
    \qquad\alpha \, \vec{b} + \beta \, \vec{c} &= \vec{a}
    \\[3mm]
    \alpha(2,\,4,\,-1) + \beta(1,\,-3,\,2) &= (0,10,-5)
    \\[3mm]
    \left( 2\alpha + \beta  ,\, 4 \alpha - 3 \beta ,\, - \alpha + 2\beta \right) &= (0,10,-5) 
\end{array}$$
Al aplicar la igualdad entre vectores obtenemos el siguiente sistema de ecuación:

\medskip 

$\left\{\begin{array}{rl}
    2\alpha + \beta &= 0  \indicador{x}{(.1,.15)}{y}{(1.3,.15)}
    \\[4mm]
    4\alpha - 3\beta &= 10 \marka{a}{(.1,.15)}
    \\[4mm]
    - \alpha + 2\beta &= \marka{xx}{(-.15,-.15)}
    -5
\end{array}\right.$



\tikz[overlay,remember picture] {
\draw [overlay,->,very thick,red,opacity=.5](x)--(y);
\node  at (y) [ right ]{
$\beta
= \marka{XX}{(-.25,-.15)}%%%%%
-2 \alpha 
$ \indicador{x}{(.1,.15)}{y}{(1.5,-.25)}\marka{z}{(3.3,-.25)}
};
%%%%%%%%%%%%%%%%%%%%%%%%
\draw [overlay,->,very thick,red,opacity=.5](x)-|(y)--(z);
\draw [overlay,very thick,red,opacity=.5](a)-|(y);
\node  at (z) [text width=40mm, below right , yshift= 8mm ]{ 
\[
\begin{array}{rl}
4\alpha - 3(-2\alpha) 
&=%%%%%%%%%%%%%%%%%%%%%%
10
 \\[3mm]%%%%%%%%%%%%%%%%%%%%%%%%%%%%%%%%%%%%%%%%%%%%%%%%%%%%%%%%%%%%%%%%%%
4\alpha + 6\alpha  
  &=%%%%%%%%%%%%%%%%%%%%%%
 10
  \\[3mm]%%%%%%%%%%%%%%%%%%%%%%%%%%%%%%%%%%%%%%%%%%%%%%%%%%%%%%%%%%%%%%%%%%
10\alpha  
  &=%%%%%%%%%%%%%%%%%%%%%%
  10
  \\[4mm]%%%%%%%%%%%%%%%%%%%%%%%%%%%%%%%%%%%%%%%%%%%%%%%%%%%%%%%%%%%%%%%%%% 
  \marka{x1}{(-.1,-.2)}
  \alpha
  &= \indicador{YY}{(-.3,-.2)}{ZZ}{(-7,-.4)}\marka{TT}{(-5.5,-1.3)}%%%%%%%%%%%%%%%%%%%%%%
  1 \marka{x2}{(.1,.5)}  
\end{array} 
\]
};
%%%%%%%%%%%%%%%%%%%%%%%%%%%%%%%%%%%%%%%%%%%%%%%%%%%%%%%%%%%%%%%%%%%%%%%%%%%%%%%%%%%%%%%%%%%%%%%%%%%%%%%%%%%%%%%%%%%%%%%%%%%%%%%%%%%%%%%%%%%%%%%%%%%%%%%%%%%%%%%%%%%%%%%%%%%%%%%%%%%%
\draw[color=black!50!,rounded corners=1mm] (x1) rectangle (x2);
%%%%%%%%%%%%%%%%%%%%%%%%%%%%%%%%%%%%%%%%%%%%%%%%%%%%%%%%%%%%%%%%%%%%%%%%%%%%%%%%%%%%%%%%%%%%%%%%%%%%%%%%%%%%%%%%%%%%%%%%%%%%%%%%%%%%%%%%%%%%%%%%%%%%%%%%%%%%%%%%%%%%%%%%%%%%%%%%%%%%
\draw [overlay,->,very thick,blue,opacity=.5](YY)|-(ZZ)|-(TT);
\draw [overlay,very thick,blue,opacity=.5](XX)|-(ZZ);
\node  at (TT) [text width=23mm, below right , yshift= 9mm ]{
\[
\begin{array}{rl}
 \beta
 &=%%%%%%%%%%%%%%%%%%%%%%
  -2 \cdot 1  
 \\[4mm]%%%%%%%%%%%%%%%%%%%%%%%%%%%%%%%%%%%%%%%%%%%%%%%%%%%%%%%%%%%%%%%%%%
  \marka{a1}{(-.1,-.2)} 
   \beta
  &=\marka{YYY}{(-.15,-.15)} \indicador{ZZ}{(-1,-.5)}{TT}{(.5,-1.4)}%%%%%%%%%%%%%%%%%%%%%%
  -2 \marka{a2}{(.1,.5)}
\end{array}
\]
};
\draw[color=black!50!,rounded corners=1mm] (a1) rectangle (a2);
%%%%%%%%%%%%%%%%%%%%%%%%%%%%%%%%%%%%%%%%%%%%%%%%%%%%%%%%%%%%%%%%%%%%%%%%%%%%%%%%%%%%%%%%%%%%%%%%%%%%%%%%%%%%%%%%%%%%%%%%%%%%%%%%%%%%%%%%%%%%%%%%%%%%%%%%%%%%%%%%%%%%%%%%%%%%%%%%%%%%
\draw [overlay,->,very thick,teal,opacity=.5](YY)++(.1,0)|-(ZZ)|-(TT);
\draw [overlay,very thick,teal,opacity=.5](YYY)|-(ZZ);
\draw [overlay,very thick,teal,opacity=.5](xx)++(0,-.1)to[out=-60,in=180]node[pos=.6, below left,text width=57mm]{\footnotesize Sustituimos $\alpha = 1$ y $\beta=-2$ en la tercer ecuación:}(ZZ);
\node  at (TT) [right ]{
$-1 + 2(-2) = -1 - 4 = -5 \quad \checkmark$
};
}
 \vspace{55mm}

  Por lo tanto, para $\alpha = 1$ y $\beta = -2$, se satisfacen las tres ecuaciones del sistema. Es decir, estos son los escalares que estábamos buscando. En consecuencia, $\vec{a}$ {\bf sí} es una combinación lineal de $\vec{b}$ y $\vec{c}$. 

\end{enumerate}









\item Encontrar las componentes cartesianas de cada vector. Datos: $\|\vec{F}_1\| = 100$ \quad y \quad $\|\vec{F}_2\| = 120$.


\vspace{-3mm}
\begin{center}
\begin{tikzpicture}[scale=.3, >=stealth]

    % Ejes coordenados
    \draw[->] (-10,0) -- (10,0) node[right] {$x$};
    \draw[->] (0,-1) -- (0,10) node[above] {$y$};

    % Vectores
    \draw[->, thick] (0,0) -- (60:10) node[above right] {$\vec{F}_1$};
    \draw[->, thick] (0,0) -- (140:12) node[above left] {$\vec{F}_2$};

    % Ángulos
    \draw[->, dashed, blue] (60:5) arc[start angle=60,end angle=90,radius=5] node[pos=.4,above]  {30$^{\circ}$}; 

    \draw[->, dashed, blue] (180:5) arc[start angle=180,end angle=140,radius=5]node[pos=.4,left] {$40^{\circ}$};
    
\end{tikzpicture}
\end{center}

\vspace{-3mm}





\textbf{Solución:}

Las componentes de un vector $\vec{u}$ con magnitud $\|\vec{u}\|$ y un ángulo $\theta$ respecto al eje $x$ positivo se calculan como:
\[
\vec{u} = \|\vec{u}\| \big(\cos\theta ,\  \sin\theta \,  \big),
\]
$\theta$ se conoce como  el ángulo principal del vector.

En lo que sigue marcaremos los ángulos $\theta$ respecto al eje $x$ positivo 

\vspace{-3mm}
\begin{center}
\begin{tikzpicture}[scale=.35, >=stealth]

    % Ejes coordenados
    \draw[->] (-10,0) -- (10,0) node[right] {$x$};
    \draw[->] (0,-1) -- (0,10) node[above] {$y$};

    % Vectores
    \draw[->, thick] (0,0) -- (60:10) node[above right] {$\vec{F}_1$};
    \draw[->, thick] (0,0) -- (140:12) node[above left] {$\vec{F}_2$};

    % Ángulo 
    \draw[->, dashed, blue] (60:5) arc[start angle=60,end angle=90,radius=5] node[pos=.4,above]  {30$^{\circ}$}; 

    % Ángulo theta
    \draw[->, dashed, teal] (0:7.8) arc[start angle=0,end angle=60,radius=7.8] node[pos=.4,right, teal]  {$\theta_1 =60^{\circ}$}; 

    % Ángulo 
    \draw[->, dashed, blue] (180:5) arc[start angle=180,end angle=140,radius=5]node[pos=.4,left] {$40^{\circ}$};

    % Ángulo theta
    \draw[->, dashed, red] (0:2.3) arc[start angle=0,end angle=140,radius=2.3]node[pos=.25,right,red] {$\theta_2 =140^{\circ}$};

    
\end{tikzpicture}
\end{center}

\vspace{-3mm}

\begin{enumerate}[a)]
\begin{multicols}{2}
    \item Para el vector del inciso (a), tenemos:
    \[
        \|\vec{u}\| =100, \quad \theta = 60^\circ.
    \]
    Las componentes son:
    \begin{align*}
    \vec{u}  &= 100 (\cos 60^\circ \, , \sin 60^\circ ).
    \\ 
    &= 100 \left(\tfrac{1}{2} \, ,\  \tfrac{\sqrt{3}}{2}  \right) 
    \\
    &=  \left(50, 50\sqrt{3} \right).
    \end{align*}

    \item Para el vector del inciso (b), tenemos:
    \[
    \|\vec{u}\| = 10, \quad \theta = 140^\circ.
    \]
    Las componentes son:
    \begin{align*}
    \vec{u} &= 120 \left(\cos 140^\circ \, ,\ \sin 140^\circ \,  \right) \\
    &= 120 \left(-0.766\, ,\ 0,642\right) 
    \\
    &= {\bf \big(}-91,925 {\bf ;} \, 77,134  {\bf \big)}
    \end{align*}
\end{multicols}
\end{enumerate}




\medskip





















\item Considerando los vectores del plano $\vec{u}$ y $\vec{v}$:

\begin{enumerate}[a)]
    \item Calcular el producto escalar entre ambos vectores sabiendo que:
    \begin{enumerate}[i)]
    \begin{minipage}{.35\textwidth}
    
        \item $\vec{u} = (5,1)$ y $\vec{v} = (-8,3)$.
    \end{minipage}\quad\begin{minipage}{.6\textwidth}
        \item $\|\vec{u}\| = 5$, $\|\vec{v}\| = 4$ y entre ellos se forma un ángulo $\theta = \tfrac{\pi}{3}$.
    
    \end{minipage}
    \end{enumerate}

    \item Si $\vec{u} = (-2, \alpha)$ y $\vec{v} = (\beta,3)$, hallar el valor de $\alpha$, $\beta \in \mathbb{R}$ tal que $\vec{u} \cdot \vec{v} = 2$ y $\|\vec{u}\| = \sqrt{20}$.
\end{enumerate}







{\bf Solución}


Recordemos: 


\medskip 

{\bf \large Producto Escalar Entre Dos Vectores}

\medskip 

\noindent\fbox{\parbox{\textwidth}{
\textbf{Definición:} Sean $\vec{u}$ y $\vec{v}$ vectores y $\theta = \text{ang}(\vec{u}, \vec{v})$, entonces el producto escalar entre los vectores $\vec{u}$ y $\vec{v}$ es un número real y se define como:
\[
\vec{u} \cdot \vec{v} = 
\begin{cases}
\|\vec{u}\| \cdot \|\vec{v}\| \cdot \cos \theta & \text{si } \vec{u} \neq \vec{0} \text{ y } \vec{v} \neq \vec{0} \\
0 & \text{si } \vec{u} = \vec{0} \text{ ó } \vec{v} = \vec{0}
\end{cases}
\]
}
}

\medskip

{\bf \large Producto Escalar En Función De Sus Componentes}

\medskip

\noindent\fbox{\parbox{\textwidth}{
\textbf{Teorema:} Sean $\vec{u}$ y $\vec{v}$ vectores entonces:

\begin{itemize}
    \item[$\blacklozenge$] Si $\vec{u}, \vec{v} \in \mathbb{R}^2$, $\vec{u} = (u_1, u_2)$ y $\vec{v} = (v_1, v_2)$, el producto escalar es el número real definido por:
    \[
    \vec{u} \cdot \vec{v} = (u_1, u_2) \cdot (v_1, v_2) = u_1 \cdot v_1 + u_2 \cdot v_2
    \]

    \item[$\blacklozenge$] Si $\vec{u}, \vec{v} \in \mathbb{R}^3$, $\vec{u} = (u_1, u_2, u_3)$ y $\vec{v} = (v_1, v_2, v_3)$, el producto escalar será el número real definido por:
    \[
    \vec{u} \cdot \vec{v} = (u_1, u_2, u_3) \cdot (v_1, v_2, v_3) = u_1 \cdot v_1 + u_2 \cdot v_2 + u_3 \cdot v_3
    \]
\end{itemize}
}
}
 



\begin{enumerate}[a)]
    \item Calcular el producto escalar entre ambos vectores sabiendo que:

   \medskip 
\begin{enumerate}[i) ]
\begin{minipage}{.45\textwidth}

    \item Producto escalar entre $\vec{u} = (5, 1)$ y $\vec{v} = (-8, 3)$
   
    \begin{align*}
    \vec{u} \cdot \vec{v} &= 
    (5\marka{x1}{(-.15,-.05)} , 1 \marka{y1}{(-.15,-.05)}) 
    \cdot 
    (-8 \marka{x2}{(-.15,-.05)},3\marka{y2}{(-.15,-.05)})
    \\[2mm]
    &= 5\cdot (-8) + 1 \cdot 3 
    \\
    &= -40 + 3 
    \\
    &= -37 
    \end{align*}
    \tikz[overlay,remember picture] {
  \draw[thick,blue,opacity=.3, line width=1pt] (x1) to[out=-85,in=200] (x2) ;
  \draw[thick,teal,opacity=.3, line width=1pt] (y1) to[out=-20,in=285] (y2) ;
} 


    Por lo tanto,  $\vec{u} \cdot \vec{v} = -37.$
     
\end{minipage}\quad\qquad\begin{minipage}{.45\textwidth}

    \item  Producto escalar usando la fórmula con el ángulo entre los vectores:    
    \begin{align*}
    \vec{u} \cdot \vec{v} &= \|\vec{u}\| \cdot  \|\vec{v}\| \cdot  cos\left(\frac{\pi}{3}\right)
    \\
    &= 5 \cdot 4 \cdot \frac{1}{2}.
    \\
    & = 10. 
    \end{align*}

     Por lo tanto, $\vec{u} \cdot \vec{v} = 10.$\\
     \,
\end{minipage}
\end{enumerate}

\bigskip

    \item Si $\vec{u} = (-2, \alpha)$ y $\vec{v} = (\beta,3)$, hallar el valor de $\alpha,\, \beta \in \mathbb{R}$ tal que: \quad $\vec{u} \cdot \vec{v} = 2 \quad \text{y} \quad \|\vec{u}\| = \sqrt{20}$ 

   \medskip
    
\begin{itemize}

\begin{minipage}{.3\textwidth}
    \item  $\vec{a} = (-2, \alpha)$\marka{A}{(.1,.05)} 

\bigskip
    
    \item  $\vec{b} = (\beta, 3)$ \marka{B}{(.1,.15)} \indicador{X}{(.8,.85)}{Y}{(2.3,1.25)}

\bigskip
    
    \item  $\vec{a} \cdot \vec{b} = 2$ \marka{C}{(.1,.15)} 

\bigskip
    
    \item $\|\vec{a}\| = \sqrt{20}$ \marka{D}{(.1,.15)} \indicador{Z}{(.9,.65)}{T}{(1.8,.65)}

\end{minipage}
\tikz[overlay,remember picture] {
\draw [overlay,very thick,red,opacity=.5]
(A)++(0,.15)-|(X)|-(C);
\draw [overlay,very thick,->,red,opacity=.5]
(B)-|(X)--(Y);
\node  at (Y) [text width=34mm, below right ,xshift= -9mm,yshift= 9mm]{
\begin{align*}
    \vec{a} \cdot \vec{b} \qquad  &= 2
    \\ 
    (-2\marka{x1}{(-.15,-.05)} , \alpha\marka{y1}{(-.15,-.05)}) 
    \cdot 
    ( \beta \marka{x2}{(-.15,-.05)},   3\marka{y2}{(-.15,-.05)})
    &= 2
    \\[2mm]
     -2\beta + 3\alpha \ \ &= 2 \indicador{a}{(.1,.15)}{b}{(1.2,-.45)} \marka{c}{(1.8,-.45)}
    \end{align*}
    
};
 \draw[thick,blue,opacity=.3, line width=1pt] (x1) to[out=-85,in=200] (x2) ;
  \draw[thick,teal,opacity=.3, line width=1pt] (y1) to[out=-20,in=285] (y2) ;
%%--------------------------------------------------------------------------------------------------------------------
\draw [overlay,very thick,teal,opacity=.5]
(A)-|(Z)--(T);
\draw [overlay,very thick,->,teal,opacity=.5]
(D)-|(Z)--(T);
\node  at (T) [text width=58mm, below right ,xshift= -19mm,yshift= 9mm]{
\begin{align*}
    \|(-2,\alpha)\|&= \sqrt{20}
    \\[2mm]
    \sqrt{(-2)^2 + \alpha^2}   &= \sqrt{20} 
    \\[2mm]
    \left( \sqrt{(-2)^2 + \alpha^2} \right)^2  &= \left(  \sqrt{20} \right)^2
    \\[2mm]
    4 + \alpha^2 &= 20
    \\[2mm]
    \alpha^2 &= 16 
    \\[2mm]
    \alpha &= \pm 4 \marka{d}{(.1,.15)}
    \end{align*}  
    };
%%-------------------------------------------------------------------------------------------------------------------- 
\draw [overlay,very thick,purple,opacity=.5]
(a)-|(b)|-(d); 
\draw [overlay,->,very thick,purple,opacity=.5]
(b)--(c);
%%--------------------------------------------------------------------------------------------------------------------
\node  at (c) [text width=66mm, below right ,xshift= 5mm,yshift= 9mm]{
\begin{itemize} 
    \begin{multicols}{2}
        \item Si \marka{a1}{(-.1,-.15)}$\alpha = 4$\marka{a2}{(.1,.4)} : 
        \begin{align*}
        -2\beta + 3\cdot 4 &= 2 
        \\
        -2\beta + 12 &= 2.
        \\
        -2\beta &= -10 
        \\
        \marka{b1}{(-.1,-.15)}\beta& = 5\marka{b2}{(.1,.4)}        
        \end{align*} 

        \item Si \marka{aa1}{(-.1,-.15)}$\alpha = -4$\marka{aa2}{(.1,.4)} :
        \begin{align*}
        -2\beta + 3\cdot (-4) &= 2 
        \\
        -2\beta -12 &= 2.
        \\
        -2\beta &= 14 
        \\
        \marka{bb1}{(-.1,-.15)}\beta &= -7\marka{bb2}{(.1,.4)} 
        \end{align*} 
    \end{multicols}
    \end{itemize}
    };

\draw[color=black!50!] (a1) rectangle (a2);

\draw[color=black!50!] (aa1) rectangle (aa2);

\draw[color=black!50!] (b1) rectangle (b2);

\draw[color=black!50!] (bb1) rectangle (bb2);
}


\end{itemize}



\vspace{36mm}

\textbf{Resultados finales:} \\
    Los valores posibles de $(\alpha, \beta)$ son:\quad  )$\alpha = 4$ y $\beta = 5 ) \quad \Or \quad (\alpha = 4$ y $\beta = 5$)
     



    Entonces, las soluciones posibles son:
    \begin{itemize}
        \item Si $\alpha = 4$ y $\beta = 5$ tenemos: $\vec{u}=(-2,4)$ y $\vec{v}= (5,3)$.
        
        \item Si $\alpha = -4$ y $\beta = -7$ tenemos: $\vec{u}=(-2,-4)$ y $\vec{v}= (-7,3)$
    \end{itemize}


    
\end{enumerate} 






\item En cada caso, hallar el ángulo que forman los vectores $\vec{r}$ y $\vec{s}$ y la $\text{Proy}_{\vec{r}} \vec{s}$:

\begin{enumerate}[a)] 
\begin{multicols}{3}
    \item $\vec{r} = (1, \sqrt{3})$ y $\vec{s} = (-2, 2\sqrt{3})$.
    \item $\vec{r} = (2,1,1)$ y $\vec{s} = (1,-2,1)$.
    \item $\vec{r} = (1,2,1)$ y $\vec{s} = (-2,3,-4)$. 
\end{multicols}
\end{enumerate}  


{\bf Solución:}

\begin{enumerate}[a)]

    \item $\vec{r} = (1, \sqrt{3})$, $\vec{s} = (-2, 2\sqrt{3})$

\begin{minipage}{.55\textwidth}

    Producto escalar:
    \begin{itemize} 
    \item $
    \vec{r} \cdot \vec{s} = (1, \sqrt{3})\cdot (-2, 2\sqrt{3})= 1 \cdot (-2) + \sqrt{3} \cdot 2\sqrt{3} = -2 + 6 = 4
    $

   \hspace{-9mm} Normas:

    \item $
    \|\vec{r}\| = \sqrt{1^2 + (\sqrt{3})^2} = \sqrt{1 + 3} = 2, 
    $

    \item $
    \|\vec{s}\| = \sqrt{(-2)^2 + 2\sqrt{3})^2} = \sqrt{4 + 12} = \sqrt{16} = 4
    $

    \hspace{-9mm}  Ángulo:

    \item  $ \cos \theta = \frac{\vec{r} \cdot \vec{s}}{\|\vec{r}\| \cdot \|\vec{s}\|} = \frac{4}{2 \cdot 4} = \frac{1}{2} \quad\Rightarrow\quad  \theta= \arccos\left( \frac{1}{2} \right)  = \frac{\pi}{3}
    $

    \hspace{-9mm} Proyección:

    \item $
    \text{Proy}_{\vec{r}} \vec{s} = \frac{\vec{r} \cdot \vec{s}}{\|\vec{r}\|^2} \cdot \vec{r}
    = \frac{4}{4} \cdot (1, \sqrt{3}) = (1, \sqrt{3})
    $
    \end{itemize}

\end{minipage}\begin{minipage}{.45\textwidth}

\begin{center}

\definecolor{gris}{rgb}{0.5,0.5,0.5}
\begin{tikzpicture}[line cap=round,line join=round,>=triangle 45,scale=1.4]
    % Ejes
    \draw [color=gris!10!,dash pattern=on 1pt off 1pt, xstep=1cm,ystep=1cm] (-2.3,-.5) grid (2.3,4.3);
    \draw[->] (-2.3,0) -- (2.3,0) node[below] {\small $x$};
    \draw[->] (0,-0.5) -- (0,4.3) node[left] {\small $y$};
    \foreach \x in {-2 ,-1,1, 2}
\draw[shift={(\x,0)},color=black] (0pt,2pt) -- (0pt,-2pt) node[below] {\tiny $\x$};

\foreach \y in {1,2,...,4}
\draw[shift={(0,\y)},color=black] (2pt,0pt) -- (-2pt,0pt) node[left] {\tiny $\y$};



    \pgfmathsetmacro{\misqrttres}{sqrt(3)}

    \coordinate (A) at (0,0);
    \coordinate (R) at (1,\misqrttres);
    \coordinate (S) at (-2,2*\misqrttres);
    \coordinate (P) at(R);

    % Vectores
    \draw[->, thick, blue] (0,0) -- (R) node[midway, above left,pos=.7] {\small $\vec{r}$};
    \draw[->, thick, red] (0,0) -- (S) node[midway, above] {\small $\vec{s}$};

    % Proyección
    \draw[-> ,line width=1.5pt, teal] (0,0) -- (P) node[midway,  right,pos=.7] {\small $\operatorname{Proy}_{\vec{r}} \vec{s}$};
    \draw[dashed,rounded corners=0mm,opacity=.3] (S) -- (P);

    % Ángulo
     % Dibujar el ángulo
    \pic[draw, color=magenta, fill=magenta!50!,   angle eccentricity=1.3, angle radius=.9cm, "$\theta = \tfrac{\pi}{3}$",opacity=.8] {angle = R--A--S};
\end{tikzpicture}
\end{center}

\end{minipage}

\medskip


    \item $\vec{r} = (2,1,1)$, $\vec{s} = (1,-1,2)$

    Producto escalar:
    \[
    \vec{r} \cdot \vec{s} = (2,1,1)\cdot (1,-1,1)= 2\cdot1 + 1\cdot(-1) + 1\cdot 2 = 2 - 1 +2= 3
    \]

    Normas:
    \[
    \|\vec{r}\| = \sqrt{2^2 + 1^2 + 2^2} = \sqrt{4 + 1 + 1} = \sqrt{6}  , \qquad\qquad \qquad
    \|\vec{s}\| = \sqrt{1^2 + (-2)^2 + 1^2} = \sqrt{1 + 4 + 1} = \sqrt{6}
    \]

    Ángulo:
    \[
    \cos \theta = \dfrac{\vec{r} \cdot \vec{s}}{\|\vec{r}\| \|\vec{s}\|} = \frac{3}{\sqrt{6} \cdot \sqrt{6}} = \frac{3}{6} = \frac{1}{2} \qquad\Longrightarrow\qquad
    \theta = \cos^{-1}\left(\frac{1}{2 }\right) = \frac{\pi}{3}
    \]

    Proyección:
    \[
    \text{Proy}_{\vec{r}} \vec{s} = \frac{\vec{r} \cdot \vec{s}}{\|\vec{r}\|^2} \cdot \vec{r}
    = \frac{3}{6} \cdot (2,1,2) =  \frac{1}{2} \cdot (2,1,1) = \left(1, \tfrac{1}{2}, \tfrac{1}{2}\right)
    \]

    \begin{center}
        
\begin{asypicture}{name=Surful2}
import three;
size(8cm);

// Configuración de la proyección
currentprojection = perspective(.1, 2, 3);

// Vectores
triple A = (0,0,0);
triple r = (2, 1, 1); // Vector r
triple s = (1, -1, 2); // Vector s

// Calcular la proyección de r sobre s
triple p = ((dot(r,s) / dot(r,r)) * r);

// Dibujar ejes
draw((-1,0,0)--(3,0,0), gray+linewidth(0.5), EndArrow3); // Eje X
label("$x$", (3,0,0), SE);
draw((0,-2,0)--(0,2,0), gray+linewidth(0.5), EndArrow3); // Eje Y
label("$y$", (0,2,0), NE);
draw((0,0,-1)--(0,0,3), gray+linewidth(0.5), EndArrow3); // Eje Z
label("$z$", (0,0,3), N);

// Dibujar vectores
draw((0,0,0)--r, blue+linewidth(1),Arrow3); // Vector r
label("$\vec{r}$", r, NW);
draw((0,0,0)--s, red+linewidth(1),Arrow3); // Vector s
label("$\vec{s}$", s, NE);
draw((0,0,0)--p, purple+linewidth(1),Arrow3);  // Proyeccion
label("${Proy}_{\vec{r}} \vec{s}$", p, NW);
draw(s--p,  black+linewidth(0.5)+dashed+opacity(0.2));  // Proyeccion

// Diagrama de referencia
dot((0,0,0));


// Dibujar el arco del ángulo manualmente
triple T = .3*r; // Punto en el vector r
draw(surface(A--arc(A, T,s)--cycle), lightblue+opacity(0.5)); // Arco relleno
draw(arc(A, T, s), gray+linewidth(0.5));
triple tt = .23*r+.23*s;
label("$\theta $", tt, SE);
\end{asypicture}
    \end{center}    

    \item $\vec{r} = (1,2,1)$, $\vec{s} = (-2,3,-4)$

    Producto escalar:
    \[
    \vec{r} \cdot \vec{s} = 1\cdot(-2) + 2\cdot3 + 1\cdot(-4) = -2 + 6 - 4 = 0
    \]

    En este caso como ninguno de los vectores es el vector nulo, tenemos que el ángulo es un recto:
    \[
    \vec{r} \cdot \vec{s} = 0 \quad\Rightarrow \quad\cos \theta = 0 \quad\Rightarrow \quad \theta = \frac{\pi}{2}
    \]

    Como son ortogonales, la proyección es el vector nulo:
    \[
    \text{Proy}_{\vec{r}} \vec{s} = \vec{0}
    \]


\begin{center}
    % --- Ejercicio (c) ---
\begin{asypicture}{name=Surful3}
import three;
size(8cm);

// Configuración de la proyección
currentprojection = perspective(-2, .5, 3);

// Vectores
triple A = (0,0,0); 
triple r = (1, 2, 1); // Vector r
triple s = (-2, 3, -4); // Vector s

// Calcular la proyección de r sobre s
triple p = ((dot(r,s) / dot(s,s)) * s);

// Dibujar ejes
draw((-3,0,0)--(2,0,0), gray+linewidth(0.5), EndArrow3); // Eje X
label("$x$", (2,0,0), NE);
draw((0,-1,0)--(0,4,0), gray+linewidth(0.5), EndArrow3); // Eje Y
label("$y$", (0,4,0), NE);
draw((0,0,-5)--(0,0,2), gray+linewidth(0.5), EndArrow3); // Eje Z
label("$z$", (0,0,2), N);

// Dibujar vectores
draw((0,0,0)--r, blue+linewidth(1),Arrow3); // Vector r
label("$\vec{r}$", r, NW);
draw((0,0,0)--s, red+linewidth(1),Arrow3); // Vector s
label("$\vec{s}$", s, NE);
draw((0,0,0)--p, purple+linewidth(1),Arrow3);  // Proyeccion
//label("${Proy}_{\vec{s}} \vec{r}$", p, NW);
//draw(r--p,  red+linewidth(0.5)+dashed+opacity(0.2));  // Proyeccion

// Diagrama de referencia
dot((0,0,0));


// Dibujar el arco del ángulo manualmente
triple T = .3*r; // Punto en el vector r
draw(surface(A--arc(A, T,s)--cycle), lightblue+opacity(0.5)); // Arco relleno
draw(arc(A, T, s), gray+linewidth(0.5));
triple tt = .23*r+.1*s;
label("$\theta=90^\circ $", tt, SE);
\end{asypicture}
\end{center}


\end{enumerate}

































\item En Física se considera al {\bf trabajo ($W$)} realizado sobre una partícula por una fuerza $\vec{F}$ aplicada en un desplazamiento $\vec{r}$ como $W = \vec{F} \cdot \vec{r}$.

    \begin{itemize}
        \item[a)] Calcular el trabajo realizado por la fuerza de componentes $(1, -2, 1)$ al desplazar un sólido desde un punto $A = (3, 2, -1)$ hasta $B = (2, 3, 1)$.
        \item[b)] Si una fuerza aplicada a un sólido que se desplaza una cierta distancia realizó un trabajo nulo, ¿qué puede decir de las direcciones de los vectores que representan la fuerza y el desplazamiento?
        \item[c)] Determinar el trabajo realizado por una fuerza de $5\,\text{N}$ que actúa en la dirección del vector $\vec{u} = (1, 1)$ al mover un objeto una distancia de un metro desde $P = (0, 0)$ hasta $Q = (1, 0)$.
    \end{itemize}


\textbf{Solución:}

Recordar: La unidad de medida del trabajo en el Sistema Internacional de Unidades (SI) es el Joule (J).

\bigskip 

\begin{itemize}
\begin{minipage}{.43\textwidth}
\item[a)] Trabajo realizado por  $\vec{F} = (1, -2, 1)$: \\
    El desplazamiento $\vec{r}$ se calcula como punto final menos punto inicial:
    \begin{align*}
            \vec{r} &= \vec{B} - \vec{A} 
            \\ 
            &= (2, 3, 1) - (3, 2, -1) 
            \\
            &= (-1, 1, 2).
    \end{align*}
    
\, \\  
\, \\  
\end{minipage}\qquad\begin{minipage}{.43\textwidth}
  El trabajo realizado es:    
  \begin{align*}
    W &= \qquad\qquad \vec{F} \cdot \vec{r} 
    \\ 
    &= \quad\   (\,1\marka{x1}{(-.1,-.05)},\, -2\marka{y1}{(-.15,-.05)},\, 1\marka{z1}{(-.15,-.05)}) \cdot (-1\marka{x2}{(-.15,-.05)},\, 1\marka{y2}{(-.15,-.05)},\, 2\marka{z2}{(-.15,-.05)}) 
    \\[3mm]
    &= 1\cdot(-1) + (-2)\cdot 1+ 1\cdot 2 
    \\[2mm]
    &= -1 - 2 + 2 
    \\[2mm]
    &= -1.
    \end{align*} 
    \tikz[overlay,remember picture] {
  \draw[thick,blue,opacity=.3, line width=1pt] (x1) to[out=-95,in=200] (x2) ;
  \draw[thick,purple,opacity=.3, line width=1pt] (y1) to[out=-50,in=230] (y2) ;
  \draw[thick,teal,opacity=.3, line width=1pt] (z1) to[out=-20,in=285] (z2) ;
} 

    
\end{minipage}\begin{comment}
\qquad\begin{minipage}{.3\textwidth}
   
  
\begin{asy}
import three;
size(6cm);

// Configuración de la proyección
currentprojection = perspective(3, 3, 2);

// Puntos y vectores
triple A = (3, 2, -1);
triple B = (2, 3, 1); 
triple F = (1, -2, 1); // Vector fuerza 
triple X = B+(1, -2, 1);
triple M; = .7*(3, 2, -1)+.3*(2, 3, 1);
triple Y = M+(1, -2, 1);

// Dibujar ejes
draw((-1,0,0)--(4,0,0), gray+linewidth(0.5), EndArrow3); // Eje X
draw((0,-1,0)--(0,4,0), gray+linewidth(0.5), EndArrow3); // Eje Y
draw((0,0,-2)--(0,0,2), gray+linewidth(0.5), EndArrow3); // Eje Z

// Dibujar puntos
dot(A, blue); label("$A$", A, W);
dot(B, blue); label("$B$", B, E); 

// Dibujar vectores
draw(A--B, blue+linewidth(1)); label("$\vec{r}$", midpoint(A--B), SE); 
draw((0,0,0)--F, purple+linewidth(1), EndArrow3); label("$\vec{F}$", F, NE);

draw(B--X, purple+linewidth(1)+opacity(0.5), EndArrow3); label("$\vec{F}$", F, NE);
draw(M--Y, purple+linewidth(1)+opacity(0.5), EndArrow3); label("$\vec{F}$", F, NE);
\end{asy}

\end{minipage}
    
\end{comment}

\medskip 

    Por lo tanto, el trabajo realizado es: \quad$ W = -1\,\text{J}$.
    

    \item[b)] Trabajo nulo: \\
    Si el trabajo realizado es $W = 0$, entonces:
    \[
    \vec{F} \cdot \vec{r} = 0.
    \]
    Esto implica que los vectores $\vec{F}$ y $\vec{r}$ son ortogonales, si los dos son distintos del nulo tenemos que son perpendiculares (forman un ángulo de $90^\circ$).


\bigskip 

    \item[c)] El trabajo realizado por una fuerza de $5\,\text{N}$ en la dirección de $\vec{u}$ lo calcularemos en lo que sgue. Primero, determinaremos la fuerza mencionada. Dado que se trata de una fuerza de $5\,\text{N}$ en la dirección de $\vec{u}$, esta se calcula como el producto de $5$ por el versor en la dirección y sentido de $\vec{u}$.  \\



    
    Primero, normalizamos $\vec{u}$ para obtener un versor:\quad $ \|\vec{u}\| = \sqrt{1^2 + 1^2} = \sqrt{2}.$ 

    
    El versor es: \qquad ${\bm{\check{u}}} = \dfrac{\vec{u}}{\|\vec{u}\|} = \left(\tfrac{1}{\sqrt{2}}, \tfrac{1}{\sqrt{2}}\right).$ 

    
    En consecuencia: \quad $\vec{F} = 5 {\bm{\check{u}}} = 5\left(\tfrac{1}{\sqrt{2}}, \tfrac{1}{\sqrt{2}}\right) = \left(\tfrac{5}{\sqrt{2}}, \tfrac{5}{\sqrt{2}}\right).$ 

    
    El desplazamiento es:
    \[
    \vec{r} = \vec{Q} - \vec{P} = (1, 0) - (0, 0) = (1, 0).
    \]
    El trabajo realizado es:
    \[
    W = \vec{F} \cdot \vec{r} = \left(\tfrac{5}{\sqrt{2}}\right)(1) + \left(\tfrac{5}{\sqrt{2}}\right)(0) = \tfrac{5}{\sqrt{2}}.
    \]
    Por lo tanto:
    \[
    W = \tfrac{5}{\sqrt{2}}\,\text{J}.
    \]
\end{itemize}



\begin{comment}
    
\textbf{Gráfico de la situación:}

\begin{tikzpicture}[scale=1.5, >=triangle 45]

    % Ejes coordenados
    \draw [color=gris!10!,very thin] (-.5,-.5) grid (4.3,4.3);
    \draw[->] (-0.5,0) -- (4.3,0) node[below right] {\small $x$};
    \draw[->] (0,-0.5) -- (0,4.3) node[left] {\small $y$};
    
    % Puntos P y Q
    \fill[blue] (0,0) circle (0.03) node[below left] {\small $P(0,0)$};
    \fill[blue] (1,0) circle (0.03) node[below right] {\small $Q(1,0)$};
    
    % Vector r (desplazamiento)
    \draw[->, thick, red] (0,0) -- (1,0) node[midway, above,pos=.9] {\small $\vec{r}$};
    
    % Vector F (fuerza)
    \pgfmathsetmacro{\minum}{5*sqrt(2)/2}
    \draw[->, thick, blue] (0,0) -- (\minum,\minum) node[midway, above left] {\small $\vec{F} = \left(\tfrac{5}{\sqrt{2}}, \tfrac{5}{\sqrt{2}}\right)$};
    
    % Vector u (dirección inicial)
    \draw[->, dashed, purple] (0,0) -- (1,1) node[midway, above left,pos=.9] {\small $\vec{u}$};
    
     

    % Etiquetas adicionales
    \node at (0.5,-0.2) {\small $\vec{r} = \vec{Q} - \vec{P}$};
    
\end{tikzpicture}


\end{comment}












  \item Sea el vector $\vec{u} = (3, -4)$. Hallar otro vector:

    \begin{enumerate}[a) ]
        \item Perpendicular y con el mismo módulo que el vector dado.
        \item Ortogonal al vector dado y que tenga módulo tres.
    \end{enumerate}

   \textbf{Resolución:}

\textbf{Módulo del vector $\vec{u}$:}
\[
\|\vec{u}\| = \sqrt{3^2 + (-4)^2} = \sqrt{9 + 16} = \sqrt{25} = 5.
\]

\begin{enumerate}[a) ]
    \item \textbf{ Vector perpendicular y con el mismo módulo:}

Un vector perpendicular a $\vec{u} = (3, -4)$ puede obtenerse intercambiando las componentes y cambiando el signo de una de ellas. Por lo tanto:
\[
\vec{v} = (4, 3).
\]
Verificamos que $\vec{v}$ es perpendicular a $\vec{u}$ mediante el producto escalar:
\[
\vec{u} \cdot \vec{v} = (3)(4) + (-4)(3) = 12 - 12 = 0.
\]
Además, el módulo de $\vec{v}$ es:
\[
\|\vec{v}\| = \sqrt{4^2 + 3^2} = \sqrt{16 + 9} = \sqrt{25} = 5.
\]

Por lo tanto, $\vec{v} = (4, 3)$ cumple las condiciones.

\medskip

\begin{minipage}{.6\textwidth}

\item {\bf Vector ortogonal a $\vec{u}$ y con módulo tres:}

Primero, obtenemos el versor perpendicular a $\vec{u}$:
\[
{\bm{\check{v}}}  = \frac{\vec{v}}{\|\vec{v}\|} = \frac{(4, 3)}{5} = \left(\frac{4}{5}, \frac{3}{5}\right).
\]
Multiplicamos este versor por 3 para obtener el vector con \linebreak módulo tres:
\[
\vec{w} = 3 \cdot \vec{v}_\perp = 3 \cdot \left(\frac{4}{5}, \frac{3}{5}\right) = \left(\frac{12}{5}, \frac{9}{5}\right).
\]

Por lo tanto, $\vec{w} = \left(\frac{12}{5}, \frac{9}{5}\right)$ cumple las condiciones.

\textbf{Resultados:}
\begin{itemize}
    \item[a)] Uno de los vector perpendicular con el mismo módulo es: $\vec{v} = (4, 3)$.
    \item[b)] Uno de los vector ortogonal con módulo tres es: $\vec{w} = \left(\frac{12}{5}, \frac{9}{5}\right)$.
\end{itemize}
\end{minipage}\qquad\begin{minipage}{.37\textwidth}

\textbf{Gráfico de la situación:}

\begin{center}
\definecolor{gris}{rgb}{0.5,0.5,0.5}
\begin{tikzpicture}[scale=.8, , >=triangle 45]

    % Ejes coordenados
    \draw [color=gris!10!,very thin] (-1,-4.3) grid (4.3,4.3);
    \draw[->] (-1,0) -- (5.3,0) node[below] {\small $x$};
    \draw[->] (0,-4.3) -- (0,4.3) node[left] {\small $y$};

    % Vector u
    \draw[->, thick, blue] (0,0) -- (3,-4) node[midway, above right] {\small $\vec{u} = (3, -4)$};

    % Vector v (perpendicular con mismo módulo)
    \draw[->, thick, red] (0,0) -- (4,3) node[midway, above right,pos=1.01] {\small $\vec{v} = (4, 3)$};

    % Vector w (ortogonal con módulo 3)
    \draw[->, thick, purple] (0,0) -- (2.4,1.8) node[midway, above left,pos=.9] {\small $\vec{w} = \left(\frac{12}{5}, \frac{9}{5}\right)$};

    % circulo en los puntos
    \clip(-1,-4.3) rectangle (5.3,4.3);
    \draw[blue!30!,dash pattern=on 1pt off 1pt,] (0,0) circle (5);
    %\fill[blue] (3,-4) circle (0.05);
    %\fill[red] (4,3) circle (0.05);
    %\fill[purple] (2.4,1.8) circle (0.05);
\end{tikzpicture}
\end{center}
\end{minipage}
    \end{enumerate}














    

























    

\item Dados los vectores $\vec{u} = \jv + \kv$ y $\vec{v} = (1, -1, 2)$:

\begin{enumerate}[a)]
    \item Hallar un vector perpendicular a $\vec{u}$ y que sea unitario. ¿Es único?

    \item Hallar un vector perpendicular a $\vec{u}$ y de norma cinco.

    \item Hallar todos los vectores perpendiculares a $\vec{u}$ y $\vec{v}$.

    \item Hallar un vector perpendicular a ambos vectores.

    \item Hallar un vector perpendicular a ambos y de módulo igual a $4\sqrt{11}$. ¿Es único este vector?

    \item Determinar si existe $\vec{w} = (w_1, w_2, w_3)$ tal que verifique simultáneamente: \  $\vec{w} \perp \vec{u}, \  \|\vec{w}\| = 3\sqrt{3}$  y   $w_2 = 4.$
\end{enumerate}



 

\begin{enumerate}[a)]
    \item \textbf{Hallar un vector perpendicular a $\vec{u}$ y que sea unitario. ¿Es único?}

Buscamos un vector $\vec{w} = (x, y, z)$ tal que:\quad $ \vec{w} \perp \vec{u}  \quad \text{y} \quad \|\vec{w}\| = 1$

\noindent Resolvamos esta situación:

\bigskip 


\begin{itemize} 

\begin{minipage}{.7\textwidth}
    \item Como $\vec{w} \perp \vec{u} \ \Longrightarrow \ 
     \vec{w} \cdot \vec{u} =\marka{A}{(-.2,-.15)}\indicador{X}{(-.2,-.6)}{Y}{(2.3,-.9)}  0 , \qquad \vec{u} \neq \vec{\bf O}
     \qquad$ y $\quad  \vec{w} \neq \vec{\bf O}
    $ 

\bigskip
    
    \item  $\vec{w} = (x, y, z)$ \marka{B}{(.1,.15)} 

\bigskip
    
    \item  $\vec{u}= \jv +\kv $

    \, \hspace{2mm}$= (0, 1, 0) + (0,0,1)$
    
    \, \hspace{2mm}$= (0, 1, 1)$ \marka{C1}{(.1,.15)} \marka{C2}{(.6,.15)} 

\bigskip
    
    \item $\|\vec{w}\| = \marka{D}{(-.15,-.15)}
    1$ \ 

\end{minipage}
\tikz[overlay,remember picture] {
\draw [overlay,very thick,red,opacity=.5]
(A)-|(X)to[out=-90,in=0](C2)--(C1);
\draw [overlay,very thick,->,red,opacity=.5]
(B)-|(X)--(Y);
\node  at (Y) [text width=80mm, below right ,xshift= -29mm,yshift= 7mm]{
\begin{align*}
    \vec{w} \cdot \vec{u} \qquad  &= 0
    \\ 
    ( x,y,z) 
    \cdot 
    (0,1,1)
    &=0
    \\[2mm]
     y+z &= 0
     \\[2mm]
     y &= -z
    \end{align*}
\ \ Por lo tanto, todo vector ortogonal a $\vec{u}$ tiene la forma:
$$ \marka{E}{(-.1,.15)}\vec{w}= (x,-z,z), \text{ con } x,\, z\in \R$$
\noindent de tal forma que $x$ y $z$ no son simultáneamente nulos.
};
}


\end{itemize}



\vspace{40mm}

\marka{F}{(3,.05)}


\tikz[overlay,remember picture] {
\draw [overlay,very thick,teal,opacity=.5]
(E)-|(D);
\draw [overlay,very thick,->,teal,opacity=.5]
(D)|-(F);
\node  at (F) [text width=.75\textwidth, below right ,xshift= 0mm,yshift=3mm]{
Podemos tomar la anterior opción de $\vec{w}$ y hallarle el versor asociado, esa  seria la respuesta: 
\begin{align*}
    \vec{w} &= \frac{1}{\,\sqrt{x^2+(-z)^2+z^2\,}\,} (x,-z,z)
    \\
    &=\frac{1}{\,\sqrt{x^2+2z^2\,}\,} (x,-z,z)
\end{align*} 
};
}
\vspace{27mm} 

\begin{itemize}
    \item Tomamos por ejemplo $x = 1, \ z = 0$, entonces tenemos:
    \begin{align*}
        \vec{w}  
        &= 
        \frac{1}{\,\sqrt{1^2+2\cdot 0^2\,}\,} (1,-0,0)
        \\
        &= (1,0,0)
    \end{align*}

    \item Tomamos por ejemplo $x = 0, \ z = 1$, entonces tenemos:
    \begin{align*}
        \vec{w}  
        &= 
        \frac{1}{\,\sqrt{0^2+2\cdot 1^2\,}\,} (0,-1,1)
        \\
        &= \frac{1}{\,\sqrt{2\,}\,} (0,-1,1)
        \\
        &= (0,-\tfrac{1}{\,\sqrt{2\,}\,},\tfrac{1}{\,\sqrt{2\,}\,})
    \end{align*}
    

    \item Tomamos por ejemplo $x = 1, \ z = 1$, entonces tenemos:
    \begin{align*}
        \vec{w}  
        &= 
        \frac{1}{\,\sqrt{1^2+2\cdot 1^2\,}\,} (1,-1,1)
        \\
        &= \frac{1}{\,\sqrt{3\,}\,} (0,-1,1)
        \\
        &= (\tfrac{1}{\,\sqrt{3\,} \,},-\tfrac{1}{\,\sqrt{3\,}\,},\tfrac{1}{\,\sqrt{3\,}\,} )
    \end{align*}
\end{itemize}
      
    No existe un único vector que cumpla con lo solicitado, ya que, como vimos, hay infinitos vectores unitarios ortogonales a $\vec{u}$. Basta con elegir cualquier par de valores $x$ y $z$ que no sean simultáneamente nulos, y al normalizar el vector resultante se obtiene una nueva opción para $\vec{w}$. En nuestro ejemplo mostramos tres posibilidades distintas para $\vec{w}$, todas ellas válidas y hay más opciones.

    \item \textbf{Hallar un vector perpendicular a $\vec{u}$ y de norma cinco.}

    A partir de cualquiera de los versores hallados en el punto anteriores, lo multiplicamos por $5$ y obtenemos un vector que cumple lo pedido:
    \[
    \vec{w} = 5 \cdot \left( 1, 0, 0 \right) = \left( 5,0,0 \right)
    \]

    \item \textbf{Hallar todos los vectores perpendiculares a $\vec{u}$ y $\vec{v}$.}

    Buscamos $\vec{w} = (x, y, z)$ tal que: $ \vec{w} \perp \vec{u}  \quad \text{y} \quad \vec{w} \perp \vec{v}$

\noindent Resolvamos esta situación:

\bigskip 


\begin{itemize} 

    \item Como $\vec{w} \perp \vec{u} \ \Longrightarrow \ 
     \vec{w} \cdot \vec{u} =\marka{A}{(-.2,-.15)}\indicador{X}{(-.2,-.6)}{Y}{(2.3,-.9)}  0 , \qquad \vec{u} \neq \vec{\bf O}
     \qquad$ y $\quad  \vec{w} \neq \vec{\bf O}
    $ 

\bigskip
    
    \item Como $\vec{w} \perp \vec{v} \ \Longrightarrow \ 
     \vec{w} \cdot \vec{v} =\marka{A}{(-.2,-.15)}\indicador{X}{(-.2,-.6)}{Y}{(2.3,-.9)}  0 , \qquad \vec{v} \neq \vec{\bf O}
     \qquad$ y $\quad  \vec{w} \neq \vec{\bf O}
    $ 


\end{itemize}

Con lo cual tenemos: 
    \[
    \vec{w} \cdot \vec{u}= (x,y,z) \cdot (0,1,1) = y + z = 0 , \qquad \vec{w} \cdot \vec{v} = (x,y,z) \cdot (1,-1,2) = x - y + 2z = 0 \quad \text{y} \quad  \vec{w}\neq \vec{\bf O}.
    \]
Es decir tenemos que resolver el siguiente sistema de ecuaciones: 

\bigskip

$\left\{ \begin{array}{rl}
    y+z &=  0  \indicador{x}{(.1,.15)}{y}{(1.3,.15)}
    \\[4mm]
     x - y + 2z &= 0 \marka{a}{(.1,.15)}
\end{array}\right.$




\tikz[overlay,remember picture] {
\draw [overlay,->,very thick,red,opacity=.5](x)--(y);
\node  at (y) [ right ]{\marka{a1}{(-.1,-.2)}  
$y
= \marka{XX}{(-.25,-.15)}%%%%%
-z
$\marka{a2}{(.1,.5)} \indicador{x}{(.1,.15)}{y}{(1.5,-.25)}\marka{z}{(2.3,-.25)}
}; 
\draw[color=black!50!,rounded corners=1mm] (a1) rectangle (a2);
%%%%%%%%%%%%%%%%%%%%%%%%%%%%%%%%%%%%%%%%%%%%%%%%%%%%%%%%%%%%%%%%%%%%%%%%%%%%%%%%%%%%%%%%%%%%%%%%%%%%%%%%%%%%%%%%%%%%%%%%%%%%%%%%%%%%%%%%%%%%%%%%%%%%%%%%%%%%%%%%%%%%%%%%%%%%%%%%%%%%%%%%%%%%%%%%%%%%%%%%%%%%
\draw [overlay,->,very thick,red,opacity=.5](x)-|(y)--(z);
\draw [overlay,very thick,red,opacity=.5](a)-|(y);
\node  at (z) [text width=40mm, below right , yshift= 8mm]{ 
\[
\begin{array}{rl}
x - (-z)+ 2z 
&=%%%%%%%%%%%%%%%%%%%%%%
 0
 \\[2mm]%%%%%%%%%%%%%%%%%%%%%%%%%%%%%%%%%%%%%%%%%%%%%%%%%%%%%%%%%%%%%%%%%%
x+ 3z
  &=%%%%%%%%%%%%%%%%%%%%%%
 0
  \\[3mm]%%%%%%%%%%%%%%%%%%%%%%%%%%%%%%%%%%%%%%%%%%%%%%%%%%%%%%%%%%%%%%%%%% 
  \marka{x1}{(-.1,-.2)}
  x
  &= %%%%%%%%%%%%%%%%%%%%%%
  -3z \marka{x2}{(.1,.5)}  
\end{array}
\]
};
%%%%%%%%%%%%%%%%%%%%%%%%%%%%%%%%%%%%%%%%%%%%%%%%%%%%%%%%%%%%%%%%%%%%%%%%%%%%%%%%%%%%%%%%%%%%%%%%%%%%%%%%%%%%%%%%%%%%%%%%%%%%%%%%%%%%%%%%%%%%%%%%%%%%%%%%%%%%%%%%%%%%%%%%%%%%%%%%%%%%
\draw[color=black!50!,rounded corners=1mm] (x1) rectangle (x2);
%%%%%%%%%%%%%%%%%%%%%%%%%%%%%%%%%%%%%%%%%%%%%%%%%%%%%%%%%%%%%%%%%%%%%%%%%%%%%%%%%%%%%%%%%%%%%%%%%%%%%%%%%%%%%%%%%%%%%%%%%%%%%%%%%%%%%%%%%%%%%%%%%%%%%%%%%%%%%%%%%%%%%%%%%%%%%%%%%%%%
}
 \vspace{5mm}

 \noindent   Por lo tanto, los vectores que cumplen son:
    \[
    \vec{w} = (-3z, -z, z), \quad z \in \R-\{0\} . 
    \]

\medskip


    \item \textbf{Hallar un vector perpendicular a ambos vectores.}

    El producto vectorial $\vec{u} \times \vec{v}$ da un vector perpendicular a ambos:
    \begin{align*}
    \vec{u} \times \vec{v} = 
    \begin{vmatrix}
    \iv & \jv & \kv \\
    0 & 1 & 1 \\
    1 & -1 & 2
    \end{vmatrix}
    &= \iv (1 \cdot 2 - 1 \cdot (-1)) - \jv (0 \cdot 2 - 1 \cdot 1) + \kv (0 \cdot (-1) - 1 \cdot 1)
    \\ 
    &=
    \iv (2 + 1) - \jv  (-1) + \kv (-1) = 3\vec{\imath} + \vec{\jmath} - \vec{k}         
    \end{align*}
    Entonces:
    \[
    \vec{n} = (3, 1, -1)
    \]

\medskip


    \item \textbf{Hallar un vector perpendicular a ambos y de módulo $4\sqrt{11}$. ¿Es único este vector?}

    Partimos del vector del punto anterior: $\vec{n} = (3,1,-1)$. Su norma es:
    \[
    \|\vec{n}\| = \sqrt{3^2 + 1^2 + (-1)^2} = \sqrt{9 + 1 + 1} = \sqrt{11}
    \]
    Lo multiplicamos por $4$:
    \[
    \vec{w} = 4 \cdot (3,1,-1) = (12,4,-4)
    \]
y por propiedad de modulo tenemos que este vector es el buscado. No es único, ya que también su opuesto $(-12,-4,4)$ cumple las mismas condiciones.

\medskip


    \item \textbf{¿Existe $\vec{w} = (w_1, w_2, w_3)$ tal que $\vec{w} \perp \vec{u}$, $\|\vec{w}\| = 3\sqrt{3}$ y $w_2 = 4$?}

    Condiciones:

\begin{itemize} 

    \item Como $\vec{w} \perp \vec{u} \ \Longrightarrow \ 
     \vec{w} \cdot \vec{u} =\marka{A}{(-.1,-.15)}\indicador{X}{(-.2,-.6)}{Y}{(2.3,-.9)}  0 , \qquad \vec{u} \neq \vec{\bf O}
     \qquad$ y $\quad  \vec{w} \neq \vec{\bf O}
    $ 

\bigskip  
    
    \item  $\vec{w} = (w_1, w_2, w_3)$ \marka{B}{(.1,.15)} \indicador{XX}{(1.1,-.6)}{YY}{(2.3,-.6)}

\bigskip
    
    \item $\|w\| 
    = \marka{C}{(-.2,-.15)}
    3 \sqrt{3}$

\bigskip
    
    \item $w_2 = 4$ \marka{D}{(.1,.15)} 


\end{itemize}
\tikz[overlay,remember picture] {

\draw [overlay,very thick,->,red,opacity=.5]
(B)-|(XX)--(YY);
\draw [overlay,very thick,red,opacity=.5]
(D)-|(XX);
\node  at (YY) [ right ]{
$\vec{w} = (w_1,4,w_3)$\marka{BB}{(.1,.15)}\indicador{X}{(.7,.6)}{Y}{(2.3,.6)}
};
%%%%%%%%%%%%%%%%%%%%%%%%%%%%%%%%%%%%%%%%%%%%%%%%%%%%%%%%%%%%%%%%%%%%%%%%%%%%%%%%%%%%%%%%%%%%%%%%%%%%%%%%%%%%%%%%%%%%%%%%%%%%%%%%%%%%%%%%%%%%%%%%%%%%%%%%%%%%%%%%%%%%%%%%%%%%%%%%%%%%%%%%%%%%%%%%%%%%%%%%%%%%%%%%%%%%%%%%%
\draw [overlay,very thick,teal,opacity=.5]
(A)to[out=-30,in=120](X);
\draw [overlay,very thick,->,teal,opacity=.5]
(BB)-|(X)--(Y);
\node  at (Y) [text width=80mm, below right ,xshift= -29mm,yshift=9mm]{
\begin{align*}
    \vec{w} \cdot \vec{u} \quad \quad  &= 0
    \\ 
    (w_1,4,w_3) 
    \cdot 
    (0,1,1)
    &=0
    \\[2mm]
     4+ w_3 &= 0
     \\[2mm]
     w_3 &= -4
    \end{align*} 
};
}

\vspace{3mm}

    Entonces $\vec{w} = (w_1, 4, -4)$ y su norma es:
    \[
    \|\vec{w}\| = \sqrt{w_1^2 + 4^2 + (-4)^2} = \sqrt{w_1^2 + 16 + 16} = \sqrt{w_1^2 + 32}
    \]
    Queremos que $\|\vec{w}\| =  4$ y nos queda:
    \[
    \sqrt{w_1^2 + 32} = 3\sqrt{3} \qquad\Longrightarrow\qquad w_1^2 + 32 = 27 \qquad\Longrightarrow\qquad w_1^2 = -5
    \]
    No tiene solución real. \textbf{No existe tal vector}.

\end{enumerate}


 















\item Sabiendo que $\vec{u} \times \vec{v} = (2,5,-4)$:

\begin{enumerate}[a)]
    \item Hallar el área del paralelogramo determinado por los vectores $\vec{u}$ y $\vec{v}$. Mostrar con un dibujo la situación.
    
    \item Usando propiedades, hallar: \qquad\qquad 
    \begin{minipage}{.4\textwidth}
    \begin{enumerate}[i)]
    \begin{multicols}{2}
        \item $\vec{v} \times \vec{u}$
        \item $\vec{u} \times (2\vec{v})$         
    \end{multicols}
    \end{enumerate}
        
    \end{minipage}
    
    \item Representar en cada diagrama el vector resultado (utilizar la regla de la mano derecha).

\vspace{-3mm}
\begin{multicols}{2} 
    \begin{center}
    Figura 1: $\vec{u} \times \vec{v}$\\[3mm]
    % Figura 1: u x v
    \begin{tikzpicture}[xscale=1, yscale=.6,>=stealth]
        % Plano
        \draw[gray!80] (-1.5,-.5) -- (4,-.5) -- (5,2.5) -- (-.5,2.5) -- cycle;

        \coordinate (O) at (0,.3);
        
        \draw[fill=black] (O) circle (1.5pt);
        % Vectores
        \draw[->, thick] (O) --node[above] {$\vec{v}$} (2,1.5) ;
        \draw[->, thick] (O) --node[below] {$\vec{u}$} (3,0) ;
    \end{tikzpicture}    
    \end{center}

    \vspace{1cm}

    \begin{center}
    Figura 2: $\vec{v} \times \vec{u}$\\[3mm]
    % Figura 2: v x u 
    \begin{tikzpicture}[xscale=1, yscale=.6,>=stealth]
        % Plano
        \draw[gray!80] (-1.5,-1) -- (4,-1) -- (5,2) -- (-.5,2) -- cycle;

        \coordinate (O) at (0,.3);
        
        \draw[fill=black] (O) circle (1.5pt);
        % Vectores
        \draw[->, thick] (O) -- node[below] {$\vec{u}$} (3,0);
        \draw[->, thick] (O) -- node[above] {$\vec{v}$}(2,1.5);
    \end{tikzpicture}
    \end{center}
\end{multicols}




    \item Determinar el valor de $k \in \mathbb{R}$ tal que $\| \text{Proy}_{\vec{u} \times \vec{v}} \vec{w}\| = 2\sqrt{5} $ siendo $\vec{w} = (k,-1,-k)$.
\end{enumerate}








{\bf Solución:}

\begin{enumerate}[a)]

\item \textbf{Hallar el área del paralelogramo determinado por los vectores $ \vec{u} $ y $ \vec{v} $. Mostrar con un dibujo la situación.}

    El área del paralelogramo formado por dos vectores se obtiene mediante el módulo del producto vectorial:
    \[
    \text{Área} = \| \vec{u} \times \vec{v} \| = \|(2,\, 5,-4) \|  = \sqrt{2^2 + 5^2 + (-4)^2} = \sqrt{4 + 25 + 16} = \sqrt{45} = 3\sqrt{5}
    \]

    \begin{center}
    \begin{asy}
    import three;
size(5cm);

// Configuración de la proyección
currentprojection = perspective(6, .5, 3);

// Ejes coordenados draw(Label("$x$", EndPoint), O--(1,0,0), Arrow3);draw(Label("$y$", EndPoint), O--(0,1,0), Arrow3);draw(Label("$z$", EndPoint), O--(0,0,1), Arrow3);

// Definición de vectores
triple u = (.6, 1.1, 0);
triple v = (-.35, .75, 0);
triple uxv = cross(u, v); // Producto cruzado
triple vxu = cross(v, u); // Producto cruzado
triple A = (.0, .0, .0);
triple uu = A+u;
triple vv = A+v;
triple  uuxvv = A+ uxv ;
triple  vvxuu = A+ vxu ;
// Dibujar el plano donde están los vectores u y v
draw(surface(A--uu--(u+v+A)--vv--cycle), gray+opacity(0.5));

// Dibujar vectores
draw(A--uu, red+linewidth(1.5), Arrow3);
label("$\vec{u}$", uu, align=Z);
draw(A--vv, red+linewidth(1.5), Arrow3);
label("$\vec{v}$", vv, align=Z);
draw(A--uuxvv, red+linewidth(1.5), Arrow3);
label("$\vec{u} \times \vec{v}$", uuxvv, align=Z); 

// Dibujar sistema de ejes (base canónica) 
\end{asy}
    \end{center}

    \item \textbf{Usando propiedades, hallar:}

    \begin{enumerate}[i)]
    \begin{multicols}{2}
        \item $ \vec{v} \times \vec{u} = - (\vec{u} \times \vec{v}) = (-2, -5, 4) $

        \item $ \vec{u} \times (2\vec{v}) = 2 (\vec{u} \times \vec{v}) = 2(2,5,-4) = (4,10,-8) $
        
    \end{multicols}
    \end{enumerate}

    \item \textbf{Representar en cada diagrama el vector resultado (utilizar la regla de la mano derecha).}


\begin{comment}
    \begin{multicols}{2}
    \textbf{Figura 1: $ \vec{u} \times \vec{v} $}

    
\begin{asy}
    import three;
size(5cm);

// Configuración de la proyección
currentprojection = perspective(6, .5, 3);

// Ejes coordenados
draw(Label("$x$", EndPoint), O--(1,0,0), Arrow3);
draw(Label("$y$", EndPoint), O--(0,1,0), Arrow3);
draw(Label("$z$", EndPoint), O--(0,0,1), Arrow3);

// Definición de vectores
triple u = (.5, -0.25, -.25);
triple v = (.25, .75, .25);
triple uxv = cross(u, v); // Producto cruzado
triple vxu = cross(v, u); // Producto cruzado
triple A = (.5, .5, .5);
triple uu = A+u;
triple vv = A+v;
triple  uuxvv = A+ uxv ;
triple  vvxuu = A+ vxu ;
// Dibujar el plano donde están los vectores u y v
draw(surface(A--uu--(u+v+A)--vv--cycle), gray+opacity(0.2));

// Dibujar vectores
draw(A--uu, red+linewidth(1.5), Arrow3);
label("$\vec{u}$", uu, align=Z);
draw(A--vv, red+linewidth(1.5), Arrow3);
label("$\vec{v}$", vv, align=Z);
draw(A--uuxvv, red+linewidth(1.5), Arrow3);
label("$\vec{u} \times \vec{v}$", uuxvv, align=Z); 

// Dibujar sistema de ejes (base canónica)
 
\end{asy} 

    \columnbreak

    \textbf{Figura 2: $ \vec{v} \times \vec{u} $}

    \begin{asy}
    import three;
size(5cm);

// Configuración de la proyección
currentprojection = perspective(6, .5, 3);

// Ejes coordenados
draw(Label("$x$", EndPoint), O--(1,0,0), Arrow3);
draw(Label("$y$", EndPoint), O--(0,1,0), Arrow3);
draw(Label("$z$", EndPoint), O--(0,0,1), Arrow3);

// Definición de vectores
triple u = (.5, -0.25, -.25);
triple v = (.25, .75, .25);
triple uxv = cross(u, v); // Producto cruzado
triple vxu = cross(v, u); // Producto cruzado
triple A = (.5, .5, .5);
triple uu = A+u;
triple vv = A+v;
triple  uuxvv = A+ uxv ;
triple  vvxuu = A+ vxu ;
// Dibujar el plano donde están los vectores u y v
draw(surface(A--uu--(u+v+A)--vv--cycle), gray+opacity(0.2));

// Dibujar vectores
draw(A--uu, red+linewidth(1.5), Arrow3);
label("$\vec{u}$", uu, align=Z);
draw(A--vv, red+linewidth(1.5), Arrow3);
label("$\vec{v}$", vv, align=Z); 
draw(A--vvxuu, red+linewidth(1.5), Arrow3);
label("$\vec{v} \times \vec{u}$", vvxuu+(0,0,-.1), align=Z);

// Dibujar sistema de ejes (base canónica)
 
\end{asy}
    \end{multicols}
\end{comment}
    

\begin{multicols}{2} 
    \begin{center}
    Figura 1: $\vec{u} \times \vec{v}$\\[3mm]
    % Figura 1: u x v
    \begin{tikzpicture}[xscale=1, yscale=.6,>=stealth]
        % Plano
        \draw[gray!80] (-1.5,-.5) -- (4,-.5) -- (5,2.5) -- (-.5,2.5) -- cycle;

        \coordinate (O) at (0,.3);
        
        \draw[fill=black] (O) circle (1.5pt);
        % Vectores
        \draw[->, thick] (O) --node[above] {$\vec{v}$} (2,1.5) ;
        \draw[->, thick] (O) --node[below] {$\vec{u}$} (3,0) ;
         \draw[->, thick] (O) --node[right,pos=.8]{$\vec{u}\times\vec{v}$}(0,3);
    \end{tikzpicture}    
    \end{center}

    \vspace{1cm}

    \begin{center}
    Figura 2: $\vec{v} \times \vec{u}$\\[3mm]
    % Figura 2: v x u 
    \begin{tikzpicture}[xscale=1, yscale=.6,>=stealth]
        % Plano
        \draw[gray!80] (-1.5,-1) -- (4,-1) -- (5,2) -- (-.5,2) -- cycle;

        \coordinate (O) at (0,.3);
        
        \draw[fill=black] (O) circle (1.5pt);
        % Vectores
        \draw[->, thick] (O) -- node[below] {$\vec{u}$} (3,0);
        \draw[->, thick] (O) -- node[above] {$\vec{v}$}(2,1.5);
         \draw[->, thick] (O) --node[right,pos=.8] {$\vec{v}\times\vec{u}$} (0,-3);
    \end{tikzpicture}
    \end{center}
\end{multicols}





    \item \textbf{Determinar el valor de $ k \in \mathbb{R} $ tal que $ \left\| \text{Proy}_{\vec{u} \times \vec{v}} \vec{w} \right\| = 2\sqrt{5} $, siendo $ \vec{w} = (k, -1, -k) $.}

    Sabemos que:
    \[
    \| \text{Proy}_{\vec{n}} \vec{w} \| = \left| \frac{\vec{w} \cdot \vec{n}}{\|\vec{n}\|} \right|
    \quad \text{con } \vec{n} = \vec{u} \times \vec{v} = (2,5,-4)
    \]
    Entonces:
    \[
    \vec{w} \cdot \vec{n} = (k, -1, -k) \cdot (2,5,-4) = 2k + 5(-1) + (-4)(-k) = 2k - 5 + 4k = 6k - 5
    \]
    \[
    \| \vec{n} \| = \sqrt{2^2 + 5^2 + (-4)^2} = \sqrt{45} = 3\sqrt{5}
    \]
    \[
    \left\| \text{Proy}_{\vec{u} \times \vec{v}} \vec{w} \right\| = 2\sqrt{5} \qquad\Rightarrow\qquad \left| \frac{6k - 5}{3\sqrt{5}} \right| = 2\sqrt{5}
    \qquad\Rightarrow\qquad |6k - 5| = 6\sqrt{5}  
    \]
    Resolviendo:
    \[
    \begin{cases}
    6k - 5 = 6\sqrt{5} \qquad\Rightarrow\qquad  k = \dfrac{5 + 6\sqrt{5}}{6} \\[1mm]
    o
    \\
    6k - 5 = -6\sqrt{5} \qquad\Rightarrow\qquad  k = \dfrac{5 - 6\sqrt{5}}{6}
    \end{cases}
    \]

    Por lo tanto, los valores posibles de $ k $ son:\quad $k = \frac{5 - 6\sqrt{5}}{6} \qquad\text{ o }\qquad  k = \frac{5 + 6\sqrt{5}}{6}$.

\end{enumerate} 





 









\item Sean los vectores $\vec{v} = (a+1, -1, 4)$ y $\vec{w} = (2b+2, -2, 8)$, determinar el o los valores de $a \in \mathbb{R}$ y $b \in \mathbb{R}$ para que los vectores $\vec{v}$ y $\vec{w}$ sean paralelos. Resolver de dos maneras diferentes.


\textbf{Solución:}

\begin{itemize}
    \item Método 1: Se basa en la definición de paralelismo. Dos vectores son paralelos si existe un escalar $\lambda$, no nulo, tal que:

\[
\vec{v} = \lambda \vec{w}.
\]

Esto significa que las componentes de los vectores deben ser proporcionales:

\[
\frac{a+1}{2b+2} \ =\indicador{A}{(-.3,-.1)}{B}{(-.7,-.9)} \  \frac{-1}{-2} \  =\indicador{X}{(-.2,-.1)}{Y}{(.6,-.9)} \  \frac{4}{8}.
\]





\tikz[overlay,remember picture] {
\node  at (B) [text width=31mm, below ,xshift=-6mm ,yshift=4mm ]{
\begin{align*} 
\frac{a+1}{2b+2} \ &=  \  \frac{1}{2}
    \\[2mm]
    2(a+1) &= 2b+2
    \\[2mm]
    2a + 2 &= 2b + 2
    \\[2mm]
    2a &= 2b
    \\[2mm]
    a &= b 
    \end{align*}  
    };
\draw [overlay,very thick,->,teal,opacity=.5]
(A)--(B);
%%-------------------------------------------------------------------------------------------------------------------- 
\draw [overlay,very thick,<-,teal,opacity=.5]
(X)|-(Y);
\node[fill=blue!80!white!10!, rectangle,rounded corners=3mm] at (Y) [text width=38mm,right ]{
 ya es una identidad verdadera, las dos fracciones 
};
}


\vspace{36mm} 

\hfill Entonces, la solución obtenida es: $a = b.$


\item  Método 2: Uso del producto vectorial

En $\R^3$ si dos vectores son paralelos, su producto vectorial debe ser el vector nulo:

\[
\vec{v} \times \vec{w} = \mathbf{0}.
\]

Calculamos:

\begin{align*}
\vec{v} \times \vec{w} &=
\begin{vmatrix}
\iv & \jv & \kv \\
a+1 & -1 & 4 \\
2b+2 & -2 & 8
\end{vmatrix}
\\
 &=
\iv \begin{vmatrix} -1 & 4 \\ -2 & 8 \end{vmatrix}
- \jv \begin{vmatrix} a+1 & 4 \\ 2b+2 & 8 \end{vmatrix}
+ \kv \begin{vmatrix} a+1 & -1 \\ 2b+2 & -2 \end{vmatrix}
\\
& =
\iv \left(-1\cdot 8 - 4\cdot (-2) \rule{0mm}{4mm} \right) 
- \jv \left((a+1)\cdot 8 - 4\cdot (2b+2)  \rule{0mm}{4mm}\right) 
+ \kv \left( (a+1)\cdot(-2) - (-1)\cdot(2b+2) \rule{0mm}{4mm} \right)
\\
& =
\iv \cdot 0  
- \jv \left(8a + 8 - 8b - 8  \rule{0mm}{4mm}\right) 
+ \kv \left(  -2a - 2 + 2b + 2 \rule{0mm}{4mm} \right)
\\
& =
\iv \cdot 0  
- \jv \left( 8a - 8b  \rule{0mm}{4mm}\right) 
+ \kv \left( -2a + 2b \rule{0mm}{4mm} \right)
\end{align*}

Para que $\vec{v} \times \vec{w} = \mathbf{0}$, todas las componentes deben ser 0:

\[
8a - 8b = 0 \quad \Rightarrow \quad 8a = 8b \quad \Rightarrow \quad a = b.
\]

\[
-2a + 2b = 0 \quad \Rightarrow \quad -2a = -2b \quad \Rightarrow \quad a = b.
\]

La  solución es: \qquad $a = b.$ 


\textbf{Resultado final:} Para que los vectores sean paralelos, $ a = b $.
 
\end{itemize}









\item  En cada caso, dados los vectores $\vec{u}$, $\vec{v}$ y $\vec{w}$, calcular:

\begin{enumerate}[ a)] 
    \item  El volumen del paralelepípedo que ellos determinan.
    \item  El área del paralelogramo que forman los vectores $\vec{v}$ y $\vec{w}$.
    \item  La altura y el vector altura correspondientes al paralelepípedo cuya base está formada por los vectores $\vec{v}$ y $\vec{w}$.
        \begin{enumerate}[ i)] 
            \item  $\vec{u} = (2, 3, -2)$,\quad $\vec{v} = (4, -1, -1)$\quad y \quad $\vec{w} = (3, 1, 2)$.

            
            \item  $\vec{u} = (1, 6, -3)$,\quad $\vec{v} = (8, 8, -4)$ \quad y \quad $\vec{w} = (2, -4, 2)$. 
        \end{enumerate}
    \item Representar con un dibujo (esquema plano) el paralelepípedo y todos sus elementos (área de la base, altura y
vector altura)
\end{enumerate}











\textbf{Solución:}

\textbf{Recordemos:}
\begin{enumerate}[a)]
    \item El \textbf{volumen del paralelepípedo}, esta dado por: \quad$V = |\vec{u} \cdot (\vec{v} \times \vec{w})|$.

    \item Calcular el \textbf{área del paralelogramo} que forman $\vec{v}$ y $\vec{w}$:\quad$ A = \|\vec{v} \times \vec{w}\|$ 

    \item Calcular la \textbf{altura} y el \textbf{vector altura} respecto de la base formada por $\vec{v}$ y $\vec{w}$: \quad $h = \frac{V}{A}, \quad \vec{h} = \frac{V}{\|\vec{v} \times \vec{w}\|^2} (\vec{v} \times \vec{w})$.

    \item \textbf{Representar con un esquema} el paralelepípedo y sus elementos (no incluido en este documento).
\end{enumerate}

\vspace{0.5cm}

\textbf{i)} $\vec{u} = (2, 3, -2)$,\quad $\vec{v} = (4, -1, -1)$,\quad $\vec{w} = (3, 1, 2)$

\begin{itemize}
    \item Producto vectorial: \quad$\vec{v} \times \vec{w} = 
    \begin{vmatrix}
    \mathbf{i} & \mathbf{j} & \mathbf{k} \\
    4 & -1 & -1 \\
    3 & 1 & 2
    \end{vmatrix}
    = (-1)\mathbf{i} -11\mathbf{j} +7\mathbf{k} = (-1, -11, 7)$
     
\medskip
    \item Producto mixto:\quad$\vec{u} \cdot (\vec{v} \times \vec{w}) = (2,3,-2) \cdot (-1,-11,7) = -2 - 33 -14 = -49
    \quad\Rightarrow\quad V = 49$.
    
\medskip

    \item Área de la base:\quad$A = \|(-1, -14, 7)\| = \sqrt{1 + 121 + 49} = \sqrt{171}$.
    
\medskip

    \item Altura:\quad$h = \frac{49}{\sqrt{171}\, } = \frac{58}{\sqrt{171}} \approx 3.74$. 
    
\medskip

    \item Vector altura:\quad$ \displaystyle\vec{h} = \text{Proy}_{\vec{v} \times \vec{w}} \vec{u} =   
\frac{ \vec{u} \cdot ( \vec{v}\times \vec{w})}{ \|  \vec{v}\times \vec{w}\|^2} \  \vec{v} \times \vec{w} =  \frac{-49}{171} (-1, -11, 7) = \left(\frac{49}{171} , \frac{539}{171}, -\frac{343}{171} \right)$.
\end{itemize}




    \begin{center}
    \begin{asy} 
    import three;
size(5cm);
transform3 t = rotate(20, X)*rotate(-15, Y);
// Configuración de la proyección
currentprojection = perspective(5, -3, 3); 

// Ejes coordenados draw(Label("$x$", EndPoint), O--(1,0,0), Arrow3);draw(Label("$y$", EndPoint), O--(0,1,0), Arrow3);draw(Label("$z$", EndPoint), O--(0,0,1), Arrow3);

// Definición de vectores
triple u = (2,3,-2);
triple v = (4,-1,-1);
triple w = (3,1,2);
triple vxu = .25*cross(v, u); // Producto cruzado 
triple O= (.0, .0, .0);
triple uu = O+u;
triple vv = O+v;
triple ww = O+w;
triple  vvxuu = O+ vxu ; 

// Dibujar el planos
draw(t*surface(O--u--(u+v)--v--cycle), gray+opacity(0.2));
draw(t*surface(O--u--(u+w)--w--cycle), gray+opacity(0.2));
draw(t*surface(O--v--(v+w)--w--cycle), gray+opacity(0.2)); 

draw(t*surface(w--(w+u)--(w+u+v)--(w+v)--cycle), gray+opacity(0.2));
draw(t*surface(v--(v+u)--(v+u+w)--(v+w)--cycle), gray+opacity(0.2));
draw(t*surface(u--(u+v)--(u+v+w)--(u+w)--cycle), gray+opacity(0.2)); 



// Dibujar arista
draw(t*u--t*(u+v)--t*v, black+linewidth(0.5)+opacity(0.2));
draw(t*w--t*(w+v)--t*v, black+linewidth(0.5)+opacity(0.2));
draw(t*u--t*(u+w)--t*w, black+linewidth(0.5)+opacity(0.2));
draw(t*(u+v)--t*(u+w+v)--t*(w+v), black+linewidth(0.5)+opacity(0.2));
draw(t*(u+w)--t*(u+w+v), black+linewidth(0.5)+opacity(0.2));



// Dibujar vectores
draw(t*O--t*uu, red+linewidth(1), Arrow3);
label("$\vec{u}$", t*uu, align=Z);
draw(t*O--t*vv, red+linewidth(1), Arrow3);
label("$\vec{v}$", t*vv, align=Z);
draw(t*O--t*ww, red+linewidth(1), Arrow3);
label("$\vec{w}$", t*ww, align=Z);
draw(t*O--t*vvxuu, grey+linewidth(1), Arrow3);
label("$\vec{v} \times \vec{u}$", t*vvxuu, align=Z);
\end{asy}
    \end{center}




\vspace{0.5cm}

\textbf{ii)} $\vec{u} = (1, 6, -3)$,\quad $\vec{v} = (8, 8, -4)$,\quad $\vec{w} = (2, -4, 2)$

\begin{itemize}
    \item Producto vectorial:\quad$\vec{v} \times \vec{w} = 
    \begin{vmatrix}
    \mathbf{i} & \mathbf{j} & \mathbf{k} \\
    8 & 8 & -4 \\
    2 & -4 & 2
    \end{vmatrix}
    = (0, -24, -48)$.
    
\medskip 

    \item Producto mixto:\quad$\vec{u} \cdot (\vec{v} \times \vec{w}) = (1,6,-3) \cdot (0, -24, -48) = -144 + 144 = 0
    \quad \Rightarrow\quad V = 0$.
    
\medskip 

    \item Área de la base:\quad$A = \|(0, -24, -48)\| = \sqrt{576 + 2304} = \sqrt{2880}$.
    
\medskip 

    \item Altura:\quad$h = \frac{0}{\sqrt{2880}} = 0$.
    
\medskip
 

    \item Vector altura:\quad$\vec{h} = \vec{0}$. 
\end{itemize} 





    \begin{center}
    \begin{asy} import three;
size(5cm);
transform3 t = rotate(40, X)*rotate(-15, Y);
// Configuración de la proyección
currentprojection = perspective(5, -3, 3); 

// Ejes coordenados draw(Label("$x$", EndPoint), O--(1,0,0), Arrow3);draw(Label("$y$", EndPoint), O--(0,1,0), Arrow3);draw(Label("$z$", EndPoint), O--(0,0,1), Arrow3);

// Definición de vectores
triple u = (1,6,-3);
triple v = (8,8,-4);
triple w = (2,-4,2);
triple vxu = .25*cross(v, u); // Producto cruzado 
triple O= (.0, .0, .0);
triple uu = O+u;
triple vv = O+v;
triple ww = O+w;
triple  vvxuu = O+ vxu ; 

// Dibujar el planos
draw(t*surface(O--u--(u+v)--v--cycle), gray+opacity(0.2));
draw(t*surface(O--u--(u+w)--w--cycle), gray+opacity(0.2));
draw(t*surface(O--v--(v+w)--w--cycle), gray+opacity(0.2)); 

draw(t*surface(w--(w+u)--(w+u+v)--(w+v)--cycle), gray+opacity(0.2));
draw(t*surface(v--(v+u)--(v+u+w)--(v+w)--cycle), gray+opacity(0.2));
draw(t*surface(u--(u+v)--(u+v+w)--(u+w)--cycle), gray+opacity(0.2)); 



// Dibujar arista draw(t*u--t*(u+v)--t*v, black+linewidth(0.5)+opacity(0.2)); draw(t*w--t*(w+v)--t*v, black+linewidth(0.5)+opacity(0.2)); draw(t*u--t*(u+w)--t*w, black+linewidth(0.5)+opacity(0.2)); draw(t*(u+v)--t*(u+w+v)--t*(w+v), black+linewidth(0.5)+opacity(0.2)); draw(t*(u+w)--t*(u+w+v), black+linewidth(0.5)+opacity(0.2));



// Dibujar vectores
draw(t*O--t*uu, red+linewidth(1), Arrow3);
label("$\vec{u}$", t*uu, align=Z);
draw(t*O--t*vv, red+linewidth(1), Arrow3);
label("$\vec{v}$", t*vv, align=Z);
draw(t*O--t*ww, red+linewidth(1), Arrow3);
label("$\vec{w}$", t*ww, align=Z);
draw(t*O--t*vvxuu, grey+linewidth(1), Arrow3);
label("$\vec{v} \times \vec{u}$", t*vvxuu, align=Z);
\end{asy}
    \end{center}


Lo que esta pasando es que los tres vectores están sobre un mismo plano, por eso no determinan ningún volumen, es decir {\bf los vectores son coplanares}.





\item Dados los vectores $ \vec{a} =\iv + x \jv - \kv $ y $ \vec{b} = 2 \jv + \iv $. Hallar todos los valores $ x \in \mathbb{R} $ para que el área del paralelogramo determinado por ellos sea $ 3 \, \text{u}^2 $.
    


{\bf Solución} 

El área del paralelogramo es igual al módulo del producto vectorial:

\[
A = \left\| \vec{a} \times \vec{b} \right\|
\]

Calculamos el producto vectorial:

\begin{align*}
\vec{a} \times \vec{b} &=
\begin{vmatrix}
\iv & \jv & \kv \\
1 & x & -1 \\
1 & 2 & 0 \\
\end{vmatrix}
= \iv \cdot
\begin{vmatrix}
x & -1 \\
2 & 0 \\
\end{vmatrix}
- \jv \cdot
\begin{vmatrix}
1 & -1 \\
1 & 0 \\
\end{vmatrix}
+ \kv \cdot
\begin{vmatrix}
1 & x \\
1 & 2 \\
\end{vmatrix}
\\
&= \iv (x \cdot 0 - (-1) \cdot 2) - \jv (1 \cdot 0 - (-1) \cdot 1) + \kv (1 \cdot 2 - x \cdot 1)
\\
&= 2\iv - \jv (1) + (2 - x)\kv
\\
\vec{a} \times \vec{b} &= 2\iv - \jv + (2 - x)\kv 
\end{align*}

Ahora calculamos el módulo:

\[
\left\| \vec{a} \times \vec{b} \right\| = \sqrt{2^2 + (-1)^2 + (2 - x)^2}
= \sqrt{4 + 1 + (2 - x)^2}
= \sqrt{5 + (2 - x)^2}
\]

Queremos que esta norma sea 3:

\[
\sqrt{5 + (2 - x)^2} = 3
\qquad\Rightarrow\qquad 5 + (2 - x)^2 = 9
\qquad\Rightarrow\qquad  (2 - x)^2 = 4
\qquad\Rightarrow\qquad \]
\[
2 - x = 2 \Rightarrow x = 0 \qquad\textbf{ o }\qquad
2 - x = -2 \Rightarrow x = 4
\]

\textbf{Respuesta:} Los valores de $ x \in \mathbb{R} $ que satisfacen la condición son:
\[
\boxed{x = 0 \qquad \text{ o } \qquad x = 4}
\]



\item 
    \begin{enumerate}[a)]
        \item Determinar si los puntos dados son coplanares:  \ $A = (1, 1, 6), \  B = (2, 3, 5), \  C = (8, 4, 6), \  D = (2, 1, 3)$.
        
        \item Dados tres vectores $ \vec{u} = 2 \jv - \iv + 5 \kv $, $ \vec{v} = \alpha \jv - \kv $, $ \vec{w} = 4 \kv - \jv - \iv $. Hallar el valor de $ \alpha \in \mathbb{R} $ de modo tal que los vectores resulten coplanares.
    \end{enumerate} 









{\bf Solución:}

\begin{enumerate}[a)]
    \item \textbf{¿Los puntos $A = (1, 1, 6), \  B = (2, 3, 5), \  C = (8, 4, 6), \  D = (2, 1, 3)$ son coplanares?}

    Para que cuatro puntos sean coplanares, el volumen del paralelepípedo formado por tres vectores que parten de un mismo punto debe ser cero. Tomamos como origen común el punto $A$ y formamos los vectores:

    \[
    \overrightarrow{AB} = \overrightarrow{OB} - \overrightarrow{OA} = (2{-}1,\ 3{-}1,\ 5{-}6) = (1,\ 2,\ -1)
    \]
    \[
    \overrightarrow{AC} = \overrightarrow{OC} - \overrightarrow{OA} = (8{-}1,\ 4{-}1,\ 6{-}6) = (7,\ 3,\ 0)
    \]
    \[
    \overrightarrow{AD} = \overrightarrow{OD} - \overrightarrow{OA} = (2{-}1,\ 1{-}1,\ 3{-}6) = (1,\ 0,\ -3)
    \] 
    Calculamos el producto mixto: $\overrightarrow{AB} \cdot (\overrightarrow{AC} \times \overrightarrow{AD})$.

    Primero el producto vectorial:
    \[
    \overrightarrow{AC} \times \overrightarrow{AD} =
    \begin{vmatrix}
        \iv & \jv & \kv \\
        7 & 3 & 0 \\
        1 & 0 & -3
    \end{vmatrix}
    = \iv(3\cdot(-3) - 0\cdot0) - \jv(7\cdot(-3) - 0\cdot1) + \kv(7\cdot0 - 3\cdot1)
    \]
    \[
    = \iv(-9) - \jv(-21) + \kv(-3) = (-9,\ 21,\ -3)
    \]

    Ahora el producto escalar:
    \[
    \overrightarrow{AB} \cdot (\overrightarrow{AC} \times \overrightarrow{AD}) = (1,\ 2,\ -1) \cdot (-9,\ 21,\ -3) = 1\cdot(-9) + 2\cdot21 + (-1)\cdot(-3) = -9 + 42 + 3 = 36
    \]

    Como el producto mixto es diferente de cero, los puntos \textbf{no son coplanares}.

    \vspace{5mm}

    \item \textbf{Hallar $\alpha \in \mathbb{R}$ para que los vectores $\vec{u} = 2\jv - \iv + 5\kv$, $\vec{v} = \alpha \jv - \kv$, $\vec{w} = 4\kv - \jv - \iv$ sean coplanares.}

    Los vectores serán coplanares si el producto mixto es cero:
    \[
    \vec{u} \cdot (\vec{v} \times \vec{w}) = 0,
    \]

    \noindent consideremos los vectores: \quad $\vec{u} = (-1,\ 2,\ 5), \quad \vec{v} = (0,\ \alpha,\ -1), \quad \vec{w} = (-1,\ -1,\ 4)$ 

    Calculamos el producto vectorial $\vec{v} \times \vec{w}$:
    \[
    \vec{v} \times \vec{w} =
    \begin{vmatrix}
        \iv & \jv & \kv \\
        0 & \alpha & -1 \\
        -1 & -1 & 4
    \end{vmatrix}
    = \iv(\alpha\cdot4 - (-1)\cdot(-1)) - \jv(0\cdot4 - (-1)\cdot(-1)) + \kv(0\cdot(-1) - \alpha\cdot(-1))
    \]
    \[
    = \iv(4\alpha - 1) - \jv(0 - 1) + \kv(\alpha) = (4\alpha - 1,\ 1,\ \alpha)
    \]

    Ahora calculamos el producto escalar:
    \[
    \vec{u} \cdot (\vec{v} \times \vec{w}) = (-1)(4\alpha - 1) + 2(1) + 5(\alpha)
    = -4\alpha + 1 + 2 + 5\alpha = \alpha + 3
    \]

    Para que sean coplanares:
    \[
    \alpha + 3 = 0 \quad \Rightarrow \quad \alpha = -3
    \]

    \textbf{Respuesta:} Los vectores son coplanares si y sólo si $\boxed{\alpha = -3}$.
\end{enumerate}
   








\end{enumerate}




\hfill{\bf \huge Fin.}

\end{document}